\documentclass[11pt]{article}

\usepackage[margin=1in]{geometry}
\usepackage{amsmath,amssymb,amsthm,mathtools}
\usepackage{cite}
\usepackage{enumitem}
\usepackage{hyperref}

\newtheorem{theorem}{Theorem}[section]
\newtheorem{lemma}[theorem]{Lemma}
\newtheorem{proposition}[theorem]{Proposition}
\newtheorem{corollary}[theorem]{Corollary}
\newtheorem{definition}[theorem]{Definition}
\newtheorem{remark}[theorem]{Remark}

\newcommand{\C}{\mathcal C}
\newcommand{\F}{\mathcal F}
\newcommand{\E}{\mathcal E}
\newcommand{\J}{\mathcal J}
\newcommand{\K}{\mathcal K}
\newcommand{\R}{\mathcal R}
\newcommand{\T}{\mathcal T}
\newcommand{\Prb}{\mathbb P}
\newcommand{\EE}{\mathbb E}
\newcommand{\bfH}{\mathbf H}
\newcommand{\KL}{D}
\newcommand{\Coll}{\operatorname{Coll}}
\newcommand{\Split}{\operatorname{Split}}
\newcommand{\Law}{\operatorname{Law}}
\newcommand{\rank}{\operatorname{rank}}

\title{A Public Normal-Form Framework for Entropy Bookkeeping in Sunflower
Bounds}
\author{
Yuli Li\\
Department of Chemistry and Chemical Biology, Harvard University\\
Cambridge, MA 02138, USA
\and
Jingping Yang\\
ByteDance Ltd.\\
Douyin Building, Building 2, Yard 18, West Third Ring Road North\\
Haidian District, Beijing, China, 100098
\and
Renhao Zhang\\
Manning College of Information \& Computer Sciences\\
University of Massachusetts Amherst\\
140 Governors Dr., Amherst, MA 01003, USA
}
\date{May 2026}

\begin{document}
\maketitle

\begin{abstract}
We describe a public normal-form framework for entropy-based approaches to the
three-petal sunflower problem.  The framework organizes tuple-level stopping
rules, multi-sample projected shadows, conditioned copying with public
transcript cost, deterministic DRC block selection, bounded active reservoirs,
and delayed-window potential bookkeeping into a single finite-state interface.
The main formal statement is an assembly theorem: assuming the local exits
satisfy the publicness, copying, pruning, and ledger certificates specified
below, the terminal shadow obeys the entropy criterion needed for a
constant-base rank recurrence for balanced two-collision codes.  Together with
the standard random-coloring reduction, a complete verification of these
interfaces would imply a constant-base bound for three-petal sunflower-free set
families.  The purpose of the manuscript is to isolate the proposed
alphabet-free fixed-window mechanism and to make the remaining verification
obligations explicit.
\end{abstract}

\tableofcontents

\section{Introduction and Conditional Reduction}

This paper develops a public-normal-form route to a possible constant-base
sunflower bound.  The emphasis is on the method and on the interfaces that a
complete proof would have to verify, not on claiming an unconditional
settlement of the sunflower conjecture.  The central assembly statement is an
interface theorem: once all local exits are available in the stated public
normal form, the global entropy bookkeeping closes.

Sunflowers, or $\Delta$-systems, were introduced by Erdos and Rado
\cite{ER60}.  Their classical sunflower lemma gives a bound of order
$k!(r-1)^k$ for $r$-petal sunflowers in $k$-uniform families, and they
conjectured that for each fixed $r$ this can be replaced by $C_r^k$.  For
three petals, the conjecture remains the relevant benchmark for constant-base
entropy or encoding arguments.  Kostochka obtained the first asymptotic
improvement over the Erdos--Rado bound \cite{Kostochka99}.  The modern
breakthrough was the spread/robust-sunflower method of Alweiss, Lovett, Wu,
and Zhang \cite{ALWZ21}; subsequent entropy and probabilistic presentations by
Rao \cite{Rao20}, Tao \cite{Tao20}, and Bell, Chueluecha, and Warnke
\cite{BCW21} gave the current logarithmic-base form of the method.  Related
threshold work of Frankston, Kahn, Narayanan, and Park \cite{FKNP21} shows how
robust sunflower ideas interact with expectation-threshold problems.
Classical lower-bound constructions and finite-universe variants also play an
important role in the surrounding picture \cite{AHS72,ES78}.

There is also a bounded-alphabet route through the polynomial method and slice
rank.  The Croot--Lev--Pach and Ellenberg--Gijswijt cap-set breakthroughs
\cite{CLP17,EG17} and Tao's slice-rank formulation \cite{Tao16}, together
with the sunflower-free applications of Alon,
Shpilka, and Umans \cite{ASU13} and Naslund--Sawin \cite{NS17}, show that
strong bounds are available when the alphabet is controlled.  Those arguments,
however, retain an alphabet-dependent base.  The present framework is meant to
separate a different question: whether a public transcript system can route all
large-alphabet information either to terminal projected shadows or to bounded
local ledgers, leaving the active recurrence alphabet-free.

The technical language used below is standard.  Entropy, relative entropy,
data processing, Pinsker-type estimates, and chain rules are used in the
usual information-theoretic sense \cite{Shannon48,KullbackLeibler51,
Pinsker64,CoverThomas06,CsiszarKorner11}; Renyi entropy enters only through
conditioned-copying estimates \cite{Renyi61}.  The random-coloring reduction
is a first-moment probabilistic-method step \cite{AlonSpencer16}, and the DRC
blocks are finite coordinate versions of dependent random choice
\cite{FoxSudakov11}.

\subsection*{Self-contained interface convention}

For the framework to be externally checkable, every statement used in the
conditional assembly must be one of the following:
\begin{enumerate}[label=\textup{(\arabic*)}]
\item stated and proved in this manuscript;
\item cited as a standard external theorem in the ordinary mathematical way.
\end{enumerate}
The elementary, structural, product, residual, old-support, and bookkeeping
interfaces used below are therefore stated explicitly.  The reduction from the
original set-family formulation to balanced two-collision codes is also
included, so that the conditional conclusion can be read in the original
sunflower language once the interface hypotheses have been verified.

\subsection*{Proposed fixed-window mechanism}

The candidate alphabet-free mechanism is the fixed, labelled hash window of
Theorem~\ref{thm:public-hash}.  In a usual random-restriction or encoding
argument, a useful coordinate is often found together with a value or a long
restriction transcript; exposing that value costs on the order of the
available alphabet entropy.  The proposed normal form separates the two operations.  A
high-mass value is terminalized immediately as a projected shadow and is paid
by pinned coordinates.  A low-collision coordinate is not exposed by value:
it is placed into a fixed hash window for one public sample label.  After
only $h$ surviving coordinates, with $h$ fixed once and for all, the window
returns a distinct near-sunflower block at cost
\[
  O(B_{\rm win}+h),
\]
including the public transcript, the Renyi conditioned-copy tax, and the exact
split exposure.  This cost depends on the fixed thresholds, not on the
alphabet sizes or on $k$.  Repetition is controlled separately by the
fresh/old/scale potentials, so the hash window is not an additional reusable
budget.  This is the proposed point at which alphabet-dependent value exposure
would be removed from the active recurrence, if the local hash and ledger
interfaces are fully verified.

\subsection*{Key mechanism trace}

The framework uses the fixed window only as one local progress source, not as a
replacement for all near-sunflower production.  The branch accounting is as
follows.

First, a positive-KL or collision output never exposes an alphabet value
unless it terminalizes as a projected shadow.  Otherwise it contributes one
low-collision coordinate to the labelled hash queue.  The hash queue has a
fixed cumulative transcript budget $B_{\rm win}$: prefix, tail/pruning, and
coordinate-production transcript are appended only while the cumulative cost
remains below that budget.  If the next set-valued tail/pruning increment
would exceed the budget, it is not appended; the branch terminalizes on the
already paid value prefix.  Prefixes which exceed the remaining budget are
terminal projected shadows and do not enter the hash queue.

Second, the normal-form hash window is applied only after enough distinct
coordinates for one public sample label have accumulated before the cumulative
budget closes.  If the budget closes first, the branch is a terminal or
bounded-support exit and is charged by the same fixed hash ledger.  In the
actual transition system the destruction budget in Theorem~\ref{thm:public-hash} is zero:
whenever a produced coordinate later has an atom of mass at least $\beta$,
the branch stops as a value shadow.  Thus a nonterminal window with at least
$2h$ produced coordinates has $h$ surviving coordinates deterministically;
the Markov form of Theorem~\ref{thm:public-hash} is kept only to make the
fixed-window statement robust.

Third, on those $h$ surviving coordinates the one-sample conditional law has
coordinate collision at most $\beta$ and codeword atoms below $\beta$.  Hence
three independent conditioned copies are distinct and have at most
$6\beta h$ split coordinates with fixed positive probability.  The Renyi
copying tax and the exact split exposure cost
\[
  O(B_{\rm win}+h),
\]
with no dependence on the alphabet sizes or on $|\C|$.

Finally, the returned split block is processed by the residual decision
order.  Fresh coordinates can be promoted only once.  Old incidences are
recorded by raw revisit counters and charged to the weighted debit
$N_{\rm old}$; the raw counters are capped by the high-revisit rule.  Scale-only
returns decrease active scale geometrically, and every new scale is born with
a paid certificate.  Consequently repeated use of fixed windows is paid by
the fresh, old, and scale potentials; there is no additional union bound over
the whole code family.

\begin{theorem}[Public normal-form interface reduction]
\label{thm:main-reduction}
Assume that the public-normal-form transition system satisfies the local
interfaces stated in Sections~\ref{sec:elementary}--\ref{sec:constant-hierarchy}:
public tuple exits, fixed-window hash returns, residual repair, old-window
projection, bounded reservoirs, pruning normalization, global charge summation,
and scale-birth certificates.  Then there is a constant $C<\infty$, depending
only on the fixed hierarchy of parameters, such that every balanced
two-collision code of length $k$ started from the root state of
Lemma~\ref{lem:root-state} has size at most $C^k$.
Consequently, a complete verification of these interfaces in the intended
sunflower setting would give that every rank at most $k$ set family with no
three-petal sunflower has size at most $(eC)^k$, after increasing the constant
for the finitely many small values of $k$.
\end{theorem}

\begin{proof}
Let $\C$ be a balanced two-collision code of length $k$ and let $\mu$ be the
uniform law on $\C$.  By Lemma~\ref{lem:root-state}, this gives an admissible
public root state with rank $k$, empty old support, and root potential
bounded by $O(k)$.

Run the normal-form transition system of Definition~\ref{def:transition-system}
with the constants chosen as in Section~\ref{sec:constant-hierarchy}.  Under
the stated interface hypotheses, every nonterminal mode has an admissible
public exit: close-pair exits are handled by
Theorem~\ref{thm:close-pair-exit}, product-KL exits by
Theorem~\ref{thm:product-kl-branch}, positive-KL outputs by
Theorem~\ref{thm:positive-kl-public}, public hash returns by
Theorem~\ref{thm:public-hash} and Propositions~\ref{prop:hash-ledger} and
\ref{prop:hash-accumulation}, residual repair by
Proposition~\ref{prop:residual-public} and
Lemma~\ref{lem:residual-support-promotion}, old-supported repair by
Lemma~\ref{lem:old-residual-routing} and
Proposition~\ref{prop:sparse-residual-transition}, old windows by
Theorem~\ref{hyp:old-window}, bounded reservoirs by
Proposition~\ref{prop:reservoir}, and active scale birth by
Proposition~\ref{prop:birth}.  Therefore Theorem~\ref{thm:assembly} gives a
terminal projected shadow whose entropy is controlled by pinned coordinates,
sample labels, pruning loss, and a telescoping potential.

Theorem~\ref{thm:shadow-potential} converts this entropy bound into a shadow
atom with the required mass.  The strengthened induction absorbs the
potential drop, and Theorem~\ref{thm:multi-shadow} gives
\[
  |\C|\le D^k\exp(A\Phi(\mathsf S_0)).
\]
Since the root has empty old support and
\[
  \Phi(\mathsf S_0)\le (B_{\rm rk}+B_f+B_s)k,
\]
the right-hand side is at most
\[
  \bigl(D\exp(A(B_{\rm rk}+B_f+B_s))\bigr)^k.
\]
Taking this quantity as $C^k$ proves the balanced-code bound.  Finally
Theorem~\ref{thm:balanced-code-reduction} transfers the bound back to
ordinary rank at most $k$ set families, with the factor $e^k$ coming from
random coloring.
\end{proof}

\begin{remark}
Theorem~\ref{thm:main-reduction} is a conditional bookkeeping theorem.  Its
role is to separate the global entropy accounting from the local mathematical
interfaces.  The rest of the manuscript records a proposed verification of
those interfaces in public normal form, replacing private terminal transcripts
by projected shadows, potential drops, pruning charges, or fixed public hash
windows.
\end{remark}

\section{Reduction to Balanced Two-Collision Codes}
\label{sec:front-reduction}

A three-petal sunflower is a triple of distinct sets
$A,B,C$ satisfying
\[
  A\cap B=A\cap C=B\cap C .
\]
The rank of a family is the maximum size of one of its members.

\begin{definition}[Balanced two-collision code]
\label{def:balanced-code}
A code $\C\subseteq\prod_{i=1}^k\Omega_i$ is a balanced two-collision code of
length $k$ if for every three distinct codewords $x,y,z\in\C$ there is a
coordinate $i$ such that exactly two of $x_i,y_i,z_i$ are equal.
\end{definition}

\begin{lemma}[Padding to uniformity]
\label{lem:padding-uniformity}
Let $\F$ be a rank at most $k$ family with no three-petal sunflower.  Then
there is a $k$-uniform family $\F^+$ with $|\F^+|=|\F|$ and no three-petal
sunflower.
\end{lemma}

\begin{proof}
For each $A\in\F$, add a private set $D_A$ of $k-|A|$ new vertices, disjoint
from all old vertices and from all $D_B$ with $B\ne A$, and set
$A^+=A\cup D_A$.  The family
$\F^+=\{A^+:A\in\F\}$ is $k$-uniform and has the same cardinality as $\F$.

If $A^+,B^+,C^+$ formed a sunflower, then every element in any pairwise
intersection of two of them would be an old vertex, because each dummy vertex
belongs to exactly one padded set.  Hence
\[
  A^+\cap B^+=A\cap B,\quad
  A^+\cap C^+=A\cap C,\quad
  B^+\cap C^+=B\cap C .
\]
Thus $A,B,C$ would form a sunflower in $\F$, a contradiction.
\end{proof}

\begin{lemma}[Random coloring to a partite code]
\label{lem:random-coloring-code}
Let $\F$ be a $k$-uniform family with no three-petal sunflower.  Then there
is a balanced two-collision code $\C$ of length $k$ with
\[
  |\C|\ge \frac{k!}{k^k}\,|\F| .
\]
\end{lemma}

\begin{proof}
This is the standard random-coloring first-moment argument; see
\cite{AlonSpencer16}.
Color the ground set of $\F$ independently and uniformly by colors
$1,\ldots,k$.  A fixed $k$-set is colorful, meaning that it receives each
color exactly once, with probability $k!/k^k$.  Therefore some coloring has a
colorful subfamily $\F_0\subseteq\F$ with
$|\F_0|\ge (k!/k^k)|\F|$.

For this coloring, let $\Omega_i$ be the set of vertices of color $i$ which
appear in $\F_0$.  Encode every $A\in\F_0$ by the word
$x(A)=(x_1(A),\ldots,x_k(A))$, where $x_i(A)$ is the unique vertex of color
$i$ in $A$.  This gives an injective map from $\F_0$ into
$\prod_i\Omega_i$; let $\C$ be its image.

Take three distinct words $x,y,z\in\C$ and their corresponding colorful sets
$A,B,C\in\F_0$.  In a fixed color class $i$, the three symbols
$x_i,y_i,z_i$ are either all equal, all distinct, or exactly two are equal.
If every coordinate is all equal or all distinct, then the contribution of
each color class to the three pairwise intersections is the same; hence
$A,B,C$ form a sunflower.  Since $\F_0$ is sunflower-free, some coordinate
has exactly two equal symbols.  Thus $\C$ is a balanced two-collision code.
\end{proof}

\begin{theorem}[Balanced-code reduction]
\label{thm:balanced-code-reduction}
Suppose that every balanced two-collision code of length $k$ has size at most
$C^k$ for an absolute constant $C$.  Then every rank at most $k$ family with
no three-petal sunflower has size at most $(eC)^k$.
\end{theorem}

\begin{proof}
By Lemma~\ref{lem:padding-uniformity}, it is enough to consider a
$k$-uniform sunflower-free family $\F$.  Lemma~\ref{lem:random-coloring-code}
gives a balanced two-collision code $\C$ of length $k$ with
\[
  |\C|\ge \frac{k!}{k^k}|\F|.
\]
The assumed code bound gives
\[
  |\F|\le \frac{k^k}{k!}C^k.
\]
Using $k!\ge (k/e)^k$, we obtain $|\F|\le (eC)^k$.
\end{proof}

\section{Encoded Codes and Balanced States}

We work with a finite ambient code
\[
  \C \subseteq \prod_{i=1}^k\Omega_i
\]
and with public active coordinate blocks $I\subseteq [k]$.  A random codeword
under a law $\mu$ is denoted by $X$.
At a recursive cylinder state the pinned coordinates have been removed from
the ambient product, and the remaining coordinates are relabelled as the
current ambient set.  Thus every later occurrence of $[k]$ inside a state,
including $[k]\setminus U$ in the potential, means the current unpinned ambient
coordinate set, not necessarily the original root set.

\begin{definition}[Admissible two-collision state]
A state is a tuple
\[
  \mathsf S=(I,\C,\mu,U,\mathcal P,\mathcal R),
\]
where $I$ is the active coordinate set, $U\subseteq [k]$ is the global old
support,
$\mathcal P$ is the public state, and $\mathcal R$ is a bounded active
reservoir of sampled codewords.  The code $\C$ is a balanced two-collision
code after the public value pins already present in $\mathcal P$.  Heavy
coordinate fibres are immediate exits of the proof tree, not additional
entry assumptions.
\end{definition}

The final proof must be insensitive to analysis-only data.  Thus we separate
public state from private transcript.  Formally, each node $v$ carries a
public sigma-algebra $\mathcal F_v$ generated by the fixed constants, the root
coordinate order, and all previously paid transcript variables.  A block,
law, tie-break, or transition selector is public at $v$ if it is
$\mathcal F_v$-measurable.  If a non-public child selector $\Gamma_v$ is
revealed, then the tree ledger pays $\bfH(\Gamma_v\mid E_v)$ and the child
public sigma-algebra is
\[
  \mathcal F_{v'}=\mathcal F_v\vee \sigma(\Gamma_v)
\]
on that child event.  Variables used only to analyze the branch are not
adjoined to the terminal sigma-algebra unless they are paid in this way or
are terminalized as part of a projected shadow.

\begin{definition}[Public normal form]
A proof node is in public normal form if:
\begin{enumerate}[label=\textup{(\roman*)}]
\item every coordinate block used for a hash, DRC, close-pair, or
near-sunflower transition is $\mathcal F_v$-measurable after the already paid
public transcript at that node;
\item every value pin terminalizes immediately as a projected value-shadow;
if the pin is selected by a high-mass existence rule rather than by exposing a
full value selector, the retained atom is chosen by the fixed public tie-break
among atoms above the threshold and the discarded complement is charged as a
pruning-normalization loss; no part of that complement is inherited by a
nonterminal child;
\item every DRC block is selected by a deterministic tie-break from the
current public law;
\item every hash window carries one public active sample label, and all hash
coordinates accumulated in that window are ambient coordinates of that same
sample label under the paid tuple event; the fixed budget $B_{\rm win}$ is a
per-labelled-window budget, not a global per-branch budget, and a different
sample label opens a different window with its own stopped transcript;
\item every old-window queue carries one public sample-label signature, namely
the ordered active triple/reservoir labels whose equality pattern is being
recorded; old-pattern atoms are nested only while this signature is unchanged;
\item nonterminal nodes carry only a bounded active reservoir;
\item every conditional one-sample law used at a tuple node is the law of an
active sample label under a paid tuple event, and every conditioned copy is
realized back in independent root samples by Lemma~\ref{lem:fixed-copy} or
Lemma~\ref{lem:random-copy} before it can contribute to a terminal shadow;
\item private analysis data either never enters the terminal sigma-algebra, is
charged to a potential drop or pruning event, or is terminalized into a
projected shadow.
\end{enumerate}
\end{definition}

\begin{lemma}[Root state]
\label{lem:root-state}
Let $\C\subseteq\prod_{i=1}^k\Omega_i$ be a balanced two-collision code and
let $\mu$ be the uniform law on $\C$.  Then
\[
  \mathsf S_0=([k],\C,\mu,\varnothing,\mathcal P_0,\varnothing)
\]
is an admissible public root state, where $\mathcal P_0$ contains only the
fixed coordinate order and the fixed threshold constants.
\end{lemma}

\begin{proof}
There are no old coordinates, no private transcript, and no active reservoir
at the root.  The coordinate order and constants are deterministic, hence
public.  The defining two-collision property of $\C$ is exactly the
admissibility condition before any public value pins are made.
\end{proof}

\begin{lemma}[Root activation dichotomy]
\label{lem:root-activation}
Let $\mu$ be a public law on a balanced two-collision code on an active block
$I$, $|I|=m$.  Fix constants $0<\alpha<1$ and $0<\tau<1$.  For independent
$X,Y,Z\sim\mu$, at least one of the following occurs:
\begin{enumerate}[label=\textup{(\roman*)}]
\item terminal atom: either a codeword has $\mu$-mass at least $1/12$, or
some coordinate value has $\mu(X_i=a)\ge\tau$;
\item near-sunflower activation:
  \[
    \Prb\!\left[X,Y,Z\text{ distinct and }
      |\Split_I(X,Y,Z)|\le 6\alpha m\right]\ge 1/4;
  \]
\item close-pair activation: for every integer
  \[
    1\le s\le (1-\tau)\alpha m/4
  \]
  there is a public block $J\subseteq I$, $|J|=s$, such that
  \[
    \Prb[X_J=Y_J,\ Z_j\ne X_j\ (j\in J)]
    \ge
    \frac{(1-\tau)\alpha}{2}
    \left(\frac{(1-\tau)\alpha}{4}\right)^s .
  \]
\end{enumerate}
\end{lemma}

\begin{proof}
If the terminal atom alternative holds, choose the first such atom in the fixed
public order, retain that atom as a projected terminal shadow, and charge the
constant pruning-normalization loss of the retained atom as in the public
normal-form convention.  No complement branch is inherited by a nonterminal
child.  Assume from now on that no terminal atom alternative holds.  In
particular,
\[
  \max_x\mu(x)<1/12,\qquad
  \max_a\Prb[X_i=a]<\tau\quad(i\in I).
\]

If the average collision on $I$ is at most $\alpha$, then
Lemma~\ref{lem:collision-split} gives
\[
  \Prb[|\Split_I(X,Y,Z)|\le 6\alpha m]\ge 1/2.
\]
The probability that $X,Y,Z$ are not all distinct is at most
$3\sum_x\mu(x)^2\le 3\max_x\mu(x)<1/4$.  Hence the near-sunflower activation
has mass at least $1/4$.

It remains to consider the case
\[
  \frac1m\sum_{i\in I}\Coll_i(\mu)>\alpha.
\]
For a coordinate $i$, write $p_a=\Prb[X_i=a]$.  Then
\[
  \Prb[X_i=Y_i\ne Z_i]
  =
  \sum_a p_a^2(1-p_a)
  \ge
  (1-\tau)\sum_a p_a^2
  =
  (1-\tau)\Coll_i(\mu).
\]
Thus the events $E_i=\{X_i=Y_i\ne Z_i\}$ have average mass at least
$(1-\tau)\alpha$.  Apply Lemma~\ref{lem:drc-row-block} to the probability
space of triples and the events $E_i$.  The block is chosen by the fixed
public tie-break among all valid DRC blocks, so it is public.
\end{proof}

\begin{definition}[Normal-form transition system]
\label{def:transition-system}
The proof tree is the smallest adaptive tree generated by the following
public modes.  A fixed integer $m_0$ is chosen in the constant hierarchy; any
node whose active rank and all active block sizes are $<m_0$ is declared a
bounded leaf.
\begin{enumerate}[label=\textup{(\arabic*)}]
\item Inactive code-law mode.  The node has an admissible public state but no
active block.  Apply Lemma~\ref{lem:root-activation}.  Terminal atoms stop,
near-sunflower activations enter residual mode, and close-pair activations
enter close-pair mode after choosing, for active rank $m\ge m_0$, the first
valid DRC block of size
\[
  \left\lfloor (1-\tau)\alpha m/8\right\rfloor .
\]
\item Close-pair mode.  The node carries a public block $J$ and a tuple law
with $X_J=Y_J=A$ and $Z_j\ne A_j$ on $J$.  Apply
Theorem~\ref{thm:close-pair-exit}.  Low-entropy and atom exits stop as
projected shadows; KL/collision exits route through terminal projected tuple
pins, charged pruning, or hash mode; the law-level near-sunflower exit enters
residual mode.
\item Hash-window mode.  The node carries a public sample label, a public hash
transcript prefix whose length is bounded by the per-window budget
$B_{\rm win}$, and a bounded list of public low-collision coordinates for
that label.  Each active sample label has its own window; a coordinate produced
for another label is routed to that label's separate stopped window, subject to
the same per-window budget and the single-token rules below.  Apply
Propositions~\ref{prop:hash-ledger} and \ref{prop:hash-accumulation}, and
Theorem~\ref{thm:public-hash}.  The
output is a terminal value shadow, a bounded-support exit, or a public
near-sunflower block which enters residual mode.
\item Residual mode.  The node carries a public near-sunflower block and an
exact split exposure.  Apply the public order of
Proposition~\ref{prop:residual-decision-order}.  The output is fresh
progress, scale descent, a close-pair block, a bounded leaf, or an old-window
block.
\item Old-window mode.  The node carries public old blocks of one scale and one
public sample-label signature.
Apply Theorem~\ref{hyp:old-window}.  The output is a terminal low-symbol
shadow, a high-revisit terminal/pruning/hash exit, or a low-revisit
continuation whose non-value transcript has been registered in the
old-incidence ledger.
Bounded short windows remain in old-window mode through the delayed queue.
\item Terminal mode.  A terminal node outputs a projected shadow and projects
away all unpaid analysis-only data before the final entropy count.
\item Bounded-leaf mode.  A bounded leaf is counted by the crude constant
$D_0^{\rank}$ after $D_0$ is chosen large enough for the finitely many ranks
and block sizes below $m_0$.  The required finite bounds come from the
classical Erdos--Rado sunflower bound in these finitely many fixed ranks.
\end{enumerate}
All coordinate and block choices in these rules are deterministic functions
of the public state and the fixed thresholds, except for value atoms that are
immediately terminalized.
\end{definition}

\begin{proposition}[Formal mode and exit table]
\label{prop:formal-mode-table}
The nonterminal part of Definition~\ref{def:transition-system} is equivalently
the following priority-ordered table.  Each row is read from left to right; the
first satisfied exit in the active mode is taken, and all block/value choices
inside the selected exit use the fixed public order.
\begin{center}
\small
\begin{tabular}{p{0.14\linewidth}p{0.24\linewidth}p{0.35\linewidth}p{0.16\linewidth}}
mode & entry data & priority exits &
charge source\\
\hline
inactive &
public code law, no active block &
terminal atom; low-collision near-sunflower residual; DRC close-pair block &
shadow; rank/DRC\\
close-pair &
public $J$, $X_J=Y_J=A$, $Z_J=B\ne A$ &
low entropy/atom; cross/high collision; product-KL branch; law-level residual &
shadow; pruning/hash; residual\\
hash window &
one public sample label, stopped transcript, live coordinate list &
value pin; prefix overflow; bounded support; fixed-window return &
shadow; fresh/old; hash token\\
residual &
public near-sunflower block and exact split set &
fresh promotion; deletion-trapped scale descent; old-supported route; close-pair birth; bounded leaf &
fresh; scale; old; rank/birth\\
old window &
old blocks of one scale and one sample-label signature &
low-symbol terminal; high-revisit terminal/pruning/hash; low-revisit registration &
shadow; pruning/hash; old
\end{tabular}
\end{center}
Terminal and bounded-leaf exits have no nonterminal child.  A nonterminal hash
return is the only row which creates a pending token, and
Lemma~\ref{lem:token-free-children} forces that token to be assigned before any
descendant state is entered.
\end{proposition}

\begin{proof}
This is a reindexing of the mode definitions.  The entry columns are disjoint:
inactive states have no active block, close-pair states carry a certified pair
block, hash states carry one selected labelled stopped window, residual states
carry an exact split block, and old-window states carry old blocks of one scale
and one sample-label signature.  Within a row, the displayed order is precisely
the decision order in the cited local lemma or theorem.  Terminal atoms have
highest priority and stop; otherwise the remaining alternatives are complements
of the previous tests in that row.  Thus the selected exit is a deterministic
function of the public state together with the realized value atom on terminal
branches, and every nonterminal child is placed in the child mode shown by the
row.
\end{proof}

\begin{proposition}[Transition system is public and exhaustive]
\label{prop:transition-system}
Every nonterminal node generated by Definition~\ref{def:transition-system}
is in public normal form and is in exactly one of the five nonterminal modes
listed there.  The entry conditions for these modes are disjoint:
inactive nodes have no active block, close-pair nodes carry a certified pair
block, hash-window nodes process one selected sample-label window and no pending
returned hash token, residual nodes carry an exact split block returned by
activation, hash, or close-pair, and old-window nodes carry old blocks of one
scale and one sample-label signature.  If several labelled hash windows are
available, the next active window is chosen by the fixed public order; each
labelled window has its own stopped transcript and its own possible token.  The
listed exit alternatives for the active mode form an exhaustive partition after
deterministic public tie-breaks.  Every child is either terminal or again
belongs to one of those modes.
\end{proposition}

\begin{proof}
Induct on tree depth.  The root is public by Lemma~\ref{lem:root-state}.  At
each step the cited local result is used as an exhaustive list of children, not
as a diagnostic after the child has been chosen.
The mode and exit partition is the one in Proposition~\ref{prop:formal-mode-table}.

In inactive mode, Lemma~\ref{lem:root-activation} uses only the current public
law and fixed tie-breaks.  Its alternatives are terminal atom, residual
near-sunflower, and close-pair activation; if the terminal atom condition holds
it is taken first, otherwise the collision average dichotomy separates the
near-sunflower and close-pair cases.  The chosen close-pair block has the
fixed-fraction size displayed in Definition~\ref{def:transition-system}; the
choice of $m_0$ makes this size nonzero and within the DRC range of
Lemma~\ref{lem:root-activation}.  In close-pair mode, the input block
is public by the mode definition and Theorem~\ref{thm:close-pair-exit} returns
only terminal low-entropy/atom exits, product-KL exits routed by
Theorem~\ref{thm:product-kl-branch}, collision exits routed
exhaustively by Lemma~\ref{lem:collision-exit-pin}, or a law-level near-sunflower residual
return; the theorem's decision order makes these exits a partition.  In
hash-window mode, the active state is one selected labelled window.  Coordinates
for another sample label are placed in that label's separate stopped window and
are not merged into the current transcript.  Theorem~\ref{thm:public-hash}
requires coordinate production to be public for the fixed label of the active
window and pays that window's public transcript before any copying, while
Propositions~\ref{prop:hash-ledger} and \ref{prop:hash-accumulation} account
for terminal pins, bounded support, survival scans, failed accumulation, and
residual returns.

In residual mode, the exact split block is already public.  The residual
decision order first tests fresh support, then old support, then sparse
scale/close-pair alternatives; DRC blocks are public by
Definition~\ref{hyp:public-drc}, fresh support promotion is public by
Lemma~\ref{lem:residual-support-promotion}, and the remaining old-supported
witnesses are routed by Lemma~\ref{lem:old-residual-routing} and
Proposition~\ref{prop:sparse-residual-transition} to scale, old-window, or
close-pair modes.  In old-window mode, Theorem~\ref{hyp:old-window} either
terminalizes a low-symbol old value, exits by high revisit, or registers all
non-value old data before continuation; Proposition~\ref{prop:delayed-old-queue}
delays only bounded short windows of the same public scale and signature.
Terminal nodes keep only projected shadows, and bounded leaves are counted by
the finite constant $D_0$.

The tuple-node realization invariant is preserved because value pins stop
immediately, while the only operation which creates new conditioned samples is
conditioned copying, and Lemmas~\ref{lem:fixed-copy} and
\ref{lem:random-copy} realize those copies in independent root samples with
the displayed transcript cost.  Thus public normal form and the mode
classification are preserved.
\end{proof}

\begin{proposition}[Branch invariants]
\label{prop:branch-invariants}
On every retained branch of the normal-form transition system the following
invariants hold.
\begin{enumerate}[label=\textup{(\roman*)}]
\item The public sigma-algebras are increasing, and every increase is generated
by a child selector whose conditional entropy is charged in the proof-tree
ledger.
\item Each labelled hash window has its own stopped transcript budget
$B_{\rm win}$ and one public sample label.  Different sample labels have
different windows.  The global branch is not charged by a shared hash transcript
budget; it is charged by the number of completed or bounded labelled windows,
each of which is later terminalized or assigned to one hard exit.
\item Every tuple law is the law of a fixed finite list of independent root
samples conditioned on the paid transcript accumulated on the branch.  New
conditioned samples are introduced only through Lemmas~\ref{lem:fixed-copy}
and \ref{lem:random-copy}.
\item The old support $U$ is monotone increasing.  Fresh promotions are
disjoint subsets of $[k]\setminus U$, and each promoted coordinate is removed
from the fresh potential exactly once.
\item Old-window incidences increase the raw revisit counters monotonically
until the high-revisit exit is triggered, and each registered incidence debits
the weighted old account $N_{\rm old}$.  No low-revisit continuation uses an
incidence which has not been registered and debited.
\item A hash return carries at most one pending token, of value at most
$C_{\rm hash}$, attached to the returned block.  The transition system processes
that returned block until the token is assigned to the first subsequent
terminal, fresh, old, scale, or birth exit; before this assignment no child
state, descendant block, or new hash window is entered from that block.  If the
first exit is deletion-trapped scale descent or a birth transition, the token is
assigned at the parent residual transition and is not inherited by the child.
\item Every active block is born with the certificate of
Definition~\ref{hyp:birth}.  After birth, deletion-trapped returns decrease
its scale by the fixed factor in Lemma~\ref{lem:deletion-scale-descent}; no
transition increases active scale without a new birth certificate.
\item Old-pattern atoms are nested under growth of $U$ for a fixed sample-label
signature.  When $U$ is enlarged to $U^+$ without changing that signature, any
later pattern atom on $U^+$ is refined inside the previously paid atom on $U$
and exposes only the pattern on $U^+\setminus U$ as new data.  If the active
triple/reservoir signature changes, every bounded old queue for the old
signature is first registered, terminalized, or charged to the hard exit which
changes the signature, and then cleared.
\end{enumerate}
\end{proposition}

\begin{proof}
The first two invariants are exactly the public-normal-form definition and the
copying clauses in Proposition~\ref{prop:transition-system}.  Fresh promotion
is performed only by Lemma~\ref{lem:residual-support-promotion} and by bounded
hash-support exits, both of which add public subsets of the current complement
$[k]\setminus U$ to $U$; hence the promoted sets are disjoint along a branch.
Old incidences are registered in Lemma~\ref{lem:old-residual-routing},
Theorem~\ref{hyp:old-window}, and Proposition~\ref{prop:delayed-old-queue};
these are precisely the old-window continuations.  Each registration increments
the raw revisit counter and debits the weighted old account, while the
high-revisit exit stops the low-revisit regime when the raw counter reaches the
cap.
Proposition~\ref{prop:hash-ledger} and Lemma~\ref{lem:token-free-children}
give the single-token and token-free-child rules for hash returns.  Finally,
Proposition~\ref{prop:birth} enumerates the sources of
active blocks and assigns their certificates, while
Lemma~\ref{lem:deletion-scale-descent} gives the only scale-decreasing
nonterminal continuation.  These are the only nonterminal modes by
Proposition~\ref{prop:transition-system}.  The nested old-pattern rule is a
normal-form convention for refining public sigma-algebras for the same active
sample-label signature: a pattern already paid on $U$ remains part of the public
state for that signature, and a later pattern on $U^+$ is the intersection of
that old atom with a newly paid pattern on $U^+\setminus U$.  Signature-changing
transitions cannot inherit that atom; Proposition~\ref{prop:delayed-old-queue}
flushes the bounded unpaid queue before such a change is entered.
\end{proof}

\begin{proposition}[Finite truncation]
\label{prop:finite-truncation}
After bounded leaves are imposed at scale $m_0$, every branch of the
normal-form transition system has a finite truncation whose discarded tail has
total mass bounded by the pruning estimates already charged in the transition
system.
\end{proposition}

\begin{proof}
The only zero-entropy transitions are deterministic public choices: fixed
tie-breaks for blocks, mode changes after a charged selector has already been
revealed, and bookkeeping relabellings such as unchanged-signature old-queue
growth.  Such a choice cannot repeat in the same state.  It either enters a
terminal/bounded leaf, enters a hash window with a fixed remaining transcript
budget, enters an old queue whose cumulative incidence is counted, or enters an
active block with a rank/scale certificate.  Hence every infinite retained path
would have to contain infinitely many charged or monotone-budget steps of one
of the types listed below.

Fresh promotions and fresh/rank descents decrease
$\Phi_{\rm fresh}+\Phi_{\rm old}$ by a fixed positive amount per promoted
coordinate.  Old-window incidences are either registered and debited in
$N_{\rm old}$, which decreases $\Phi_{\rm old}$ by the weighted debit, or they
trigger the high-revisit exit of
Theorem~\ref{hyp:old-window}; the delayed queue prevents bounded short windows
from restarting without increasing the same cumulative counters.
Hash windows have a fixed stopped transcript budget $B_{\rm win}$ and fixed
target size $2h$: before reaching that target they either close as terminal or
bounded-support exits, and after reaching it they return to residual mode with
a token which must be assigned before any descendant window is opened.  Thus
hash mode cannot contribute infinitely many zero-cost continuations.  Scale
descents are geometric by Lemma~\ref{lem:deletion-scale-descent}; once
the scale falls below $m_0$ the branch is a bounded leaf.  Every new scale is
born with a certificate by Proposition~\ref{prop:birth}, so the total
geometric scale available along a branch is bounded by the initial active-rank
or parent active-scale account, fresh-promotion drops, and weighted old-debit
budgets.  Bounded hash
tokens pay their fixed transcript charge only after assignment to a hard exit;
they do not create recurrence potential.  Terminal pins and public DRC
certificates do not create recurrence potential.

The only discarded branches in the finite truncation are explicit pruning
branches, such as rare positive-tail pruning, first-crossing light-complement
pruning from Lemma~\ref{lem:heavy-or-light-hash}, high-mass terminal
atom pruning from the public normal-form convention, hash-survival pruning
from Definition~\ref{def:hash-survival}, and the public fresh-reset pruning of
Lemma~\ref{lem:fresh-reset-pruning}.  Their total discarded mass is
bounded by the estimates charged at the pruning step.  In particular,
positive-tail and light-complement pruning are charged inside the same
prefix/tail transcript which then immediately produces a terminal pin or a hash
coordinate.  If the fixed tail/pruning increment would overflow the window, it
is not retained as a set-valued terminal shadow; the branch stops on the value
prefix already paid by the KL chain rule.  Thus repeated prunings consume the
fixed hash-window budget or terminate as projected value-shadows.  Therefore
truncating after all potential-drop, incidence, scale, and bounded-window
budgets up to a fixed level discards only already charged pruned mass.
\end{proof}

\begin{proposition}[Proof-tree entropy ledger]
\label{prop:tree-entropy-ledger}
Let $\mathcal T$ be a finite adaptive proof tree.  Let $\Gamma$ be its full
terminal leaf transcript, and let $\Theta=f(\Gamma)$ be the projected
terminal shadow obtained after all unpaid analysis-only data are projected
away.  For an
internal node $v$, let $E_v$ be the event that the branch reaches $v$, and
let $A_v$ be the child selector at $v$.  Then
\[
  \bfH(\Theta)
  \le
  \bfH(\Gamma)
  =
  \sum_{v\in\mathcal T^\circ}
  \Prb(E_v)\,\bfH(A_v\mid E_v).
\]
Consequently, if nonnegative local charges $c_v$ satisfy
\[
  \bfH(A_v\mid E_v)
  \le
  \EE[c_v\mid E_v]
  \qquad(v\in\mathcal T^\circ),
\]
then
\[
  \bfH(\Theta)
  \le
  \EE\sum_{v\ {\rm reached}} c_v .
\]
\end{proposition}

\begin{proof}
Since $\Theta$ is a function of $\Gamma$, the data-processing inequality gives
$\bfH(\Theta)\le\bfH(\Gamma)$.  Reveal the terminal leaf by scanning the
finite tree in a fixed topological order, so every parent is scanned before its
children, and when a reached internal node is met, reveal its child selector.
Selectors at unreached nodes are deterministic cemetery symbols.  The chain rule gives
\[
  \bfH(\Gamma)
  =
  \sum_{v\in\mathcal T^\circ}
  \bfH(A_v\mathbf 1_{E_v}\mid \text{previous selectors}).
\]
Conditioned on the previous selectors, the event $E_v$ is already known.
Thus the contribution is zero off $E_v$ and is
$\bfH(A_v\mid E_v)$ on $E_v$, giving the displayed tree identity.  Multiplying
the local charge inequality by $\Prb(E_v)$ and summing over $v$ gives the last
claim.
\end{proof}

\begin{proposition}[Pruning normalization charge]
\label{prop:pruning-normalization}
Suppose a finite proof tree is run with charged pruning steps.  At a pruning
step the current conditional law is restricted to a retained event of conditional
mass $q_v>0$, the discarded complement is not inherited by any recursive child,
and the branch records the loss
\[
  \ell_v=-\log q_v .
\]
Let $\widehat{\Prb}$ be the normalized law on retained terminal outcomes, let
\[
  L_{\rm pr}=\sum_{v\ {\rm pruned\ on\ the\ branch}}\ell_v,
\]
and let $\Prb_{\rm orig}$ denote the original, unnormalized mass of the retained
terminal shadow event before pruning projection.  Then
\[
  \widehat{\EE}\bigl[-\log \Prb_{\rm orig}(\Theta)\bigr]
  \le
  \bfH_{\widehat{\Prb}}(\Theta)+\widehat{\EE}L_{\rm pr}.
\]
Consequently, if
\[
  \bfH_{\widehat{\Prb}}(\Theta)+\widehat{\EE}L_{\rm pr}
  \le
  \widehat{\EE} B
\]
for a nonnegative branch charge $B$, then some retained terminal outcome
$\omega$ with shadow $\theta=\Theta(\omega)$ satisfies
\[
  -\log \Prb_{\rm orig}(\Theta=\theta)\le B(\omega).
\]
\end{proposition}

\begin{proof}
For a shadow atom $\theta$, the original retained mass is
\[
  \Prb_{\rm orig}(\Theta=\theta)
  =
  \widehat{\Prb}(\Theta=\theta)\,
  \widehat{\EE}\!\left[e^{-L_{\rm pr}}\mid \Theta=\theta\right].
\]
Therefore Jensen's inequality gives
\[
  -\log \Prb_{\rm orig}(\Theta=\theta)
  \le
  -\log \widehat{\Prb}(\Theta=\theta)
  +
  \widehat{\EE}\!\left[L_{\rm pr}\mid \Theta=\theta\right].
\]
Averaging over $\theta$ proves the displayed expectation bound.  Applying the
same averaging argument to
$-\log\Prb_{\rm orig}(\Theta)-B$ gives the final atom.
\end{proof}

\section{Multi-Sample Projected Shadows}

\begin{definition}[Multi-sample value-shadow]
\label{def:multi-sample-shadow}
At the recurrence interface, let $X^1,\ldots,X^r$ be independent samples from
the uniform law on the code currently being bounded.  Internal conditional
tuple laws used by the proof tree are auxiliary devices; every terminal shadow
is evaluated after pulling its event back to these root samples.  A
multi-sample projected value-shadow is a finite list
\[
  \theta=((J_s,a_s))_{s\in S},
\]
where $\varnothing\ne S\subseteq [r]$, $\varnothing\ne J_s\subseteq I$, and
$a_s\in\prod_{j\in J_s}\Omega_j$.  A terminal event $E$ is compatible with
$\theta$ if
\[
  E \subseteq \bigcap_{s\in S}\{X^s_{J_s}=a_s\}.
\]
Define
\[
  T(\theta)=\sum_{s\in S}|J_s|,\qquad R(\theta)=|S|.
\]
\end{definition}

\begin{theorem}[Multi-sample shadow descent]
\label{thm:multi-shadow}
Assume by induction that all ranks $<k$ satisfy the bound $B_j\le D^j$.
Let $\C$ be a rank-$k$ balanced two-collision code with uniform law $\mu$, and
let $\theta=((J_s,a_s))_{s\in S}$ be a nonempty terminal shadow value for
independent root samples from $\mu$.  Let $E_\theta$ be an event compatible with
$\theta$ and suppose it has mass
\[
  \delta=\Prb[E_\theta]
  \ge \exp(-\lambda T(\theta)-cR(\theta)).
\]
If $\log D\ge \lambda+c$, then $|\C|\le D^k$.
\end{theorem}

\begin{proof}
Let $M=|\C|$.  For each sample label $s$, let
\[
  \eta_s=\mu(\{x:x_{J_s}=a_s\}).
\]
Since the samples are independent and the terminal event is contained in the
intersection of the cylinders,
\[
  \delta\le \prod_{s\in S}\eta_s.
\]
The cylinder $\{x:x_{J_s}=a_s\}$ is again a balanced two-collision code after
the coordinates in $J_s$ are deleted: if three remaining codewords are
distinct, the original balanced property gives a $2+1$ coordinate, and it
cannot be one of the fixed coordinates.  Hence the cylinder has rank at most
$k-|J_s|$, and
\[
  \eta_s M\le D^{k-|J_s|}.
\]
Multiplying over $s$ gives
\[
  M^{R(\theta)}
  \le \delta^{-1}D^{R(\theta)k-T(\theta)}
  \le \exp(\lambda T(\theta)+cR(\theta))
     D^{R(\theta)k-T(\theta)}.
\]
Since each nonempty sample label contributes at least one pinned coordinate,
$T(\theta)\ge R(\theta)$.  If $\log D\ge\lambda+c$, the exponential factor is
at most $D^{T(\theta)}$, and so $M^{R(\theta)}\le D^{R(\theta)k}$.
Because $R(\theta)\ge1$, taking $R(\theta)$-th roots gives $M\le D^k$.
\end{proof}

\section{The Shadow Entropy Criterion}

The proof tree should not count full leaves.  It should merge leaves with the
same projected shadow and count only the resulting random variable.

\begin{theorem}[Shadow entropy criterion]
\label{thm:shadow-entropy}
Let $\Theta$ be the terminal projected shadow after all analysis-only data is
marginalized, and suppose every retained terminal outcome has a nonempty
shadow.  Suppose
\[
  \bfH(\Theta)
  \le
  \EE[\lambda T(\Theta)+cR(\Theta)].
\]
Then some shadow value $\theta$ satisfies
\[
  -\log \Prb[\Theta=\theta]
  \le
  \lambda T(\theta)+cR(\theta),
\]
and Theorem~\ref{thm:multi-shadow} applies.
\end{theorem}

\begin{proof}
For $I(\Theta)=-\log\Prb[\Theta]$ one has
$\EE I(\Theta)=\bfH(\Theta)$.  Hence
\[
  \EE[I(\Theta)-\lambda T(\Theta)-cR(\Theta)]\le 0.
\]
At least one atom has nonpositive value of the random variable inside the
expectation.
\end{proof}

We shall use a potential version.

\begin{theorem}[Stateful shadow entropy with potential]
\label{thm:shadow-potential}
Let a parent public state $\mathsf S$ have rank $k$ and potential
$\Phi_{\rm par}$.  Along a retained terminal branch let $\Delta\Phi$ be the
nonnegative potential expenditure charged by the proof tree after all hash
tokens, birth deposits, old queues, and scale charges on that branch have been
assigned, and let $L_{\rm pr}$ be the accumulated charged pruning-normalization
loss of Proposition~\ref{prop:pruning-normalization}.
Suppose every retained terminal outcome has a nonempty projected shadow.  If
\[
  \bfH(\Theta)+\EE L_{\rm pr}
  \le
  \EE[\lambda T(\Theta)+cR(\Theta)+A\Delta\Phi],
\]
then some retained terminal outcome gives a shadow atom satisfying the
corresponding pointwise inequality with respect to its original pre-pruning
mass.  For a
shadow value
\[
  \theta=((J_s,a_s))_{s\in S},
\]
let $\mathsf S_{\theta,s}$ be the cylinder child state obtained from the
terminal public ledgers on the branch by restricting the $s$-th root sample to
$x_{J_s}=a_s$ and deleting those pinned coordinates from the current ambient
coordinate set, all active supports, old supports, queues, and incidence
counters.  The underlying code of $\mathsf S_{\theta,s}$ is the full parent
value cylinder $\{x\in\C:x_{J_s}=a_s\}$ with $J_s$ deleted; non-value
analysis atoms from the terminal leaf are inherited only as ledger/support data
and do not further shrink this cylinder.  Assume the charged branch and the
inherited cylinder children
satisfy the recurrence-compatibility
inequality
\[
  \Delta\Phi+\sum_{s\in S}\Phi(\mathsf S_{\theta,s})
  \le R(\theta)\Phi_{\rm par}.
\]
If the induction hypothesis is strengthened to
\[
  B(\mathsf S)\le D^{\rank(\mathsf S)}\exp(A\Phi(\mathsf S)),
\]
and $\log D\ge\lambda+c$, then the potential term is compatible with the
multi-shadow descent.
\end{theorem}

\begin{proof}
Set
\[
  I_{\rm orig}(\Theta)=-\log\Prb_{\rm orig}[\Theta],
\]
where $\Prb_{\rm orig}$ is the retained shadow mass in the original
pre-pruning probability space.  By
Proposition~\ref{prop:pruning-normalization},
\[
  \EE I_{\rm orig}(\Theta)\le \bfH(\Theta)+\EE L_{\rm pr}.
\]
Hence
\[
  \EE\!\left[
    I_{\rm orig}(\Theta)-\lambda T(\Theta)-cR(\Theta)-A\Delta\Phi
  \right]\le 0,
\]
so some retained terminal outcome $\omega$ with shadow
$\theta=\Theta(\omega)$ satisfies
\[
  -\log\Prb_{\rm orig}[\Theta=\theta]
  \le
  \lambda T(\theta)+cR(\theta)+A\Delta\Phi(\omega).
\]

For such an atom the shadow descent proof is the stateful version of
Theorem~\ref{thm:multi-shadow}.  Write $R=|S|$.  For each sample label $s$,
the cylinder child $\mathsf S_{\theta,s}$ has as underlying code the full
value cylinder $\{x:x_{J_s}=a_s\}$, with the terminal public ledgers inherited
only as accounts after deleting $J_s$.  Thus the usual cylinder probability
bound is unchanged, and the child is a balanced two-collision code of rank at
most $k-|J_s|$.  The induction hypothesis gives
\[
  |\C_{\theta,s}|
  \le
  D^{k-|J_s|}\exp(A\Phi(\mathsf S_{\theta,s})).
\]
Multiplying the usual cylinder estimates gives the same inequality as in
Theorem~\ref{thm:multi-shadow}, but with the additional factor
\[
  \exp\!\left(A\Delta\Phi
       + A\sum_{s\in S}\Phi(\mathsf S_{\theta,s})\right).
\]
By the recurrence-compatibility condition,
\[
  \Delta\Phi+\sum_{s\in S}\Phi(\mathsf S_{\theta,s})
  \le
  R\Phi_{\rm par}.
\]
Hence, after taking $R$-th roots, the strengthened induction closes with the
parent factor $\exp(A\Phi_{\rm par})$.  The child in this paragraph is the
cylinder state carrying the terminal public ledgers; it is not an intermediate
transition child and it is not a freshly reset root state.

The drop $\Delta\Phi$ and the inherited child potentials need not be functions
of the projected shadow.  The averaging argument is applied on the underlying
terminal outcome space: it produces an outcome $\omega$ with shadow
$\theta=\Theta(\omega)$ and the displayed inequality using the branch ledgers
of that outcome.  The event $\{\Theta=\theta\}$ is still contained in the same
value cylinders defining the full cylinder children above, because all
analysis-only data has been projected away.  Therefore replacing the full leaf
by the merged shadow only increases the compatible cylinder event mass and does
not change the child code sets used in the induction.
\end{proof}

\section{Elementary Lemmas Used Throughout}
\label{sec:elementary}

This section records the elementary tools which appear repeatedly in the
argument.  We use the standard discrete entropy and relative-entropy
conventions from information theory \cite{Shannon48,CoverThomas06,
CsiszarKorner11}; the order-$m$ Renyi entropy used in copying lemmas is the
classical one \cite{Renyi61}.

\begin{lemma}[Entropy atom lemma]
\label{lem:entropy-atom}
Let $Y$ be a finite-valued random variable.  There is a value $y$ with
positive mass such that
\[
  -\log \Prb[Y=y]\le \bfH(Y).
\]
More generally, for every $b>1$, the set of atoms satisfying
\[
  -\log\Prb[Y=y]\le b\bfH(Y)
\]
has total mass at least $1-1/b$.
\end{lemma}

\begin{proof}
The first claim follows by averaging the information random variable
$-\log\Prb[Y]$.  The second follows from Markov's inequality:
\[
  \Prb[-\log\Prb[Y]>b\bfH(Y)]\le 1/b.
\]
\end{proof}

\begin{lemma}[Collision controls split coordinates]
\label{lem:collision-split}
Let $\nu$ be a probability law on a code $\C$, and let $X,Y,Z$ be independent
samples from $\nu$.  For a coordinate $i$ write
\[
  \Coll_i(\nu)=\Prb[X_i=Y_i].
\]
Then
\[
  \Prb[i\in \Split(X,Y,Z)]\le 3\Coll_i(\nu),
\]
where $i\in\Split(X,Y,Z)$ means that $\{X_i,Y_i,Z_i\}$ has size exactly $2$.
Consequently, for every block $I$,
\[
  \EE|\Split_I(X,Y,Z)|\le 3\sum_{i\in I}\Coll_i(\nu).
\]
If the average collision on $I$ is at most $\alpha$, then
\[
  \Prb[|\Split_I(X,Y,Z)|\le 6\alpha |I|]\ge 1/2.
\]
\end{lemma}

\begin{proof}
Write $p_a=\Prb[X_i=a]$.  The probability of a $2+1$ split at coordinate $i$
is
\[
  3\sum_a p_a^2(1-p_a)\le 3\sum_a p_a^2=3\Coll_i(\nu).
\]
Summing over $i\in I$ gives the expectation bound, and Markov's inequality
gives the final statement.
\end{proof}

\begin{lemma}[Exact sparse split exposure]
\label{lem:exact-split}
Let $W=\Split_I(X,Y,Z)$ and suppose we condition on the event $|W|\le s$.
Then some set $D\subseteq I$, $|D|\le s$, has conditional mass at least
\[
  \left(\sum_{r\le s}\binom{|I|}{r}\right)^{-1}.
\]
On the atom $W=D$, the triple is compatible on $I\setminus D$.
\end{lemma}

\begin{proof}
Under $|W|\le s$, the random set $W$ has at most
$\sum_{r\le s}\binom{|I|}{r}$ possible values.  Choose a most likely atom.
The compatibility assertion is the definition of $W$.
\end{proof}

\begin{lemma}[Role refinement]
\label{lem:role-refinement}
Suppose a triple is split on every coordinate of a block $O$.  Then after
conditioning on a role vector of mass at least $3^{-|O|}$, there is a subblock
$O'\subseteq O$ with $|O'|\ge |O|/3$ and one ordered pair, say $(X,Y)$ after
relabelling, such that
\[
  X_{O'}=Y_{O'},\qquad Z_i\ne X_i\quad (i\in O').
\]
The role-exposure cost is $O(|O'|)$.
\end{lemma}

\begin{proof}
There are at most $3^{|O|}$ role vectors.  On a fixed role vector, one of the
three roles occurs on at least one third of the coordinates.  This majority
class is $O'$.
\end{proof}

\begin{lemma}[KL chain rule and public coordinate localization]
\label{lem:kl-chain}
Let $A=(A_1,\ldots,A_m)$ have laws $P\ll Q$.  For a prefix $u=A_{<t}$ define
\[
  d_t(u)=D(P(A_t\in\cdot\mid A_{<t}=u)\|
          Q(A_t\in\cdot\mid A_{<t}=u)).
\]
Then
\[
  D(P\|Q)=\sum_{t=1}^m
    \EE_{u\sim P(A_{<t})} d_t(u).
\]
If a coordinate is chosen by the first $t$ satisfying a deterministic
threshold in a fixed public order, then the coordinate identity is public.
The only random selector is the prefix value $u$.
\end{lemma}

\begin{proof}
The KL identity is the standard chain rule for relative entropy
\cite{KullbackLeibler51,CoverThomas06,CsiszarKorner11}, here written from the
product rule for conditional probabilities:
\[
  \log\frac{P(A)}{Q(A)}
  =
  \sum_t
  \log
  \frac{P(A_t\mid A_{<t})}{Q(A_t\mid A_{<t})}.
\]
Taking expectation under $P$ gives the identity.  The publicness statement is
by construction.
\end{proof}

\begin{lemma}[Value-pin or hash dichotomy]
\label{lem:fibre-hash}
Let $p$ be a law on a finite alphabet and $S$ a set with $p(S)\ge\rho>0$.
Let $p_S$ be the conditional law on $S$.  For $0<\tau\le 1$, either some
$a\in S$ satisfies
\[
  p(a)\ge \rho\tau,
\]
or
\[
  \Coll(p_S)=\sum_a p_S(a)^2\le \tau.
\]
\end{lemma}

\begin{proof}
If $\max_{a\in S}p_S(a)\ge\tau$, then the first alternative holds.  Otherwise
\[
  \Coll(p_S)\le \max_{a\in S}p_S(a)\sum_{a\in S}p_S(a)<\tau.
\]
\end{proof}

\begin{lemma}[Heavy atoms or low-collision complement]
\label{lem:heavy-or-light-hash}
Let $p$ be a law on a finite alphabet and fix $0<\rho\le1/2$ and
$0<\tau<1$.  Put
\[
  H=\{a:p(a)\ge \rho\tau\},\qquad L=\Omega\setminus H .
\]
Then $|H|\le(\rho\tau)^{-1}$, every branch $V=a\in H$ is a terminal value pin
of conditional mass at least $\rho\tau$, and exactly one of the following
holds for the light complement:
\begin{enumerate}[label=\textup{(\roman*)}]
\item $p(L)<\rho$, in which case pruning $L$ costs at most $2\rho$;
\item $p(L)\ge\rho$, in which case the conditional law $p_L$ has
  $\Coll(p_L)\le\tau$.
\end{enumerate}
\end{lemma}

\begin{proof}
The cardinality bound on $H$ is immediate.  If $p(L)<\rho$, the pruning cost is
$-\log(1-p(L))\le2\rho$.  If $p(L)\ge\rho$, then every $a\in L$ has
$p(a)<\rho\tau$, so
\[
  p_L(a)=\frac{p(a)}{p(L)}<\tau .
\]
Therefore
\[
  \Coll(p_L)\le \max_a p_L(a)\sum_a p_L(a)<\tau .
\]
\end{proof}

\begin{lemma}[Positive KL tail]
\label{lem:positive-tail}
Let $p\ll q$ and $D(p\|q)>0$.  Then
\[
  S_+=\{a:p(a)>q(a)\}
\]
has positive $p$-mass.  If $p(S_+)\ge\rho$, Lemma~\ref{lem:fibre-hash}
applies on $S_+$.  If $p(S_+)<\rho$, pruning $S_+$ costs at most
$-\log(1-p(S_+))\le 2\rho$ for $\rho\le 1/2$, after which
Lemma~\ref{lem:fibre-hash} applies on the complement.
\end{lemma}

\begin{proof}
If $p(S_+)=0$, then $\log(p/q)\le 0$ $p$-almost surely, contradicting
positive KL.  The pruning cost follows from $-\log(1-x)\le 2x$ for
$0\le x\le 1/2$.
\end{proof}

In the transition system, this positive-tail pruning is not an independent
operation which may be repeated for free.  The retained complement
$V_t\notin S_+$ is appended to the same public prefix/tail transcript as the
subsequent value pin or hash coordinate, and its charge
$-\log(1-p(S_+))$ is included in that local terminal/hash charge.  The discarded
branch $V_t\in S_+$ is a charged pruning branch and is not inherited by any
cylinder child.  If this retained complement would exceed the remaining
fixed-window budget, the pruning step is not taken on the continuing branch:
the branch terminalizes on the already paid coordinate-value prefix before any
set-valued complement event is appended.

\begin{lemma}[One-coordinate collision exits are exhaustive tuple exits]
\label{lem:collision-exit-pin}
Let $p,q$ be one-coordinate laws and let $0<\rho\le1/2$.  If
\[
  \sum_a p(a)q(a)>\gamma_0,
\]
then either $\max_a p(a)>\gamma_0^2$ or $\max_a q(a)>\gamma_0^2$.  If
\[
  \Coll(p)=\sum_a p(a)^2>\alpha,
\]
then $\max_a p(a)>\alpha$.  Consequently, whenever
$\rho\tau\le\min(\alpha,\gamma_0^2)$, the cross-collision and
high-side-collision exits in Theorem~\ref{thm:close-pair-exit} route
exhaustively through one actual sample-coordinate law $r$: heavy atom branches
of $r$ are terminal projected tuple pins of conditional mass at least
$\rho\tau$, a complement of mass $<\rho$ is a charged pruning branch, and a
complement of mass at least $\rho$ gives a public low-collision hash coordinate
under the conditional law on that complement.
\end{lemma}

\begin{proof}
By Cauchy's inequality,
\[
  \sum_a p(a)q(a)\le \sqrt{\Coll(p)\Coll(q)}.
\]
If both maxima were at most $\gamma_0^2$, then both collisions would be at
most $\gamma_0^2$, contradicting the strict cross-collision inequality.  The
high-side statement follows from $\Coll(p)\le\max_a p(a)$.  In the
cross-collision case choose $r=p$ or $r=q$ according to the fixed public
tie-break among the laws with a large atom; in the high-side case choose
$r=p$.  Since
$\max_a r(a)>\min(\alpha,\gamma_0^2)\ge\rho\tau$,
Lemma~\ref{lem:heavy-or-light-hash} applies to $r$.  Its heavy branches are
projected tuple pins by Lemma~\ref{lem:tuple-pin-shadow}; its small complement
is charged pruning; and its large complement has low collision for the actual
conditional sample-coordinate law after the public complement event is paid in
the stopped hash transcript.
\end{proof}

\section{Closed Auxiliary Results}
\label{sec:closed-auxiliary}

This section records closed tools which are useful independently of the final
global normal-form assembly.

\begin{lemma}[Pair codegree in a 3-uniform sunflower-free family]
\label{lem:pair-codegree}
Let $\F$ be a 3-uniform family with no three-petal sunflower.  Then every
pair of vertices is contained in at most two members of $\F$.
\end{lemma}

\begin{proof}
If three distinct triples contain the same pair $\{u,v\}$, then their
pairwise intersections are all exactly $\{u,v\}$, so they form a sunflower
with kernel $\{u,v\}$.
\end{proof}

\begin{lemma}[Vertex links have at most six edges]
\label{lem:degree-six}
Let $\F$ be a 3-uniform family with no three-petal sunflower.  For every
vertex $v$, the link graph $L_v$ has maximum degree at most $2$, matching
number at most $2$, and at most $6$ edges.  Hence $\deg_{\F}(v)\le 6$.
\end{lemma}

\begin{proof}
If a vertex $u$ has degree at least $3$ in $L_v$, then three triples contain
the pair $\{u,v\}$, contradicting Lemma~\ref{lem:pair-codegree}.  Thus
$\Delta(L_v)\le 2$.  If $L_v$ contains a matching of size $3$, then the
corresponding three triples in $\F$ all meet pairwise exactly in $v$, forming
a sunflower with kernel $\{v\}$.  Hence $\nu(L_v)\le 2$.

A graph of maximum degree at most $2$ is a disjoint union of paths and cycles.
If it has more than $6$ edges, then its components contain a matching of size
at least $3$: each path or cycle with $e$ edges has matching number at least
$\lfloor e/2\rfloor$, and distributing at most two matching edges among all
components permits at most six edges, with equality only for two triangles.
Therefore $|E(L_v)|\le 6$.
\end{proof}

\begin{theorem}[Pairwise-intersecting 3-uniform local theorem]
\label{thm:pairwise-ten}
Let $\F$ be a pairwise-intersecting 3-uniform family with no three-petal
sunflower.  Then
\[
  |\F|\le 10.
\]
If equality holds, then, up to relabelling, $\F$ is the $2$-$(6,3,2)$ design.
\end{theorem}

\begin{proof}
The pair-codegree bound is Lemma~\ref{lem:pair-codegree}.  We first record a
small graph fact.  If two graphs $X,Y$ have maximum degree at most $2$ and
every edge of $X$ intersects every edge of $Y$, then $|X|+|Y|\le 6$.  Indeed,
any edge of $Y$ is a two-vertex cover of $X$, and conversely.  A graph of
maximum degree at most $2$ with a two-vertex cover has at most $4$ edges; if
one side has $4$ edges, its possible two-covers form at most two edges for
the other side, while if one side has $3$ edges the possible two-covers form
at most three edges.  The cases with at most two edges are immediate.

Let $\tau(\F)$ be the vertex-covering number.  Since any member of $\F$ meets
all other members, $\tau(\F)\le 3$.  If $\tau(\F)\le 2$, fix a two-vertex
cover $\{1,2\}$.  Apart from the at most two triples containing both $1$ and
$2$, the triples through $1$ and the triples through $2$ give two
cross-intersecting link graphs of maximum degree at most $2$.  The graph fact
gives $|\F|\le 2+6=8$.

It remains to treat $\tau(\F)=3$.  Fix a triple $T=\{1,2,3\}\in\F$.  Every
other member meets $T$.  For each pair $ij\subset T$, let $H_{ij}$ be the
possible extra triple containing that pair.  By Lemma~\ref{lem:pair-codegree}
there is at most one such extra triple for each pair.  Let $h$ be the number
of existing $H_{ij}$.

For $i=1,2,3$, let $X_i$ be the graph of outside pairs $e$ for which
\[
  \{i\}\cup e\in\F
\]
and this triple contains no other vertex of $T$.  Each $X_i$ is an
intersecting graph of maximum degree at most $2$: two disjoint edges in
$X_i$, together with the edge $T\setminus\{i\}$ in the link $L_i$, would give
a matching of size $3$ in $L_i$ and hence a sunflower with kernel $\{i\}$.
Therefore $|X_i|\le 3$, with equality only for a triangle.  Also $X_i$ and
$X_j$ are cross-intersecting for $i\ne j$, because $\F$ is pairwise
intersecting.

If $H_{12}=\{1,2,a\}$ exists, then every edge of $X_3$ contains $a$, since
the corresponding triple must meet $H_{12}$ outside $T$.  Hence
$|X_3|\le 2$.  Moreover $\deg_{X_1}(a)\le 1$ and $\deg_{X_2}(a)\le 1$,
because the pair $\{1,a\}$ or $\{2,a\}$ already occurs in $H_{12}$ and cannot
occur in two further triples.

The capacity bounds now give
\[
\begin{array}{c|c}
h& |\F|\\ \hline
0& \le 1+8\\
1& \le 1+1+(3+3+2)\\
2& \le 1+2+(3+2+2)\\
3& \le 1+3+(2+2+2).
\end{array}
\]
For $h=0$, the value $1+9$ would require three pairwise cross-intersecting
triangles.  Cross-intersecting triangles must have the same vertex set, which
would make every outside pair occur in three triples, violating
Lemma~\ref{lem:pair-codegree}.  Thus the $h=0$ case gives at most $9$.
All cases give $|\F|\le 10$.

We finally identify equality.  If $h=1$, say $H_{12}=\{1,2,a\}$, equality
would force $X_1$ and $X_2$ to be cross-intersecting triangles, hence the
same triangle.  The two edges of $X_3$ must hit that triangle and also contain
$a$, forcing an outside pair to appear too often.  Thus $h=1$ cannot give
equality.

If $h=2$, say $H_{12}=\{1,2,a\}$ and $H_{13}=\{1,3,b\}$, equality would force
$X_1$ to be a triangle and $X_2,X_3$ to be two-edge stars.  The forced-star
condition makes $X_2$ centered at $b$, and the cross-intersection with the
triangle $X_1$ forces $\deg_{X_1}(b)=2$, contradicting the degree cap from
$H_{13}$.

Thus equality has $h=3$.  Write
\[
  H_{12}=\{1,2,a\},\quad
  H_{13}=\{1,3,b\},\quad
  H_{23}=\{2,3,c\}.
\]
Then $X_1,X_2,X_3$ are two-edge stars centered respectively at $c,b,a$.
The degree caps force $a,b,c$ to be distinct.  Two cross-intersecting
two-edge stars with different centers must share the center-edge and have
their remaining edges point to the same third vertex.  Applying this to the
three pairs of stars gives, after relabelling,
\[
  X_1=\{ac,bc\},\qquad
  X_2=\{ab,bc\},\qquad
  X_3=\{ab,ac\}.
\]
Hence the ten triples are
\[
  123,\ 12a,\ 13b,\ 23c,\ 1ac,\ 1bc,\ 2ab,\ 2bc,\ 3ab,\ 3ac.
\]
Every pair of the six vertices $\{1,2,3,a,b,c\}$ occurs in exactly two
triples, so this is the $2$-$(6,3,2)$ design.  The preceding forced choices
also prove uniqueness.
\end{proof}

\begin{corollary}[Size-nine stability in the pairwise case]
\label{cor:size-nine}
Let $\F$ be a pairwise-intersecting 3-uniform family with no three-petal
sunflower and $|\F|=9$.  Then $\F$ is obtained from the $2$-$(6,3,2)$ design
by deleting one triple.
\end{corollary}

\begin{proof}
Use the notation from the proof of Theorem~\ref{thm:pairwise-ten}.  The case
$\tau(\F)\le 2$ gives at most $8$ triples, so $\tau(\F)=3$.  The case $h=0$
would require sizes $(3,3,2)$ among $X_1,X_2,X_3$; the two triangles must
coincide, and the third graph must use edges of the same triangle, producing
pair-codegree $3$.  Thus $h=0$ is impossible.  If $h=1$, then either two of
the $X_i$ are coincident triangles, giving the same pair-codegree violation,
or one triangle is crossed by a forced two-edge star, forcing degree $2$ at
the special vertex and contradicting the degree cap from the corresponding
$H_{ij}$.

For $h=2$, say $H_{12}$ and $H_{13}$ exist, the capacity bound forces all
three graphs $X_1,X_2,X_3$ to have two edges.  The forced centers and the
cross-intersecting two-star calculation used in
Theorem~\ref{thm:pairwise-ten} determine the unique configuration, namely
the $2$-$(6,3,2)$ design with the missing triple $H_{23}$ deleted.  For
$h=3$, the sizes of $X_1,X_2,X_3$ are $(2,2,1)$ up to permutation; the same
two-star calculation determines the unique configuration, now the
$2$-$(6,3,2)$ design with one $X_i$-type triple deleted.  These are exactly
the one-triple deletions of the design.
\end{proof}

\begin{proposition}[Linear rank case]
\label{prop:linear-case}
Let $\F$ be a rank at most $k$ linear family with no three-petal sunflower.
Then
\[
  |\F|\le 2k+2.
\]
\end{proposition}

\begin{proof}
Linearity implies that three sets through one vertex would form a sunflower
with that vertex as kernel, so every vertex degree is at most $2$.  The family
also has matching number at most $2$, because three disjoint sets form a
sunflower with empty kernel.

Let $M=|\F|$ and form the disjointness graph $G$ on vertex set $\F$.  Since
the matching number of $\F$ is at most $2$, the graph $G$ is triangle-free.
For any $A\in\F$, at most $k$ other members of $\F$ intersect $A$: each vertex
of $A$ is contained in at most one further member.  Hence
\[
  \deg_G(A)\ge M-1-k.
\]
Thus
\[
  |E(G)|\ge \frac{M(M-1-k)}2.
\]
By Mantel's theorem \cite{Mantel07}, $|E(G)|\le M^2/4$.  For $M>0$ this gives
$M-1-k\le M/2$, hence $M\le 2k+2$.
\end{proof}

\begin{theorem}[Projection-profile slice-rank bound]
\label{thm:projection-profile}
Let $\C\subseteq\prod_{i=1}^k\Omega_i$ be a balanced two-collision code.
For a map
\[
  \sigma:[k]\to\{0,xy,xz,yz\}
\]
define $R_x(\sigma)$ to be the coordinates where $\sigma(i)$ involves $x$,
and define $R_y(\sigma),R_z(\sigma)$ similarly.  Then
\[
  |\C|
  \le
  \sum_{\sigma}
  \min_{u\in\{x,y,z\}}|\pi_{R_u(\sigma)}(\C)|.
\]
\end{theorem}

\begin{proof}
This is a projection-profile variant of the slice-rank polynomial method
\cite{CLP17,EG17,Tao16}.
Work over $\mathbb F_3$.  On one coordinate define
\[
  h(a,b,c)=1-\mathbf 1_{a=b}-\mathbf 1_{a=c}-\mathbf 1_{b=c}.
\]
In $\mathbb F_3$ this equals $0$ exactly when exactly two of $a,b,c$ are
equal, and equals $1$ when all three are equal or all three are distinct.
Set
\[
  H(x,y,z)=\prod_{i=1}^k h(x_i,y_i,z_i).
\]
Restricted to $\C^3$, the balanced two-collision property gives
\[
  H(x,y,z)=
  \begin{cases}
  1,&x=y=z,\\
  0,&\text{otherwise}.
  \end{cases}
\]
Thus the slice rank of $H|_{\C^3}$ is $|\C|$.

Expanding the product, each term is indexed by
$\sigma:[k]\to\{0,xy,xz,yz\}$.  For such a term, the $x$-dependence is only
through $\pi_{R_x(\sigma)}(x)$, so it has slice rank at most
$|\pi_{R_x(\sigma)}(\C)|$ when sliced in the $x$ variable.  The same argument
applies to $y$ and $z$.  Hence the term has slice rank at most the minimum of
the three displayed projection sizes.  Slice rank is subadditive, giving the
claim.
\end{proof}

\begin{corollary}[Bounded-alphabet and variable-profile bounds]
\label{cor:bounded-alphabet}
If $|\Omega_i|\le Q$ for all $i$, then every balanced two-collision code
satisfies
\[
  |\C|\le (1+3Q^{2/3})^k\le (4Q^{2/3})^k.
\]
More generally, if every projection satisfies
$|\pi_R(\C)|\le D^{|R|}$, then
\[
  |\C|\le (1+3D^{2/3})^k.
\]
\end{corollary}

\begin{proof}
From Theorem~\ref{thm:projection-profile},
\[
  |\C|
  \le
  \sum_{\sigma}
  \min_{u}D^{|R_u(\sigma)|}
  \le
  \sum_{\sigma}
  D^{(|R_x(\sigma)|+|R_y(\sigma)|+|R_z(\sigma)|)/3}.
\]
For each coordinate, the choice $0$ contributes $0$ to the sum
$|R_x|+|R_y|+|R_z|$, while each of $xy,xz,yz$ contributes $2$.  Therefore the
last sum equals $(1+3D^{2/3})^k$.  The bounded-alphabet statement follows by
taking $D=Q$.
\end{proof}

\section{Tuple Nodes and Conditioned Copying}

Local close-pair and hash statements are often true only under a tuple-level
conditioning event.  They must not be converted directly into code-level
statements.

\begin{definition}[Tuple node]
A tuple node is an event
\[
  E\subseteq \C^r
\]
for independent samples $X^1,\ldots,X^r$.  All local laws at the node are
conditional laws under $\Prb(\cdot\mid E)$.
\end{definition}

\begin{lemma}[Tuple-level pins are projected shadows]
\label{lem:tuple-pin-shadow}
At a tuple node $E\subseteq\C^r$, suppose a branch terminalizes by fixing
$X^s_J=a$ for one of the active sample labels $s$.  Then the terminal event is
compatible with the projected shadow obtained by appending $(J,a)$ to the
$s$-sample list.  The mass of the tuple node and of the local atom is counted
by the proof-tree entropy ledger; no ambient-code density increment is being
asserted.
\end{lemma}

\begin{proof}
The terminal event is a subset of
\[
  E\cap\{X^s_J=a\}
  \subseteq \{X^s_J=a\}.
\]
This is exactly the compatibility condition in the definition of a
multi-sample value-shadow.  The probability of reaching $E$ and then choosing
the local atom is part of the full leaf transcript $\Gamma$ and is therefore
accounted for by Proposition~\ref{prop:tree-entropy-ledger}.  After projecting
away the analysis transcript, only the cylinder $\{X^s_J=a\}$ remains in the
projected shadow.
\end{proof}

\begin{lemma}[Tuple prefixes project to multi-sample shadows]
\label{lem:tuple-prefix-shadow}
At a tuple node $E\subseteq\C^r$, suppose a branch terminalizes by fixing a
finite public list of tuple-coordinate values
\[
  X^{s_\ell}_{j_\ell}=a_\ell\qquad(\ell=1,\ldots,L),
\]
with sample labels and coordinates determined by the paid public prefix.  First
discard from this list every exact repeat whose value is already forced by an
earlier retained tuple-coordinate value; such a repeat has zero conditional
entropy in the prefix chain.  If an exact repeat asks for a different value, the
branch has zero mass and is ignored.  For every remaining paid variable choose
one sample-label coordinate pin, grouping the chosen pins by sample label.  Then
the number of entropy-bearing prefix variables is at most $T(\Theta)$ for the
resulting projected shadow.  The terminal event is compatible with this
multi-sample projected value-shadow, and its local probability cost is the
entropy of the paid prefix list in the proof-tree ledger, not an ambient-code
density increment.
\end{lemma}

\begin{proof}
After the zero-entropy repeats are discarded, the terminal event is contained
in the intersection of the retained coordinate cylinders.  For a
tuple-coordinate variable which is already a common value of two labels,
choosing either label gives a cylinder containing the terminal event.  Grouping
the chosen pins by sample label rewrites this intersection in the form required
by Definition~\ref{def:multi-sample-shadow}.  The full prefix values, including
the deterministic repeats before projection, remain in the terminal leaf
transcript before projection and are paid by
Proposition~\ref{prop:tree-entropy-ledger}.
\end{proof}

\begin{lemma}[Fixed-node conditioned copying]
\label{lem:fixed-copy}
Let $E\subseteq\C^r$ have mass $\delta>0$, and let
\[
  \nu=\Law(X^s\mid E).
\]
If $F\subseteq\C^m$ satisfies $\nu^m(F)\ge\rho$, then in $m$ independent
copies of the whole tuple experiment the event that all copies lie in $E$ and
their selected $s$-samples satisfy $F$ has mass at least $\delta^m\rho$.
\end{lemma}

\begin{proof}
Conditioning the $m$ independent copies on $E^m$ makes the selected samples
independent with common law $\nu$.  The mass is exactly
$\delta^m\nu^m(F)$.
\end{proof}

\begin{lemma}[Random-transcript conditioned copying]
\label{lem:random-copy}
Let $K$ be a finite public key whose support has already been restricted to the
surviving key atoms, with atoms $k$ and probabilities $p_k$.  On the atom
$K=k$, let $\nu_k$ be the selected one-sample conditional law.  Suppose that for
every atom in this support there is an event $F_k$ for $m$ independent draws
from $\nu_k$ satisfying
\[
  \nu_k^m(F_k)\ge \rho .
\]
Then in $m$ independent copies of the whole stopped experiment, the event that
all copies have the same key $k$ and their selected samples satisfy $F_k$ has
mass at least
\[
  \rho\sum_k p_k^m
  =
  \rho\exp(-(m-1)\bfH_m(K)),
\]
where
\[
  \bfH_m(K)=\frac{1}{1-m}\log\sum_k p_k^m
\]
is the order-$m$ Renyi entropy.
\end{lemma}

\begin{proof}
The desired mass is
\[
  \sum_k p_k^m\nu_k^m(F_k).
\]
Since the support of $K$ consists only of surviving atoms, the uniform lower
bound $\nu_k^m(F_k)\ge\rho$ applies to every summand and gives the displayed
estimate.  In applications one first conditions the transcript on the good-key
event and pays the fixed mass loss separately.
\end{proof}

Thus for random transcripts the copy cost is the Renyi entropy of the public
key that is shared by the copies.  If $K$ is the public key and $m=3$, the copy
tax is
\[
  2\bfH_3(K)+O(1).
\]
Thus hash returns must use public keys whose entropy is already paid.

\section{Public Hash Windows}

The following theorem is the corrected form of the fixed hash return.

\begin{definition}[Labelled hash-window survival convention]
\label{def:hash-survival}
Fix $0<\tau\le\beta<1/12$.  A hash window also fixes a public active sample
label $s$.  Every produced hash coordinate is an ambient coordinate of $X^s$
under the current paid tuple event; coordinates belonging to different sample
labels are placed in different hash windows.  A produced hash coordinate is
live after a transcript prefix if its current one-coordinate conditional law
has no atom of mass $\ge\beta$.  The hash-window rule scans the public list of
produced coordinates after every transcript extension.  If a produced
coordinate is not live, the proof chooses the first public atom of mass at
least $\beta$, restricts to that atom, and stops as a terminal value-pin
shadow; the discarded complement is a charged pruning loss at most
\(\log(1/\beta)\) and is not inherited by any child.  A coordinate which is
live at the terminal transcript is called surviving.
Since collision is bounded by the largest atom, every surviving coordinate
has terminal collision at most $\beta$.  Since a codeword atom is bounded by
any of its coordinate atoms on a nonempty surviving block, the conditional
one-sample law on a surviving hash block has maximum codeword atom $<\beta$.
In a fixed-window statement,
$\mathrm{Destroyed}$ denotes the number of produced coordinates which are not
surviving, after terminal value-pin and bounded-support exits have been
removed.
\end{definition}

\begin{theorem}[Public-transcript fixed-window hash return]
\label{thm:public-hash}
Fix a public state $P_0$.  Run a bounded hash window with terminal transcript
$T$ satisfying $\bfH(T\mid P_0)\le B_{\rm win}$.  Suppose coordinate
production is deterministic from $P_0$ and the transcript prefix, and suppose
all produced coordinates are attached to the same public sample label $s$.
Suppose every non-bounded branch produces at least $2h$ distinct hash
coordinates, and the adaptive destruction budget satisfies
\[
  \EE[\mathrm{Destroyed}\mid P_0]\le h/2.
\]
Then, except for terminal value-pin or bounded-support exits, there is a
transcript event of conditional mass at least $1/2$ on which a block $H(T)$ of
$h$ surviving hash coordinates is deterministic from the paid transcript.  For
each transcript in that event, three conditioned copies of the conditional law
of the same sample label $X^s$ are distinct and give a near-sunflower event on
$H(T)$ with conditional probability at least
$\rho_{\rm hash}:=1/2-3\beta>0$.  Equivalently, the returned
near-sunflower has a fixed positive mass, and its total proof-tree cost is
\[
  O(B_{\rm win}+h).
\]
Here the cost is the local entropy/mass charge used in
Proposition~\ref{prop:tree-entropy-ledger}: the public transcript entropy,
the order-three Renyi copy tax for the transcript key, and the exact sparse
split exposure on the returned block.
\end{theorem}

\begin{proof}
If fewer than $h$ coordinates survive, more than $h$ are destroyed.  Markov's
inequality gives probability at most $1/2$ for this bad event.  On the good
event choose the first $h$ surviving produced coordinates.  This block is a
deterministic function of $(P_0,T)$, so it carries no entropy beyond the paid
transcript.

Let $G$ be the good-transcript event and let $T_G$ be $T$ conditioned on $G$.
Since $\Prb(G\mid P_0)\ge1/2$ and $\bfH(T\mid P_0)\le B_{\rm win}$,
\[
  \bfH_3(T_G\mid P_0)\le \bfH(T_G\mid P_0)\le 2B_{\rm win}+O(1).
\]
For each good transcript $t$, let
\[
  \nu_t=\Law(X^s\mid P_0,T=t).
\]
Every coordinate of $H(t)$ has collision at most $\beta$ under the
corresponding one-coordinate marginal of $\nu_t$.  Hence
Lemma~\ref{lem:collision-split} gives the near-sunflower split event for
three independent samples from $\nu_t$ with probability at least $1/2$:
\[
  \nu_t^3\{|\Split_{H(t)}(X,Y,Z)|\le 6\beta h\}\ge 1/2 .
\]
Because $H(t)$ is nonempty and all its coordinates are surviving,
\[
  \max_x\nu_t(x)<\beta .
\]
Thus the probability that three independent $\nu_t$-samples are not all
distinct is at most
\[
  3\sum_x\nu_t(x)^2\le 3\beta .
\]
With $\beta<1/12$, the distinct near-sunflower event
\[
  F_t=\{X,Y,Z\text{ distinct and }|\Split_{H(t)}(X,Y,Z)|\le 6\beta h\}
\]
has $\nu_t^3(F_t)\ge\rho_{\rm hash}:=1/2-3\beta>0$ for every good transcript
$t$.

Lemma~\ref{lem:random-copy} realizes these three conditional samples in the
original tuple experiment using the common good transcript key
$K=T_G$.  Since $\Prb(G\mid P_0)\ge1/2$, passing from the conditioned key
$T_G$ back to the original stopped experiment loses only the fixed factor
$\Prb(G\mid P_0)^3\ge 2^{-3}$, absorbed in the $O(1)$ term.  The copy charge is
therefore
\[
  2\bfH_3(T_G\mid P_0)+O(1)=O(B_{\rm win}+1).
\]
Finally the exact split set on $H(T)$ has at most $2^h$ possible values, and
under the sparse threshold used later at most
$\sum_{r\le 6\beta h}\binom hr$ values; in either case the fixed-window
exposure costs $O(h)$.  Thus the total charged cost is $O(B_{\rm win}+h)$.
\end{proof}

\begin{remark}
This theorem is only a fixed-window theorem.  It is designed for a
constant-base result, not for optimizing constants.
In the normal-form transition system it is applied in the stronger regime
$\mathrm{Destroyed}=0$, because the survival scan terminalizes as soon as a
produced coordinate becomes non-live.  The displayed destruction-budget
hypothesis is included to isolate the robust fixed-window statement from its
particular use in Proposition~\ref{prop:hash-ledger}.
\end{remark}

\begin{remark}[Worst-case transcript budget]
The bound $\bfH_3(T_G\mid P_0)\le \bfH(T_G\mid P_0)$ is deliberately
worst-case.  No saving from nonuniform transcripts is used anywhere in the
global charge.  The window rule is stronger than an a posteriori entropy
bound: a continuing branch appends prefix, tail/pruning, and production data
only while the cumulative public description remains within the fixed budget
$B_{\rm win}$.  A rare prefix whose next description would exceed the remaining
budget terminalizes as a projected value-shadow before it enters the hash
window; a set-valued tail/pruning event is never retained as the terminal
shadow.
Thus every nonterminal window key has already been paid by a fixed stopped
transcript budget, and the Renyi copy tax is bounded by the same constant even
in the worst transcript distribution.
\end{remark}

\begin{proposition}[Fixed-window hash ledger]
\label{prop:hash-ledger}
Choose fixed hash-window parameters
\[
  h,\ \beta,\ \tau,\ B_{\rm win}
\]
with $h\ge1$ and $0<\tau\le\beta<1/12$, and enforce
Definition~\ref{def:hash-survival}.
Here $B_{\rm win}$ is a budget for one labelled window.  A branch may contain
several labelled windows over its lifetime, but each such window has one public
sample label, one stopped transcript of size at most $B_{\rm win}$, and one
bounded charge later terminalized or assigned by
Proposition~\ref{prop:hash-token-matching}.  Then every nonterminal hash window generated by
Definition~\ref{def:transition-system} either satisfies the hypotheses of
Theorem~\ref{thm:public-hash}, or exits to one of:
\begin{enumerate}[label=\textup{(\roman*)}]
\item a terminal projected value pin;
\item a bounded hash-support block $B$ with $1\le |B|<2h$ whose fresh part
$B\setminus U$ is promoted immediately to $U$;
\item a bounded hash-support block $B\subseteq U$ whose incidences are
registered immediately in the old-incidence ledger, unless the registration
triggers the high-revisit exit.
\end{enumerate}
Moreover every nonterminal hash-window transcript charge is assigned to a
terminal, fresh, old-incidence, scale-descent, or close-pair birth certificate
before any child state or descendant hash window is entered from the returned
block.
\end{proposition}

\begin{proof}
Hash coordinates are produced by Theorem~\ref{thm:positive-kl-public} from a
public coordinate order after the prefix value has been paid.  Each produced
coordinate is assigned to its public sample label, and hash windows are
maintained separately for different labels.  Within one window the cumulative
coordinate-production transcript is stopped at the fixed budget $B_{\rm win}$;
any branch whose next prefix would exceed the remaining budget terminalizes on
that projected prefix instead of adding the coordinate to the window.  Hence
within one window the coordinate production rule is deterministic from the
public state, the public sample label, and the paid transcript prefix.  If a
value pin occurs at production, or later by the scan in
Definition~\ref{def:hash-survival}, the branch terminalizes with the fixed
pruning loss specified there.  Thus on a
nonterminal branch every produced coordinate that remains in the window is
surviving and has terminal collision at most $\beta$ for the conditional law
of that same sample label.

If the cumulative budget closes before $2h$ distinct surviving coordinates
are produced and none has yet been retained, the branch has already stopped on
the projected prefix described above.  If no coordinate has been accepted and
the ambient close-pair/residual decision stops supplying a coordinate-producing
exit, then no hash window has been opened and no hash token is created; the
branch follows that non-hash decision.  Otherwise let $B$ be the nonempty set of
surviving coordinates already produced.  Then $|B|<2h$.  If $B\setminus U$ is
nonempty, that public fresh part is promoted immediately to $U$ and the bounded
hash transcript is charged to the fresh drop.  If $B\subseteq U$, the
incidences of $B$ are registered in the old ledger; a counter crossing the
revisit cap triggers the high-revisit exit, and otherwise the old potential
pays the bounded transcript.  Thus the bounded-support exit is absorbed
without opening another hash window and without appealing back to this
proposition.  If at least $2h$ surviving coordinates are produced before the
budget
closes, the transcript entropy is at most $B_{\rm win}$ by construction,
coordinate production is public, and the adaptive destruction budget in
Theorem~\ref{thm:public-hash} is zero.  Hence the hypotheses of that theorem
hold.

It remains to certify the charge of nonterminal windows.  If a bounded support
contains a fresh coordinate, promote the whole fresh part to $U$ and charge
the fixed window cost to the fresh/old potential drop.  If the bounded support
is contained in $U$, register its coordinates as old incidences; repetition
above the revisit threshold triggers the high-revisit exit, while below that
threshold the incidences are debited in $N_{\rm old}$.  If a surviving hash block
returns to residual repair, its window charge is carried as a pending charge
of that returned block.  The residual node must leave residual mode before a
new hash window can be opened from it.  Its first nonterminal exit is fresh
support promotion, an old incidence/window, a deletion-trapped scale descent,
or the birth of a public close-pair block; a bounded leaf or terminal shadow
pays the charge directly.  The pending charge is a bookkeeping token of size
at most $C_{\rm hash}$: assigning it to the first exit means that this amount
is added to that exit's local charge before the branchwise entropy sum is
estimated.  In the deletion-trapped case the assignment is made at the parent
residual transition, before the smaller-scale child is entered; the token is
added to the scale ledger for that descent or to the birth certificate of the
new active block.  The child state therefore carries no hash token from its
parent.  The token is then removed from the returned block's state, not from
the entropy ledger already paid.  Thus a hard exit receives only the bounded
charge of the hash window which created the block currently being processed,
so hash transcripts do not form a separate reusable budget.
\end{proof}

\begin{proposition}[Hash-coordinate accumulation dichotomy]
\label{prop:hash-accumulation}
During one labelled hash window, the proof does not assume that the
product-KL or collision exit remains available a prescribed number of times.
The window is run with the following exit dichotomy.  Each time a
one-coordinate exit produces a public low-collision coordinate for the current
sample label, its prefix, tail/pruning, and production transcript are appended
only if the cumulative paid transcript remains within $B_{\rm win}$.  If this
process reaches $2h$ distinct live coordinates, Theorem~\ref{thm:public-hash}
is applied.  If the process stops first, then either the branch has already
terminalized on a projected value/prefix pin, or the nonempty set of live
coordinates accumulated so far is a bounded hash-support block and exits
through the fresh/old alternatives of Proposition~\ref{prop:hash-ledger}.

Consequently the global proof never needs a lower bound on the frequency of
positive-KL triggers inside a window.  A window either returns a fixed-size
hash near-sunflower, or its already produced bounded support is charged before
another hash window can be opened.  In particular, no estimate is made or
needed for the typical depth of the KL prefix search.  Long prefixes are
terminal projected value-shadows, where their entropy is paid by the pinned
prefix coordinates through $T(\Theta)$; only prefixes whose cumulative
description fits inside $B_{\rm win}$ are allowed to contribute to a hash key.
\end{proposition}

\begin{proof}
The coordinate-producing step is exactly the low-collision alternative of
Theorem~\ref{thm:positive-kl-public} or the corresponding high-revisit/collision
exit, run under the current paid public transcript.  If the step instead gives
a value atom or a prefix whose next transcript would exceed the remaining
budget, Proposition~\ref{prop:hash-ledger} terminalizes it as a projected
shadow and no hash coordinate is added.  If a coordinate is added, the
survival scan in Definition~\ref{def:hash-survival} immediately terminalizes
any later atom of mass at least $\beta$, so every nonterminal coordinate in the
window is live.

Thus, before the window has produced $2h$ distinct live coordinates, the
retained hash state contains only a public set $B$ with $|B|<2h$ and a paid
transcript of length at most $B_{\rm win}$.  If production continues until
$2h$ coordinates are present, the hypotheses of Theorem~\ref{thm:public-hash}
hold, with destruction budget zero.  If production cannot continue inside the
budget, or if the subsequent close-pair/residual decision no longer supplies a
coordinate-producing exit for this labelled window, there are two cases.  When
$B=\varnothing$, no labelled window has yet been opened and no hash transcript
or token is retained.  When $B\ne\varnothing$,
Proposition~\ref{prop:hash-ledger} classifies the current bounded set $B$: its
fresh part is promoted, or if $B\subseteq U$ its old incidences are registered
unless high revisit triggers.  That classification is a terminal exit from the
current hash window, so the same bounded transcript cannot be charged again.
\end{proof}

\begin{proposition}[Hash-token matching]
\label{prop:hash-token-matching}
On each retained branch, the nonterminal fixed-window hash returns which create
pending charges admit a matching $\phi$ from pending hash returns to
terminal/fresh/old/scale/birth hard exits.  Each return is assigned to the first
hard exit supplied by its returned block, and every hard exit is matched to at
most one hash return.  Consequently every local hard-exit inequality only has
to pay one additional $C_{\rm hash}$ token.
\end{proposition}

\begin{proof}
A pending token is attached to a specific returned residual block.  By
Proposition~\ref{prop:hash-ledger} and the branch invariant in
Proposition~\ref{prop:branch-invariants}, that block is processed until the
present token is assigned; before this assignment no second hash-token-bearing
window is retained from the block or from a child born inside that block.  If
the block exits by a terminal shadow, fresh promotion, old-incidence
registration, scale descent, or birth of a close-pair block, define $\phi$ to
be that first exit.  A birth transition is therefore not a way to pass the token
to a descendant: the token is assigned at the parent birth step before the
newly born close-pair block is entered.  In the deletion-trapped case the
assignment is made at the parent residual transition before the smaller-scale
child is entered.  If a window accepts no coordinate, no token is created, and
if it stops as a bounded-support exit, that same fresh/old exit is its image.

Suppose two distinct pending returns were assigned to the same hard exit.  The
later return would have had to be retained before the earlier token was removed
from the same active residual block, or else the hard exit would already have
occurred before the later return was created.  Both alternatives contradict the
first paragraph.  Thus $\phi$ is injective on hard exits.
\end{proof}

\begin{lemma}[Token-free children]
\label{lem:token-free-children}
Every nonterminal child state entered by the transition system is token-free,
except for the unique residual state produced immediately by a nonterminal
fixed-window hash return.  That exceptional residual state carries exactly the
token of the window which produced it.  Before any child of that residual state
is entered, and before any new hash window is opened from data produced by that
residual state, the token is assigned to the terminal, fresh, old, scale, or
birth exit which leaves the residual state.
\end{lemma}

\begin{proof}
The only transition which creates a pending hash token is the nonterminal
return of Theorem~\ref{thm:public-hash}; by Proposition~\ref{prop:hash-ledger}
the token is attached to the returned residual block.  All other hash-window
outcomes are terminal value pins or bounded-support exits, and those are
charged immediately as terminal, fresh, or old exits.

Consider the next transition out of this returned residual state.  The
transition system lists only terminal shadow, fresh promotion, old-incidence
registration/window, deletion-trapped scale descent, bounded leaf, or birth of
a public close-pair block.  In the first three nonterminal hard exits the token
is included in the local charge of that exit.  In a deletion-trapped scale
descent it is included in the parent scale step before the smaller-scale child
is entered.  In a birth transition it is included in the parent birth
certificate before the close-pair child is entered.  A bounded leaf stops the
recurrence and a terminal shadow has no child.  Hence every entered child is
token-free, and no descendant hash window can be opened while an ancestor token
is still pending.
\end{proof}

\section{Residual Repair Lemmas}
\label{sec:residual-repair}

\begin{lemma}[DRC row block]
\label{lem:drc-row-block}
Let $(\Omega,\lambda)$ be a probability space, let $I$ be a finite set with
$|I|=n$, and let $(E_i)_{i\in I}$ be events satisfying
\[
  \frac1n\sum_{i\in I}\lambda(E_i)\ge \delta
\]
for some $0<\delta\le 1$.  If
\[
  1\le s\le \delta n/4,
\]
then there is a set $J\subseteq I$, $|J|=s$, such that
\[
  \lambda\Bigl(\bigcap_{j\in J}E_j\Bigr)
  \ge
  \frac{\delta}{2}\left(\frac{\delta}{4}\right)^s.
\]
For fixed $\delta$, conditioning on this DRC atom costs $O_\delta(s)$ in the
proof-tree ledger.
\end{lemma}

\begin{proof}
This is the one-row form of dependent random choice \cite{FoxSudakov11}.
For $\omega\in\Omega$, put
\[
  d(\omega)=|\{i\in I:\omega\in E_i\}|.
\]
Then $\EE[d(\omega)/n]\ge\delta$.  Let
\[
  G=\{\omega:d(\omega)/n\ge\delta/2\}.
\]
Since $d/n\le 1$,
\[
  \delta\le \EE[d/n]\le \lambda(G)+(1-\lambda(G))\delta/2,
\]
so $\lambda(G)\ge\delta/2$.

	Choose $J$ uniformly from the $s$-subsets of $I$.  For $\omega\in G$,
	$d(\omega)\ge\delta n/2$, and $s\le\delta n/4$ gives
	\[
	  \Prb[J\subseteq \{i:\omega\in E_i\}]
	  =
	  \frac{\binom{d(\omega)}s}{\binom ns}
	  \ge
	  \left(\frac{\delta}{4}\right)^s.
	\]
	Indeed each factor in
	$\prod_{r=0}^{s-1}(d(\omega)-r)/(n-r)$ is at least
	$(\delta n/2-s)/n\ge\delta/4$.
	Averaging over $\omega$ yields
\[
  \EE_J \lambda\Bigl(\bigcap_{j\in J}E_j\Bigr)
  \ge
  \lambda(G)\left(\frac{\delta}{4}\right)^s
  \ge
  \frac{\delta}{2}\left(\frac{\delta}{4}\right)^s.
\]
Thus some $J$ satisfies the claimed bound.
The displayed lower bound has negative logarithm
\[
  \log(2/\delta)+s\log(4/\delta)=O_\delta(s),
\]
which is the stated conditioning cost.
\end{proof}

\begin{lemma}[Residual repair pressure]
\label{lem:residual-pressure}
Let $\C$ be a balanced two-collision code.  Let
$\nu$ be a law on distinct triples $(X,Y,Z)\in\C^3$ which are compatible on a
retained block $R$, meaning that for every $r\in R$ the set
$\{X_r,Y_r,Z_r\}$ has size $1$ or $3$.  Put $L=[k]\setminus R$ and
\[
  W_L(X,Y,Z)=\{t\in L:|\{X_t,Y_t,Z_t\}|=2\}.
\]
Then $|W_L(X,Y,Z)|\ge 1$ pointwise, and hence
\[
  \sum_{t\in L}\Prb_\nu[t\in W_L]\ge 1.
\]
\end{lemma}

\begin{proof}
Because $\C$ is a two-collision code, every distinct triple of codewords has
a coordinate with a $2+1$ split.  Compatibility excludes such a split on
$R$, so at least one split coordinate lies in $L$.  Taking expectations gives
the displayed inequality.
\end{proof}

\begin{lemma}[Public residual support promotion]
\label{lem:residual-support-promotion}
In the setup of Lemma~\ref{lem:residual-pressure}, let $U$ be the current old
support and put
\[
  F=\{t\in L\setminus U:\Prb_\nu[t\in W_L]>0\}.
\]
The set $F$ is determined by the current public law.  Fix a constant
$q\ge 1$.  Then one of the following transitions is available:
\begin{enumerate}[label=\textup{(\roman*)}]
\item fresh support promotion: $|F|>q$, in which case the public set $F$ is
promoted to old support and gives a potential drop $(B_f-B_o)|F|$;
\item bounded fresh residual: $|F|\le q$, in which case exposing the full
five-state equality pattern on $F$ costs at most $q\log 5$.  On each resulting
atom, either the split part $F_{\rm sp}\subseteq F$ is nonempty, is promoted to
old support, and produces a fresh close-pair block inside $F_{\rm sp}$, or
every residual split witness lies in the old support $U$.
\end{enumerate}
\end{lemma}

\begin{proof}
The probabilities $\Prb_\nu[t\in W_L]$ are functions of the current public
law, so $F$ is public.  If $|F|>q$, adding all coordinates of $F$ to the old
support uses no private selector.  Each such coordinate was fresh before the
transition.  In the potential of Section~\ref{sec:potential} this
changes $\Phi_{\rm fresh}+\Phi_{\rm old}$ by
$-(B_f-B_o)|F|$, which pays for the promotion and the bounded public
bookkeeping attached to it.

Assume $|F|\le q$.  For every coordinate of $F$ expose one of the five
equality patterns
\[
  XYZ,\quad XY|Z,\quad XZ|Y,\quad YZ|X,\quad X|Y|Z .
\]
There are at most $5^{|F|}\le 5^q$ atoms.  If some coordinates of $F$ are
split, let $F_{\rm sp}$ be the nonempty public split part on that atom.  Promote
	$F_{\rm sp}$ to $U$ before entering the close-pair child.  Since
	$F_{\rm sp}\subseteq [k]\setminus U$, this gives a drop
	$(B_f-B_o)|F_{\rm sp}|$, which pays the bounded pattern exposure and deposits
	scale for the fresh close-pair block after the constant hierarchy is chosen.
	The constant $C_{\rm fr}(q)$ is chosen per promoted coordinate; since
	$|F_{\rm sp}|\ge1$ and $|F|\le q$, this single promoted coordinate already
	pays the whole bounded $F$-pattern transcript.
One of the three split roles occupies at least one third of $F_{\rm sp}$,
giving a fresh close-pair block inside $F_{\rm sp}$ after relabelling.
If no coordinate of $F$ is split on the atom, then every residual split
witness in $L$ belongs to $U$, by the definition of $F$ and the fact that
$W_L$ is nonempty pointwise under Lemma~\ref{lem:residual-pressure}.  This is
the old-support branch.
\end{proof}

\begin{proposition}[Residual split-count DRC]
\label{prop:residual-drc}
In the setup of Lemma~\ref{lem:residual-pressure}, suppose
\[
  \EE_\nu |W_L|\ge \beta |L|
\]
for $0<\beta\le 1$.  Then for every integer
$1\le s\le \beta |L|/4$ there is a block $J\subseteq L$, $|J|=s$, such that
\[
  \nu[t\in W_L\text{ for every }t\in J]
  \ge
  \frac{\beta}{2}\left(\frac{\beta}{4}\right)^s.
\]
After exposing one role vector, there is a subblock $J'\subseteq J$ with
$|J'|\ge |J|/3$ and one ordered pair of samples, say $(X,Y)$ after relabelling,
such that on a subevent of mass at least
\[
  \frac{\beta}{2}\left(\frac{\beta}{12}\right)^s
\]
one has
\[
  X_t=Y_t\ne Z_t\qquad (t\in J').
\]
\end{proposition}

\begin{proof}
Apply Lemma~\ref{lem:drc-row-block} to the probability space of triples under
$\nu$, the index set $L$, and the events $E_t=\{t\in W_L\}$.  This gives the
block $J$ and the first mass bound.

On the event that every $t\in J$ splits, each coordinate has one of three
roles: $XY|Z$, $XZ|Y$, or $YZ|X$.  There are at most $3^s$ role vectors, so
some role vector has conditional mass at least $3^{-s}$.  On this atom, one
role appears on at least $s/3$ coordinates; call that majority subblock
$J'$.  Relabeling the samples gives the displayed close-pair relation.
\end{proof}

\begin{proposition}[Residual decision order]
\label{prop:residual-decision-order}
At a residual node, after exact split exposure on an active block $I$, put
$R=I\setminus D$, $L=[k]\setminus R$, and
\[
  S=\EE_\nu |W_L|.
\]
Fix constants $\beta>0$ and $q\ge1$.  The residual transition is run in the
following public order:
\begin{enumerate}[label=\textup{(\roman*)}]
\item apply Lemma~\ref{lem:residual-support-promotion} to the fresh residual
support;
\item if that lemma produces a bounded fresh close-pair block, leave residual
mode through close-pair mode;
\item otherwise the branch is old-supported, meaning every residual split
witness is contained in the current old support $U$ pointwise; no random old
positive-support selector is exposed;
\item on an old-supported branch, if $S\ge\beta |L|$ and
$|L|\ge4/\beta$, apply Proposition~\ref{prop:residual-drc} with
\[
  s_L=\max\{1,\lfloor\beta |L|/8\rfloor\};
\]
the resulting role block is old and is registered in the old-window ledger;
\item on an old-supported branch with $S\ge\beta |L|$ but $|L|<4/\beta$, expose
the bounded set $W_L\subseteq L$ directly and route the bounded support through
the old-incidence or bounded-leaf ledger;
\item on an old-supported branch with $S<\beta |L|$, apply
Lemma~\ref{lem:old-residual-routing} to the deterministic old positive
support.
On each old-pattern atom the resulting split set $K$ is classified by
Proposition~\ref{prop:sparse-residual-transition}: deletion-trapped atoms give
scale descent, while the remaining old atoms enter the old-window/close-pair
ledger.
\end{enumerate}
No DRC block is used unless the linear split-count hypothesis
$S\ge\beta |L|$ is satisfied and the resulting DRC range contains a nonempty
block.
\end{proposition}

\begin{proof}
The fresh residual support $F$ and the number $S$ are functions of the
current public law, so the case distinction is public.  If $|F|>q$,
Lemma~\ref{lem:residual-support-promotion} promotes every coordinate of $F$
to the global old support.  After this update, every coordinate with positive
probability of belonging to $W_L$ is old.  If $|F|\le q$, the same lemma
exposes only the five-state pattern on $F$, at cost at most $q\log 5$; each
atom either promotes its nonempty split part and gives a fresh close-pair block
paid by that bounded fresh drop, or leaves all residual witnesses in $U$.

Thus no residual branch exposes the whole random set $W_L\subseteq L$ as an
unpaid sparse selector.  On the remaining old-supported branches, the dense case
$S\ge\beta |L|$ with $|L|\ge4/\beta$ satisfies the non-vacuous hypothesis of
Proposition~\ref{prop:residual-drc} for the displayed value of $s_L$.  Indeed
$s_L\ge1$ and $s_L\le\beta |L|/4$.  Since every split witness is old, the DRC
block and its majority role block are old blocks and are entered into the
old-window incidence ledger.

If $S\ge\beta |L|$ but $|L|<4/\beta$, then the whole residual domain has bounded
size.  Exposing the exact subset $W_L\subseteq L$ and its role pattern has
bounded entropy depending only on $\beta$.  Because the branch is old-supported,
this bounded support is registered in the old ledger or stops through the
bounded-leaf/high-revisit exit.  No DRC block is needed in this bounded case.

If $S<\beta |L|$, the branch is still old-supported and
Lemma~\ref{lem:old-residual-routing} applies directly to the deterministic old
positive support $O=\{u\in U:\Prb[u\in W_L]>0\}$, which is a function of the
public law rather than a random selector.  This exposes only old-coordinate
pattern data and charges it to old incidences.  On each resulting atom the old
split set $K=W_L$ is known inside the already paid old pattern.  Proposition~\ref{prop:sparse-residual-transition}
then separates deletion-trapped scale descents from old close-pair returns.
The last sentence is part of the stated order and prevents the non-linear
residual pressure $S\ge1$ from being used as a linear DRC input.
\end{proof}

\begin{lemma}[Fresh-old-deletion trichotomy]
\label{lem:fresh-old-trichotomy}
Let $J$ be the current block, $D\subseteq J$ the deletion set, $U$ the old
support, and let
\[
  K\subseteq D\cup ([k]\setminus J)
\]
have size $t$.  Define
\[
  \operatorname{Fresh}(K)=K\setminus (J\cup U),
  \qquad
  \operatorname{Old}(K)=K\cap U.
\]
Then at least one of the following holds:
\[
  t\le 4|D|,\qquad
  |\operatorname{Fresh}(K)|\ge t/4,\qquad
  |\operatorname{Old}(K)|\ge t/2.
\]
\end{lemma}

\begin{proof}
If $t\le4|D|$, the first alternative holds.  Hence assume
that $t>4|D|$.
Since $K\cap J\subseteq D$, one has
 $|K\setminus J|>3t/4$.  If also
$|\operatorname{Fresh}(K)|<t/4$, then the remaining coordinates of
$K\setminus J$ lie in $U$, giving
\[
  |\operatorname{Old}(K)|\ge |K\setminus J|-|\operatorname{Fresh}(K)|>t/2.
\]
No disjointness assumption between $U$ and $J$ is used.  The partition is only
on $K\setminus J$: on this set
$(K\setminus J)\setminus U=\operatorname{Fresh}(K)$, while
$(K\setminus J)\cap U\subseteq\operatorname{Old}(K)$.
\end{proof}

\begin{proposition}[Explicit residual split atom transition]
\label{prop:explicit-residual-atom}
At a residual node, let $I$ be the current active block and let
$D=\Split_I(X,Y,Z)$ be the exact split set already exposed on the branch.
Put $R=I\setminus D$ and $L=[k]\setminus R$.  Suppose that a further public
atom exposes
\[
  K=W_L(X,Y,Z)\subseteq L,\qquad |K|=t .
\]
Then $K\subseteq D\cup([k]\setminus I)$ and every coordinate of $K$ is split.
With respect to the current old support $U$, at least one of the following
transitions is available:
\begin{enumerate}[label=\textup{(\arabic*)}]
\item deletion-trapped:
  \[
    t\le 4|D|;
  \]
\item fresh escape:
  \[
    |\operatorname{Fresh}(K)|\ge t/4;
  \]
\item old close-pair return: after routing $K\cap U$ through
  Lemma~\ref{lem:old-residual-routing}, there is
  $O'\subseteq K\cap U$ with $|O'|\ge t/6$ and one ordered pair of samples,
  say $(X,Y)$ after relabelling, such that
  \[
    X_{O'}=Y_{O'},\qquad Z_i\ne X_i\quad(i\in O').
  \]
\end{enumerate}
\end{proposition}

\begin{proof}
Since $L=[k]\setminus(I\setminus D)=D\cup([k]\setminus I)$, the exposed atom
has $K\subseteq D\cup([k]\setminus I)$.  Apply
Lemma~\ref{lem:fresh-old-trichotomy} with the current block $I$ and deletion
set $D$.  The first two alternatives of that lemma give the deletion-trapped
and fresh escape branches.

In the old branch, $|K\cap U|\ge t/2$.  Each coordinate of $K\cap U$ is in
$W_L$, hence is split.  Lemma~\ref{lem:old-residual-routing} applies to the
public old block $K\cap U$ and gives a close-pair subblock of size at least
$(t/2)/3=t/6$, with its role cost entered into the old-window ledger.
\end{proof}

\begin{proposition}[Sparse-exposed residual transition]
\label{prop:sparse-residual-transition}
Let a law on distinct triples $(X,Y,Z)$ satisfy
\[
  |\Split_I(X,Y,Z)|\le s
\]
on a positive-mass branch.  Expose the exact split set
$D=\Split_I(X,Y,Z)$, so $|D|\le s$ and the triples are compatible on
$I\setminus D$.  Suppose residual repair produces a split block
\[
  K\subseteq D\cup([k]\setminus I),\qquad |K|=t.
\]
With respect to an old support $U$, at least one of the following holds:
\begin{enumerate}[label=\textup{(\arabic*)}]
\item deletion-trapped:
  \[
    t\le 4s;
  \]
\item fresh escape:
  \[
    |\operatorname{Fresh}(K)|\ge t/4;
  \]
\item old close-pair return: after routing $K\cap U$ through
  Lemma~\ref{lem:old-residual-routing}, there is
  $O'\subseteq K\cap U$ with $|O'|\ge t/6$ and one ordered pair of the
  samples, say $(X,Y)$ after relabelling, such that
  \[
    X_{O'}=Y_{O'},\qquad Z_i\ne X_i\quad(i\in O').
  \]
\end{enumerate}
\end{proposition}

\begin{proof}
Apply Lemma~\ref{lem:fresh-old-trichotomy} with the deletion set $D$.  Its
first branch gives $t\le 4|D|\le 4s$, and its second branch is the fresh
escape branch.

In the third branch, $|K\cap U|\ge t/2$.  Since $K$ is a split block, every
coordinate in $K\cap U$ is split.  Lemma~\ref{lem:old-residual-routing}
applies to the public old block $K\cap U$ and gives a close-pair subblock of
size at least $(t/2)/3=t/6$.
\end{proof}

\begin{lemma}[Deletion-trapped scale descent]
\label{lem:deletion-scale-descent}
In the setup of Proposition~\ref{prop:explicit-residual-atom} or
Proposition~\ref{prop:sparse-residual-transition}, assume $|I|=m$,
 $|D|\le s=\sigma m$, and the deletion-trapped old-only branch occurs.  Then
every close-pair block produced inside the residual block has size at most
\[
  4\sigma m.
\]
In particular, if $4\sigma<1$, repeated pure deletion-trapped old-only returns
strictly decrease scale geometrically.
\end{lemma}

\begin{proof}
The deletion-trapped branch gives $|K|=t\le 4|D|\le 4s=4\sigma m$.  Any
close-pair subblock selected inside $K$ has size at most $|K|$.  Iterating the same
estimate gives geometric scale decay when $4\sigma<1$.
\end{proof}

\begin{lemma}[Nontrapped residual blocks escape the current block]
\label{lem:nontrapped-escape}
In the setup of Proposition~\ref{prop:explicit-residual-atom} or
Proposition~\ref{prop:sparse-residual-transition}, if the residual block is
not deletion-trapped, so $t>4|D|$, then
\[
  |K\setminus I|>3t/4.
\]
If all coordinates of $K$ are old and split, then after role refinement there
is an old close-pair block $L\subseteq K\setminus I$ with
\[
  |L|>t/4.
\]
\end{lemma}

\begin{proof}
Since $K\subseteq D\cup([k]\setminus I)$, every coordinate of $K$ lying in
$I$ lies in $D$.  Hence $|K\cap I|\le |D|<t/4$, and
\[
  |K\setminus I|>3t/4.
\]
On $K\setminus I$, one of the three split roles occurs on more than
$t/4$ coordinates.  That majority-role set is the required close-pair block.
\end{proof}

\section{Public DRC Blocks and Residual Repair}

DRC-generated blocks must be chosen publicly.

\begin{definition}[Public DRC convention]
\label{hyp:public-drc}
Whenever a DRC argument proves the existence of a coordinate block with a
given mass lower bound, the proof chooses the first valid block in a fixed
public order.  Thus the block is a deterministic function of the current
public law and the fixed parameters.
\end{definition}

Under this convention residual repair blocks are public.  The only private
data left in old-coordinate residuals is the role vector
\[
  \operatorname{Role}_O\in\{XY|Z,XZ|Y,YZ|X\}^O.
\]
Exposing this vector costs at most $|O|\log 3$, and a majority role gives a
close-pair subblock $O'$ with $|O'|\ge |O|/3$.  Hence the role exposure is
locally amortized by the block it creates.

\begin{proposition}[Residual old blocks are public up to role entropy]
\label{prop:residual-public}
In the residual repair route of Proposition~\ref{prop:residual-drc}, enforce
the public DRC convention.  Then every old close-pair parent block has the
following form:
\begin{enumerate}[label=\textup{(\roman*)}]
\item a residual DRC block $K$ is deterministic from the current public law;
\item its old part $O=K\cap U$ is public;
\item the only remaining private selector is a role vector on $O$, whose
entropy is at most $|O|\log 3$;
\item after role refinement, the produced close-pair block has size at least
 $|O|/3$.
\end{enumerate}
\end{proposition}

\begin{proof}
The DRC block is public by the deterministic tie-break.  Since $I,U,K$ are
public at that point, the set $O=K\cap U$ is public as well.  The role vector
has at most $3^{|O|}$ values.  Conditioning on one role vector costs at most
$|O|\log 3$, and a majority role occupies at least one third of $O$.
\end{proof}

\section{Close-Pair Product Lemmas}
\label{sec:close-product}

\begin{lemma}[Conditioned product factorization]
\label{lem:conditioned-product}
For each $j\in J$, let $p_j,q_j$ be probability laws on $\Omega_j$, and put
\[
  \gamma_j=\sum_a p_j(a)q_j(a)<1.
\]
Let $Q=\bigotimes_{j\in J}(p_j\otimes q_j)$ on pairs $(A_J,B_J)$ and let
\[
  E=\{A_j\ne B_j\text{ for all }j\in J\}.
\]
Then $Q(\cdot\mid E)$ factorizes as
\[
  Q_E=\bigotimes_{j\in J}Q_j^*,
  \qquad
  Q_j^*(a,b)=
  \frac{p_j(a)q_j(b)\mathbf 1_{a\ne b}}{1-\gamma_j}.
\]
\end{lemma}

\begin{proof}
The event $E$ is the intersection of coordinate events $E_j=\{A_j\ne B_j\}$,
and under $Q$ these coordinate pairs are independent.  Thus
$Q(E)=\prod_j(1-\gamma_j)$, and division by this product gives the displayed
coordinate factorization.
\end{proof}

\begin{lemma}[Product transfer to near-sunflowers]
\label{lem:product-near-sunflower}
Let $P=\Law(A_J,B_J)$ be supported on $A_j\ne B_j$ for every $j\in J$.  Let
$p_j=\Law(A_j)$, $q_j=\Law(B_j)$, and let $Q_E$ be the conditioned product
law of Lemma~\ref{lem:conditioned-product}.  Assume
\[
  D(P\|Q_E)\le \epsilon,\qquad
  \gamma_j\le \gamma_0<1\quad(j\in J),
\]
and
\[
  \frac1{|J|}\sum_{j\in J}\Coll(p_j)\le \alpha.
\]
If $A^{(1)},A^{(2)},A^{(3)}$ are independent samples from the true
$A$-marginal $P_A$, then
\[
  \Prb\!\left[
    |\Split_J(A^{(1)},A^{(2)},A^{(3)})|
    \le
    \frac{6\alpha |J|}{(1-\gamma_0)^2}
  \right]
  \ge
  \frac12-3\sqrt{\epsilon/2}.
\]
\end{lemma}

\begin{proof}
Under $Q_E$, the $A_j$-marginal is
\[
  p_j^*(a)=\frac{p_j(a)(1-q_j(a))}{1-\gamma_j}.
\]
Since $1-q_j(a)\le 1$ and $1-\gamma_j\ge 1-\gamma_0$,
\[
  \Coll(p_j^*)=\sum_a \frac{p_j(a)^2(1-q_j(a))^2}{(1-\gamma_j)^2}
  \le
  \frac{\Coll(p_j)}{(1-\gamma_0)^2}.
\]
Thus three independent samples from the $A$-marginal of $Q_E$ have expected
split count at most
\[
  \frac{3\alpha |J|}{(1-\gamma_0)^2}
\]
by Lemma~\ref{lem:collision-split}.  Markov's inequality gives probability at
least $1/2$ for the displayed event under that product marginal.

By Pinsker's inequality \cite{Pinsker64,CoverThomas06,CsiszarKorner11},
\[
  \operatorname{TV}(P,Q_E)\le \sqrt{\epsilon/2}.
\]
Total variation cannot increase under marginalization, and the distance
between threefold product laws is at most three times the one-sample
distance.  Hence the same event has probability at least
$1/2-3\sqrt{\epsilon/2}$ under $P_A^3$.
\end{proof}

\begin{lemma}[Projection near-sunflowers lift to code triples]
\label{lem:projection-lift}
Let $\mu$ be a law on a code $\C$, let $\pi_J$ be projection to $J$, and let
$\nu=(\pi_J)_*\mu$.  If a projected event $F\subseteq(\pi_J\C)^3$ has
\[
  \nu^3(F)\ge \rho,
\]
then at least one of the following holds:
\begin{enumerate}[label=\textup{(\roman*)}]
\item the first codeword $x$ in the fixed public order satisfying
$\mu(x)\ge \rho/6$ is retained as a terminal codeword atom;
\item for independent $X,Y,Z\sim\mu$, the event
$(\pi_JX,\pi_JY,\pi_JZ)\in F$ occurs with $X,Y,Z$ distinct with probability
at least $\rho/2$.
\end{enumerate}
\end{lemma}

\begin{proof}
The projections of independent samples from $\mu$ are independent samples
from $\nu$, so the projected event has probability exactly $\rho$ at the code
level before imposing distinctness.  Let
\[
  \Delta=\Prb[X=Y]=\sum_x\mu(x)^2.
\]
The probability that $X,Y,Z$ are not all distinct is at most $3\Delta$.  If
$\Delta\ge\rho/6$, then $\max_x\mu(x)\ge \Delta\ge\rho/6$.  The retained atom
is chosen by the fixed public order and the discarded complement is charged by
the terminal-pruning convention, so no nonterminal child inherits that
complement.  Otherwise $3\Delta<\rho/2$, and the distinct lift has mass at
least $\rho/2$.
\end{proof}

\begin{theorem}[Close-pair master exit]
\label{thm:close-pair-exit}
Let $\mu$ be the current one-sample marginal law on $\C$ at a tuple node, and
let $J$ be a public close-pair block.  Suppose $P=\Law(A_J,B_J)$ is supported on
\[
  A_j\ne B_j\qquad(j\in J),
\]
where the $A$-marginal is the projection of the current marginal law:
$A_J=\pi_J(X)$ for $X\sim\mu$.  Let
\[
  p_j=\Law(A_j),\qquad q_j=\Law(B_j),\qquad
  \gamma_j=\sum_a p_j(a)q_j(a),
\]
and, if all $\gamma_j<1$, let $Q_E$ be the conditioned product law from
Lemma~\ref{lem:conditioned-product}.

Fix parameters $h,\alpha,\epsilon,\gamma_0$ with $0\le\alpha\le 1$,
$0\le\gamma_0<1$, and
\[
  \rho=\frac12-3\sqrt{\epsilon/2}>0.
\]
Then at least one of the following holds:
\begin{enumerate}[label=\textup{(\arabic*)}]
\item low common-symbol entropy:
  \[
    \bfH(A_J)\le h|J|;
  \]
\item cross-collision:
  \[
    \gamma_j>\gamma_0\quad\text{for some }j\in J;
  \]
\item product KL excess (only in the case all $\gamma_j<1$):
  \[
    D(P\|Q_E)>\epsilon;
  \]
\item high $A$-side collision:
  \[
    \frac1{|J|}\sum_{j\in J}\Coll(p_j)>\alpha;
  \]
\item terminal codeword atom:
  the first codeword in the fixed public order with
  \[
    \mu(x)\ge \rho/6
  \]
  is retained as a terminal atom;
\item law-level near-sunflower return:
  for independent $X,Y,Z\sim\mu$,
  \[
    \Prb\!\left[
      X,Y,Z\text{ distinct and }
      |\Split_J(X,Y,Z)|
      \le
      \frac{6\alpha |J|}{(1-\gamma_0)^2}
    \right]
    \ge \rho/2.
  \]
\end{enumerate}
\end{theorem}

\begin{proof}
If some $\gamma_j=1$, then the cross-collision alternative holds because
$\gamma_0<1$.  Hence we may assume all $\gamma_j<1$, so $Q_E$ is defined by
Lemma~\ref{lem:conditioned-product}.

If any of the first four alternatives holds, there is nothing to prove.
Assume they all fail.  Then
\[
  D(P\|Q_E)\le\epsilon,\qquad
  \gamma_j\le\gamma_0\quad(j\in J),
\]
and
\[
  \frac1{|J|}\sum_{j\in J}\Coll(p_j)\le\alpha.
\]
Lemma~\ref{lem:product-near-sunflower} gives projected near-sunflower mass at
least $\rho$ for three independent samples from the $A$-marginal.  Since this
$A$-marginal is the $J$-projection of $\mu$, Lemma~\ref{lem:projection-lift}
gives either a terminal codeword atom of mass at least $\rho/6$ or a distinct
law-level lift of mass at least $\rho/2$.  These are alternatives (5) and
(6).  In the tuple proof tree, the latter is realized by conditioned copying
from the same public key, not by asserting an ambient-code density increment.
\end{proof}

\begin{theorem}[Positive KL public stopping scan]
\label{thm:positive-kl-public}
Let
\[
  V=(V_1,\ldots,V_m)
\]
be a list of tuple-coordinate variables, ordered by a public order; each
$V_t$ is one actual sample coordinate value, with its sample label and ambient
coordinate determined by the public list.  Let $P\ll Q$ be laws on the product
space of $V$ with $D(P\|Q)>0$.  Fix a prefix cost-density threshold
$\lambda>0$, a collision threshold $0<\tau<1$, and a tail threshold
$0<\rho\le 1/2$.

Then the following public stopping scan gives an exhaustive list of retained or
charged-pruned children at the current tuple node.  Along a realized branch, let
$t$ be the first index, if one exists, whose current conditional one-coordinate
KL is positive, and let $i$ be the first index, if any, for which the realized prefix
$V_{\le i}=u_{\le i}$ has
\[
  -\log P(V_{\le i}=u_{\le i})>\lambda i .
\]
If the positive-KL index $t$ occurs before such a crossing, the branch is routed
as in \textup{(1)} below.  If a first crossing occurs before the positive-KL
stop, the branch is routed as in \textup{(2)}.  If neither stop occurs through
the end of the list, the full tuple prefix is terminalized as a projected shadow
with
\[
  -\log P(V=v)\le \lambda m .
\]
In the first two cases the output is one of the following:
  \begin{enumerate}[label=\textup{(\arabic*)}]
\item an amortized projected prefix $V_{<t}=u$ with
  \[
    -\log P(V_{<t}=u)\le \lambda(t-1),
  \]
  followed by a one-coordinate positive-KL tail.  After conditioning on the
  public positive tail of mass at least $\rho$, or after pruning a positive
  tail of mass $<\rho$ at cost at most $2\rho$, this gives either a terminal
  projected tuple pin of conditional mass at least $\rho\tau$ in the sense of
  Lemma~\ref{lem:tuple-pin-shadow} or a low-collision conditional hash law
  whose paid prefix and tail/pruning transcript is within the fixed
  hash-window budget.  If the set-valued tail/pruning increment would overflow
  the budget, the branch stops on the already paid value prefix instead of
  treating that set event as a projected shadow;
\item an earlier first-crossing coordinate $i$ with an amortized prefix
  $V_{<i}=u_{<i}$.  At this coordinate the proof peels the public heavy atoms
  of the conditional law.  Heavy-atom branches terminalize as projected tuple
  pins of conditional mass at least $\rho\tau$; the light complement is either
  a charged pruning branch of mass $<\rho$, or, after conditioning on that
  complement, a low-collision conditional hash law whose paid prefix and
  complement transcript is within the fixed hash-window budget.  If the
  complement event would overflow the budget, the branch stops on the already
  paid value prefix.
  \end{enumerate}
All coordinate identities and sample labels used here are determined by the
public order and the charged prefix value.  The conclusion is tuple-level: a
pin is terminalized as a projected shadow including the amortized prefix used
to reach it, with mixed-label prefixes projected by
Lemma~\ref{lem:tuple-prefix-shadow}, and a hash law remains inside the paid
public transcript and the hash window of that sample label until conditioned
  copying is invoked.  Set-valued tail and pruning events are retained only
  inside this paid hash transcript; terminal projected shadows contain only
  coordinate-value pins.  Every
collision statement in the conclusion is for the actual
$P$-conditional law of the selected tuple-coordinate variable under the paid
prefix and tail/pruning event; the reference law $Q$ is used only to locate a
positive-KL tail.  First crossings and the final full-prefix terminal branch are
driven by the $P$-prefix mass, not by any low-collision claim about $Q$.
If an amortized prefix is nonempty but
  exceeds the fixed hash-window transcript budget, the branch terminalizes on
  that projected tuple-coordinate value prefix instead of entering hash mode.
\end{theorem}

\begin{proof}
For a realized branch write
\[
  C_r=-\log P(V_{\le r}=v_{\le r}).
\]
By Lemma~\ref{lem:kl-chain}, the conditional KL functions
\[
  d_t(u)=D(P(V_t\in\cdot\mid V_{<t}=u)\|
          Q(V_t\in\cdot\mid V_{<t}=u))
\]
are public after the prefix value $u$ is paid.  Scan the realized prefixes in
the fixed public order.  If the first positive-KL index $t$ occurs before any
prefix crossing, then $C_{t-1}\le\lambda(t-1)$ and the prefix
$u=V_{<t}$ is amortized.  Under that prefix, the one-coordinate laws $p,q$
satisfy
\[
  D(p\|q)>0.
\]
The positive tail $S_+=\{a:p(a)>q(a)\}$ has positive $p$-mass by
Lemma~\ref{lem:positive-tail}.  If $p(S_+)\ge\rho$, Lemma~\ref{lem:fibre-hash}
applied to $S_+$ gives either a value atom of conditional mass at least
$\rho\tau$ or a low-collision conditional law.  The value atom is
terminalized by Lemma~\ref{lem:tuple-pin-shadow}.  A low-collision law enters
  hash mode only after the public tail event $V_t\in S_+$ is entered into the
  transcript; the collision bound is for
  $P(V_t\in\cdot\mid V_{<t}=u,\ V_t\in S_+)$, not for the corresponding
  $Q$-conditional law.  This costs at most $\log(1/\rho)$ and is a fixed part of
  the window budget.  If the already paid prefix together with this fixed tail
  cost exceeds the fixed window budget, the set event $V_t\in S_+$ is not
  appended; the branch terminalizes on the projected value prefix
  $V_{<t}=u$ instead.  The constant hierarchy chooses $B_{\rm win}$ larger than
  every fixed tail/pruning increment, so an empty prefix never causes this
  overflow.  If
$p(S_+)<\rho$, prune $S_+$ at cost at most $2\rho$ by
Lemma~\ref{lem:positive-tail}.  This means that the branch
$V_t\in S_+$ is an explicit pruning branch with charged discarded mass, while
the retained branch is the public complement $V_t\notin S_+$.  The remaining
public set has $p$-mass at least $1-\rho\ge\rho$, so applying
Lemma~\ref{lem:fibre-hash} there gives the same value-pin/hash alternative,
  now under the public pruned event $V_t\notin S_+$ and again for the actual
  $P$-conditional law, provided the fixed pruning increment fits in the
	  remaining window budget.  If it does not fit, the branch terminalizes on the
	  projected value prefix $V_{<t}=u$ before the set event is appended.  This is
	  alternative (1).

If a prefix crossing occurs before the positive-KL stop, or if no positive-KL
stop has occurred before the first crossing, let $i$ be the first index with
$C_i>\lambda i$.  Then the previous prefix is amortized:
$C_{i-1}\le\lambda(i-1)$.  Let $p$ be the one-step conditional law of $V_i$
under that previous prefix; this law is attached to the public coordinate
position $i$.  Apply Lemma~\ref{lem:heavy-or-light-hash} to this law.  The
heavy set
\[
  H=\{a:p(a)\ge\rho\tau\}
\]
has bounded size, and each branch $V_i=a\in H$ is a genuine terminal value pin
of conditional mass at least $\rho\tau$.  These branches are converted to
projected tuple pins by Lemmas~\ref{lem:tuple-pin-shadow} and
\ref{lem:tuple-prefix-shadow} together with the amortized prefix.  On the
light complement $L$, if $p(L)<\rho$ then the branch $V_i\in L$ is an explicit
charged pruning branch and is not inherited by a cylinder child.  If
$p(L)\ge\rho$, then the conditional law of $V_i$ on $L$ has collision at most
$\tau$.  This gives a low-collision hash coordinate carried by the paid prefix
and the public complement event $V_i\in L$, provided that transcript remains
	inside the fixed window budget.  If it would overflow the budget, the set event
	is not appended and the branch terminalizes on the already paid value prefix
	$V_{<i}=u_{<i}$.  This is alternative (2).

It remains to cover branches on which neither a positive-KL stop nor a prefix
crossing occurs.  Then $C_m\le\lambda m$, so the full tuple-coordinate value
$V=v$ is an amortized projected shadow.  This terminal branch is included in the
shadow ledger and has no hash output.  Thus the scan is exhaustive: every
realized branch either stops by positive KL, by first crossing, by an explicit
charged pruning branch inside those cases, or by the final amortized full
prefix.
\end{proof}

\begin{theorem}[Product-KL branch theorem]
\label{thm:product-kl-branch}
\label{prop:product-kl-transfer}
In a close-pair tuple node, let $J=(j_1,\ldots,j_m)$ be the public close-pair
block and suppose that the product-KL alternative of
Theorem~\ref{thm:close-pair-exit} holds:
\[
  D(P\|Q_E)>\epsilon .
\]
Run the public stopping scan of Theorem~\ref{thm:positive-kl-public} on the
interleaved actual tuple-coordinate list
\[
  V=(A_{j_1},B_{j_1},A_{j_2},B_{j_2},\ldots,A_{j_m},B_{j_m}).
\]
Then every realized branch is routed to exactly one of the following outputs.
\begin{enumerate}[label=\textup{(\roman*)}]
\item \textup{Amortized prefix terminal.}  A tuple-coordinate prefix
$V_{\le r}=u_{\le r}$, or the full prefix $V=u$, is terminalized as a projected
multi-sample shadow.  Its local mass cost is
\[
  -\log P(V_{\le r}=u_{\le r})\le \lambda r
\]
for the prefix actually retained; mixed sample labels in the prefix are grouped
by Lemma~\ref{lem:tuple-prefix-shadow}.  No hash coordinate and no hash token is
created.
\item \textup{One-coordinate terminal pin.}  After an amortized prefix and,
possibly, a paid positive-tail or light-complement event, one actual
sample-coordinate value is pinned with conditional mass at least $\rho\tau$.
The terminal shadow contains the paid prefix and this coordinate value, and its
local entropy is charged to the terminal-shadow/pruning ledger.  No nonterminal
child inherits the complement.
\item \textup{Charged pruning branch.}  A positive tail or light complement has
conditional mass below $\rho$.  The discarded branch is charged by
$-\log(1-q)\le 2\rho$, where $q<\rho$, and is not inherited by a cylinder child.
The retained complement may be used only as a paid set event inside the hash
transcript; if appending it would exceed the remaining $B_{\rm win}$ budget, the
branch falls back to the amortized prefix terminal in \textup{(i)}.
\item \textup{Labelled low-collision hash coordinate.}  For some actual sample
label $s\in\{X,Y,Z\}$ and some ambient coordinate $j_r$, the paid prefix and one
public tail/complement event fit inside the labelled window budget
$B_{\rm win}$, and the selected actual one-coordinate conditional law satisfies
\[
  \Coll\bigl(\Law(X^s_{j_r}\mid \textup{paid prefix and set event})\bigr)
  \le \tau .
\]
The coordinate enters only the hash window for the public label $s$.  The
reference law $Q_E$ and the pair alphabet $(A_j,B_j)$ are not passed to the hash
ledger.
\end{enumerate}
\end{theorem}

\begin{proof}
List the public close-pair block as $J=(j_1,\ldots,j_m)$.  In close-pair mode,
after relabelling the active samples, the tuple coordinates on $J$ have the
form
\[
  X_J=Y_J=A_J,\qquad Z_j=B_j\ne A_j\quad(j\in J).
\]
Form the publicly ordered interleaved tuple-coordinate list
\[
  V=(A_{j_1},B_{j_1},A_{j_2},B_{j_2},\ldots,A_{j_m},B_{j_m}).
\]
This is only a bijective reordering of the pair data $(A_J,B_J)$.  Therefore,
if $P_V,Q_V$ are the pushforwards of $P,Q_E$ under this interleaving, then
\[
  D(P_V\|Q_V)=D(P\|Q_E)>\epsilon>0.
\]
Apply Theorem~\ref{thm:positive-kl-public} to the laws $P_V,Q_V$ and to this
public list.  Its stopping scan is exhaustive on the realized tuple-prefix
tree, so every branch is one of the four cases listed in the statement.  On
every hash alternative, the selected local variable is now an actual
tuple-coordinate value: either the common value
$A_{j_r}=X_{j_r}=Y_{j_r}$, or the third-sample value
$B_{j_r}=Z_{j_r}$.  The prefix variables are also actual tuple-coordinate
values in the same public list.  The low-collision alternative produced by
Theorem~\ref{thm:positive-kl-public} is a statement about the actual
$P_V$-conditional law of this selected tuple coordinate under the paid prefix
and tail/pruning transcript.  The reference product law $Q_V$ is not passed to
the hash ledger.

If Theorem~\ref{thm:positive-kl-public} terminalizes an amortized prefix, the
final amortized full tuple prefix, a positive-tail value atom, or a
first-crossing heavy atom, the corresponding sample-coordinate values are
entered into the terminal projected shadow and paid by the proof-tree entropy
ledger; mixed-label prefixes are grouped into multi-sample shadows by
Lemma~\ref{lem:tuple-prefix-shadow}.  First-crossing light complements are
either charged pruning branches or paid set events inside the hash transcript;
they are not terminal value-shadows.  If the theorem produces a low-collision
conditional law, that law is attached to the public sample-coordinate label
$(X,j_r)$ (equivalently $(Y,j_r)$) or $(Z,j_r)$ under the paid public
conditioning transcript.  The branch enters the hash window for that public
sample label only, and only while the paid transcript is inside the fixed
hash-window budget.

Thus the product-KL branch is routed through the KL chain rule on the actual
tuple-coordinate list.  The proof does not choose an arbitrary first
coordinate of $J$, and it does not use a low-collision statement on the pair
alphabet.
\end{proof}

\section{Old Support and Projection}
\label{sec:old}

The old support cannot be paid as an $\exp(O(|U|))$ overhead at every rank.
That would iterate to $\exp(\Theta(k^2))$.  The correct statement is a
projection-or-terminalization criterion.

\begin{definition}[Old projection]
Old role/support analysis data satisfies projection if every bit of that data is
one of the following:
\begin{enumerate}[label=\textup{(\roman*)}]
\item terminal-local, and used only to define a terminal block of comparable
size;
\item public-progress, affecting only rank/fresh/scale/high-revisit progress;
\item not retained before terminal shadow counting;
\item promoted to public state with entropy charged to a telescoping
potential.
\end{enumerate}
\end{definition}

\begin{lemma}[Old projection implies shadow control]
\label{lem:old-projection}
If every old-window segment satisfies old projection, then the old component of the
terminal shadow has entropy
\[
  \bfH(\Theta_{\rm old})
  \le C\,\EE[T_{\rm old}(\Theta)+R_{\rm old}(\Theta)].
\]
\end{lemma}

\begin{proof}
Only terminal-local old data can influence $\Theta_{\rm old}$ after
projection.  For one terminal old block, local role exposure costs $O(|O'|)$,
and the value selector is allowed only in the low-symbol case
$\bfH(A\mid E)\le h|O'|$ of Theorem~\ref{thm:old-symbol-budget}.  Hence the
terminal value contribution is $O(|O'|)$.  High-symbol old value vectors are
not exposed; they are routed as nonterminal old-window data and are not
retained in the terminal shadow unless a later low-symbol terminalization
occurs.  Summing over terminal old exits gives the bound.  Public-progress
and unretained data do not contribute to $\Theta_{\rm old}$.
\end{proof}

\begin{lemma}[High revisit local pressure]
\label{lem:high-revisit}
At a close-pair node, suppose a coordinate $i$ has role
\[
  X_i=Y_i\ne Z_i
\]
and let $A_i=X_i=Y_i$ under the current conditional law.  For every
$0<\rho\le1/2$ and $0<\tau<1$, the following exhaustive public routing is
available.  Put
\[
  H=\{a:\Prb[A_i=a]\ge \rho\tau\},\qquad L=\Omega_i\setminus H .
\]
Then $|H|\le(\rho\tau)^{-1}$, every branch $A_i=a\in H$ is a terminal
one-coordinate projected tuple pin of conditional mass at least $\rho\tau$,
and the complement satisfies exactly one of:
\begin{enumerate}[label=\textup{(\roman*)}]
\item $\Prb[A_i\in L]<\rho$, in which case the complement is a charged pruning
branch of loss at most $2\rho$;
\item $\Prb[A_i\in L]\ge\rho$, in which case, after conditioning on the public
set event $A_i\in L$, the conditional law has collision at most $\tau$ and
therefore gives a public low-collision hash coordinate.
\end{enumerate}
\end{lemma}

\begin{proof}
This is Lemma~\ref{lem:heavy-or-light-hash} applied to the one-coordinate law
of $A_i$.  The only point to note is that the coordinate identity and the sets
$H,L$ are determined by the current public conditional law.  If $L$ is retained,
the set event $A_i\in L$ is appended only as part of the paid hash transcript;
if it would exceed the remaining hash-window budget, the branch terminalizes
before entering the hash window, as in Theorem~\ref{thm:positive-kl-public}.
\end{proof}

\begin{lemma}[Old-pattern container]
\label{lem:old-pattern-container}
Fix an old support $U$ and an ordered triple $(X,Y,Z)$.  For each
$u\in U$, record the equality partition of $\{X_u,Y_u,Z_u\}$; there are five
possibilities:
\[
  XYZ,\quad XY|Z,\quad XZ|Y,\quad YZ|X,\quad X|Y|Z.
\]
After conditioning on one atom of this five-state pattern on $U$, every
future old split block $O\subseteq U$ has a deterministic subblock
$O'\subseteq O$ with $|O'|\ge |O|/3$ and one ordered pair, say $(X,Y)$ after
relabelling, such that
\[
  X_{O'}=Y_{O'},\qquad Z_o\ne X_o\quad(o\in O').
\]
\end{lemma}

\begin{proof}
On a fixed five-state atom, each coordinate of a split block has one of the
three roles $XY|Z,XZ|Y,YZ|X$.  One role occurs on at least one third of the
coordinates.  The corresponding majority set is $O'$.
\end{proof}

\begin{lemma}[Old-supported residual witness routing]
\label{lem:old-residual-routing}
At an old-supported residual node, suppose $W_L\subseteq U$ pointwise and
$W_L$ is nonempty pointwise.  Let
\[
  O=\{u\in U:\Prb_\nu[u\in W_L]>0\}.
\]
Then $O$ is public.  Before any nonterminal continuation, $O$ is registered in
the old-window incidence ledger and the five-state equality pattern on $O$ is
exposed.  On every resulting positive-mass atom, either the registration
triggers the high-revisit exit, or the atom determines the old split set
$K=W_L\subseteq O$, which is nonempty, and produces an old close-pair block
$K'\subseteq K$ with
\[
  |K'|\ge |K|/3
\]
after relabelling.  The pattern cost is charged to the registered old
incidences and hence decreases $\Phi_{\rm old}$ in the non-high-revisit case.

The same conclusion holds if a public old block $K\subseteq U$ is already
known to be split on the branch; then one registers only $K$ and exposes its
three-role vector.
\end{lemma}

\begin{proof}
The set $O$ is determined by the current public law.  Register the incidences
of $O$ and expose the five-state pattern
\[
  XYZ,\quad XY|Z,\quad XZ|Y,\quad YZ|X,\quad X|Y|Z
\]
on $O$.  Since $W_L\subseteq O$ and $W_L$ is nonempty pointwise, the split
coordinates on any positive-mass atom are exactly a nonempty set
$K=W_L\subseteq O$.  If a coordinate crosses the revisit threshold, the
exposed split label gives a close-pair revisit of that coordinate; the
high-revisit exit is then the exhaustive terminal/pruning/hash routing of
Lemma~\ref{lem:high-revisit}.  Otherwise these incidences are debited in
$N_{\rm old}$, so they lower $\Phi_{\rm old}$ in
Section~\ref{sec:potential}.  Among the split coordinates in $K$, one of the
three roles $XY|Z,XZ|Y,YZ|X$ occurs on at least one third; relabelling the
samples gives the close-pair block $K'$.

If a public split block $K$ is already available, the same proof skips the
five-state exposure on $O$ and exposes only the three-role vector on $K$.
\end{proof}

\begin{theorem}[Old-symbol values are terminal-local]
\label{thm:old-symbol-budget}
Fix an old support $U$ and condition on one old-pattern atom.  In the public
normal form, common old-symbol values are exposed only on terminal low-symbol
branches.  More precisely, if a node with event $E$ obtains, after relabelling,
an old close-pair block $O'$ and
\[
  A=X_{O'}=Y_{O'},
\]
then exactly one of the following happens.
\begin{enumerate}[label=\textup{(\roman*)}]
\item The branch is low-symbol:
\[
  \bfH(A\mid E)\le h\,|O'|.
\]
The value selector $A$ may be exposed, and the branch immediately terminates
as a projected shadow which pins the coordinates $O'$.
\item The branch is high-symbol.  The proof records only role, equality-pattern,
scale, and block-incidence data; the value vector $A$ is not exposed and is not
part of the terminal transcript.  The branch is routed to the fixed-scale
old-window analysis of Theorem~\ref{hyp:old-window}.  If a separate
close-pair exit, such as Theorem~\ref{thm:product-kl-branch}, exposes a
bounded tuple-coordinate prefix before entering hash mode, that prefix is
charged to the shadow/hash ledger of that exit and not to the old-window
ledger.
\end{enumerate}
Consequently the old-symbol value contribution to the terminal proof-tree
entropy is bounded by
\[
  \sum_{\textup{terminal low-symbol }v}
    \Prb(E_v)\,\bfH(A_v\mid E_v)
  \le
  h\,\EE T_{\rm old}(\Theta),
\]
where $T_{\rm old}$ is the number of old coordinates pinned by terminal
old-symbol shadows.  No alphabet-dependent term such as
$\bfH(X_U,Y_U,Z_U)$ is used in the global charge.
\end{theorem}

\begin{proof}
The theorem is a normal-form rule plus its entropy accounting.  On a
low-symbol branch the child selector is the terminal value vector $A_v$ (and
the finite role/block data already charged by Lemma~\ref{lem:old-pattern-container}
or Proposition~\ref{prop:old-return-public}).  Its conditional entropy is at
most $h|O'_v|$, and the branch stops with $O'_v$ included in the projected
shadow.  Summing over terminal leaves gives the displayed bound because
$T_{\rm old}(\Theta)$ counts these pinned coordinates with multiplicity in the
same terminal shadow ledger.

On a high-symbol old-window continuation the vector $A_v$ is not a selector of
the old-window proof tree.  Only finite role, pattern, scale, and incidence
information is revealed before the branch is passed to
Theorem~\ref{hyp:old-window}.  Those non-value transcripts are paid by the
old-incidence and scale ledgers.  If another exit later exposes actual
coordinate values, it does so through its own terminal-shadow or bounded
hash-prefix rule; such values are not charged to the old ledger.  Therefore no
alphabet-dependent old-symbol value entropy remains in the old-window account.
\end{proof}

\begin{proposition}[Old close-pair returns are public up to paid old data]
\label{prop:old-return-public}
Suppose a residual transition produces an old block $O\subseteq U$ after an
exact split-set atom and a public DRC choice have been fixed.  Then the
resulting old close-pair return either:
\begin{enumerate}[label=\textup{(\roman*)}]
\item is obtained from a previously fixed old-pattern atom by
Lemma~\ref{lem:old-pattern-container}, so no new coordinate selector is
private;
\item or pays local role entropy $O(|O'|)$ to produce a close-pair block
$O'\subseteq O$ with $|O'|\ge |O|/3$.
\end{enumerate}
In either case the block can be assigned to its public dyadic scale and fed
to the fixed-scale old-window analysis.
\end{proposition}

\begin{proof}
The exact split-set atom and the public DRC convention make the parent block
$O$ a function of the paid public state.  If an old-pattern atom on $U$ is
already fixed, Lemma~\ref{lem:old-pattern-container} chooses the majority-role
subblock deterministically from that atom.  Otherwise expose the local role
vector on $O$; it has at most $3^{|O|}$ values, and a majority role gives
$|O'|\ge |O|/3$, so the exposure cost is $O(|O'|)$.  The resulting coordinate
set has a definite size and is placed in the corresponding fixed-scale old
window.
\end{proof}

\section{Fixed-Scale Old Windows}

The last old-support mechanism is a fixed-scale high-symbol old-role window.

\begin{lemma}[Pooling with a fixed bounded reservoir]
\label{lem:reservoir-pooling}
Fix a reservoir $W=(X^1,\ldots,X^r)$ throughout a window.  Let
$B_1,\ldots,B_R$ be split blocks and suppose each coordinate appears in at
most $Q$ of the blocks.  Put $N=\sum_b |B_b|$.  For every incidence
$(b,i)$ with $i\in B_b$, suppose the branch determines a split label
$(s,t|u)$, meaning
\[
  X^s_i=X^t_i\ne X^u_i.
\]
Then there is one label $(s,t|u)$ and a coordinate set $L$ such that
\[
  X^s_L=X^t_L,\qquad X^u_\ell\ne X^s_\ell\quad(\ell\in L),
\]
and
\[
  |L|\ge \frac{N}{M_r Q},
  \qquad
  M_r=3\binom r3.
\]
\end{lemma}

\begin{proof}
There are $N$ incidences and at most $M_r$ possible split labels.  Some label
occurs at least $N/M_r$ times.  Since each coordinate appears in at most $Q$
blocks, the set of coordinates carrying this fixed label has size at least
$N/(M_rQ)$.  On this set the same two reservoir samples agree and the same
third sample differs.
\end{proof}

\begin{lemma}[Public fresh-reset pruning]
\label{lem:fresh-reset-pruning}
Fix a public state and public blocks $B_1,\ldots,B_R$ chosen before exposing
fresh equality gates.  Let $N=\sum_r |B_r|$.  Suppose each gate satisfies
\[
  \Prb[G_r\mid\text{previous transcript}]\le \exp(-\eta |B_r|)
\]
and that, after the public blocks are fixed and terminal value atoms are
excluded, the auxiliary non-value transcript has at most $\exp(C_{\rm aux}N)$
possible atoms.  If $\eta>C_{\rm aux}$, then the total mass of no-exit reset
windows of incidence $N$ is at most
\[
  \exp(-(\eta-C_{\rm aux})N).
\]
\end{lemma}

\begin{proof}
For one fixed auxiliary transcript atom, the conditional gate probabilities
multiply to at most
\[
  \prod_r \exp(-\eta |B_r|)=\exp(-\eta N).
\]
There are at most $\exp(C_{\rm aux}N)$ such atoms.  The union bound gives the
claim.
\end{proof}

\begin{theorem}[Fixed-scale old-incidence cap]
\label{hyp:old-window}
Let $O_1,\ldots,O_R\subseteq U$ be public old blocks of one active scale and
one public sample-label signature $\ell$ for the ordered active
triple/reservoir labels, and let $N=\sum_r |O_r|$.  The window is run under the
normal-form rule of
Theorem~\ref{thm:old-symbol-budget}: high-symbol common old values are not
exposed by the old-window ledger.  Before any nonterminal continuation from an old block, the proof
reveals only the finite non-value data needed to choose the close-pair role and
then either registers the block incidences or exits by high revisit.  Let
$M(i)$ be the cumulative number of low-revisit registered incidences of
coordinate $i$ on the current branch since $i$ entered the old support, capped
at the revisit threshold.  These counters are not reset when the old support
grows.
The signature $\ell$ is fixed throughout the window; if a transition changes the
active sample-label signature before the incidences are registered, the current
bounded queue is charged to that transition and cleared instead of being
inherited.
Then one of the following occurs:
\begin{enumerate}[label=\textup{(\arabic*)}]
\item a low-symbol old value is exposed and the branch immediately terminalizes
as in Theorem~\ref{thm:old-symbol-budget};
\item high revisit: adding the current block would make some $M(i)>Q$, and the
current close-pair coordinate is routed by Lemma~\ref{lem:high-revisit} into
terminal value-shadow branches, a charged pruned complement, or a public
low-collision hash coordinate;
\item low-revisit continuation: every $M(i)\le Q$, no high-symbol old value is
exposed, and the total non-value old transcript on the branch is bounded by
\[
  C_{\rm old}(Q,r_*)\sum_{i\in U}M(i)
  \le C_{\rm old}(Q,r_*)\,Q|U|.
\]
\end{enumerate}
Thus high-symbol old returns have no alphabet-dependent value-entropy charge,
and their remaining role, pattern, scale, and incidence data are paid by the
old-incidence potential.
\end{theorem}

\begin{proof}
The signature $\ell$ is part of the public old-window state, so all equality
patterns, role vectors, and common-symbol vectors below refer to the same
ordered active labels.  If a later transition replaces these labels before the
window is charged, it is treated as the hard exit which clears the bounded queue;
the old-pattern atom for $\ell$ is not reused for the new labels.

At a reached old block, first apply Theorem~\ref{thm:old-symbol-budget}.  If the
common old-symbol vector has entropy at most $h|O_r|$, its value selector is
terminal-local and the branch stops, giving alternative (1).  Otherwise the
branch is high-symbol and the value vector is not a child selector.

The only selectors exposed before continuation are the equality-pattern, role,
scale, reservoir-label, and incidence data listed in the old rows of the global
charge table.  By definition of $C_{\rm old}(Q,r_*)$, these cost at most
$C_{\rm old}(Q,r_*)$ per registered old incidence.  If adding the incidences of
$O_r$ would make some counter exceed $Q$, the current close-pair occurrence of
that coordinate satisfies the hypothesis of Lemma~\ref{lem:high-revisit}; it
yields terminal coordinate value pins, a charged pruned complement, or a public
low-collision hash coordinate, which is alternative (2).  The current block is
then charged to that exit and is not counted as a further low-revisit
continuation.

Otherwise register $O_r$, increasing $M(i)$ by one for each $i\in O_r$.  If no
counter exceeds $Q$, then
\[
  \sum_{r}|O_r|=\sum_{i\in U}M(i)\le Q|U|.
\]
The displayed cost bound follows.  Since high-symbol values were never exposed,
this low-revisit continuation contributes only old-incidence transcript, and
the weighted debit in $\Phi_{\rm old}=B_o|U|-N_{\rm old}$ pays it.  When these
incidences also create a continuing old active child, the same debit first
reserves its rank and scale balances.  The remaining raw capacity condition is
$N_{\rm old}\le Q\omega_{\rm old}|U|<B_o|U|$ in the constant hierarchy.
\end{proof}

Short old windows are still concatenated so that a branch cannot restart the
same fixed cost indefinitely without increasing the cumulative old-incidence
counter.

\begin{proposition}[Delayed old-window queue]
\label{prop:delayed-old-queue}
Fix an old support $U$, a dyadic old scale, and a public sample-label signature
$\ell$ for the ordered active triple/reservoir labels.  The proof maintains one
public queue $\mathcal Q(\ell,U,\mathrm{scale})$ of old blocks of that scale and
signature.  If the old support later grows from $U$ to $U^+$ while $\ell$ is
unchanged, the same queued blocks are relabelled as a queue for
$(\ell,U^+,\mathrm{scale})$; no new coordinate is added to those blocks and no
selector is re-exposed.  The old-pattern atom is refined nestedly for the same
signature: the atom on $U^+$ is required to restrict to the already paid atom on
$U$, and only the pattern on $U^+\setminus U$ for the same labels is newly
exposed.  If a transition changes the active sample-label signature before the
queue is registered, the bounded unpaid queue is assigned to that transition's
local charge and cleared; it is not inherited by the new signature.  A bounded
short window is appended to the queue and is not charged separately.  The
revisit count of a coordinate is cumulative over the whole queue, not over the
last short window alone.  Once
the queued incidence
\[
  N(\mathcal Q)=\sum_{O\in\mathcal Q}|O|
\]
reaches the fixed old-window threshold, Theorem~\ref{hyp:old-window} is
applied to the concatenated queue.  The queue is cleared after its incidences
are registered and debited in $N_{\rm old}$, or earlier if a terminal, hash, scale, or
high-revisit exit occurs.  A fresh promotion which only enlarges the old
support does not clear the queue if the signature is unchanged; it relabels the
same queued blocks as above.  If a fresh, hash, copy, scale, or birth transition
is instead used as the hard exit paying the bounded unpaid queue or changing the
signature, that bounded incidence is charged to the hard exit before the queue
is cleared.

Consequently the delayed queue has uniformly bounded unpaid incidence between
two charged exits for each signature, and all nonterminal delayed incidences are
eventually debited in $N_{\rm old}$ or assigned to the signature-changing hard
exit.
\end{proposition}

\begin{proof}
All blocks in the queue are public and have the same dyadic scale and the same
public sample-label signature $\ell$.  Monotonic growth of $U$ does not change
the old coordinates already appearing in the queue.  The nested-pattern
convention from
Proposition~\ref{prop:branch-invariants} makes the consistency statement
literal: if $P_U$ is the already paid five-state atom on $U$, then a later
atom on $U^+$ is of the form
\[
  P_U\cap P_{U^+\setminus U}.
\]
where both atoms refer to the same signature $\ell$.  Thus its restriction to
every queued block contained in $U$ is exactly the previously paid atom, and
relabelling a queue from $U$ to $U^+$ preserves the hypotheses of
Theorem~\ref{hyp:old-window}.  The only new entropy is the five-state pattern on
the newly promoted coordinates for the same labels, which is charged at the
fresh promotion or old-incidence step which introduced them.

Concatenating short windows therefore preserves both publicness and scale.
Coordinate multiplicities are counted in the concatenated list; hence high
revisit is triggered by cumulative, not local, multiplicity.  If the queue has
not reached the fixed threshold, its unpaid incidence is bounded by that
threshold.  A pure growth of $U$ with unchanged signature only changes the label
of the same queue and does not create a new unpaid window.  If the active
signature changes, the old equality atom is no longer meaningful for the new
labels, so the queue must be cleared first: if it has reached threshold, apply
Theorem~\ref{hyp:old-window}; otherwise charge its bounded unpaid incidence to
the signature-changing hard exit.  When the queue reaches the threshold, the
old-window theorem either gives a terminal or high-revisit exit, or registers
and debits the queued incidences in $N_{\rm old}$.  If some other hard exit is
declared before threshold, the bounded unpaid incidence is assigned to that
exit's local charge.  After the queue is cleared, no old incidence from it is
charged again.
\end{proof}

\section{Bounded Active Reservoir}

\begin{definition}[Bounded active reservoir convention]
\label{hyp:reservoir}
Nonterminal states retain only a bounded active reservoir of samples.  Value
pins terminalize immediately as multi-sample shadows.  A hash return may
temporarily realize three conditioned copies, but after paying the public
copy tax it retains only the active triple/pair/reservoir required for the
next transition.
\end{definition}

This prevents exponential sample-label growth under repeated conditioned
copying.

\begin{lemma}[Reservoir marginalization is law preserving]
\label{lem:reservoir-marginalization}
Let a transition create auxiliary sample labels only to certify a public event
or to realize conditioned copies.  Let $L$ be the retained active labels after
the transition, and suppose that every future nonterminal selector and every
nonterminal state variable is measurable with respect to the coordinates of
$L$ and the public transcript.  Replacing the full conditional law on all
temporary labels by its marginal law on $L$ preserves the distribution of every
future selector and of every projected terminal shadow involving retained
labels.  It also cannot increase the number of retained sample labels or the
entropy of any retained public key.
\end{lemma}

\begin{proof}
This is the elementary marginalization identity.  For every event $E$
measurable with respect to the retained labels and the public transcript, its
probability under the full conditional law equals its probability under the
marginal on those labels.  Hence all later public decisions, which by hypothesis
are such events or deterministic functions of such events, have the same law
after the auxiliary labels are erased.  Terminal constraints involving erased
labels have already been projected or terminalized before erasure; therefore no
future terminal shadow can require them.  Finally, deterministic erasure and
marginalization cannot create new retained labels and cannot increase Renyi or
Shannon entropy of a retained transcript key.
\end{proof}

\begin{proposition}[Reservoir labels stay bounded]
\label{prop:reservoir}
Under Definition~\ref{hyp:reservoir}, the number of active sample labels in
every nonterminal node is bounded by an absolute constant.  Terminal shadows
satisfy
\[
  R(\Theta)\le T(\Theta)+O(1),
\]
with the $O(1)$ term accounting only for the fixed active reservoir.
\end{proposition}

\begin{proof}
Induct over nonterminal transitions.  Fresh, coordinate-cylinder, scale, and
old transitions use the current bounded reservoir.  A value pin stops and is
counted as a
multi-sample shadow, so it is not carried forward.  A hash return may realize
three conditioned copies, but after the public copy tax is paid the auxiliary
labels are analysis-only data and are not retained.  Lemma~\ref{lem:reservoir-marginalization}
shows that erasing those labels preserves the law of all later selectors.  The
next state keeps only the triple, pair, or bounded reservoir needed for the next
transition.

Every nonempty terminal sample label pins at least one coordinate, giving
$R(\Theta)\le T(\Theta)$ apart from the fixed bookkeeping labels already
present in the reservoir.
\end{proof}

\section{The Delayed-Window Potential}
\label{sec:potential}

Define a finite-state potential
\[
  \Phi=\Phi_{\rm rank}+\Phi_{\rm fresh}+\Phi_{\rm old}+\Phi_{\rm scale}.
\]
The accounting used below is global in the current unpinned ambient coordinate
set $[k]$, not only in the current active block.  Let $m_{\rm act}$ be the total
rank of the active blocks currently carried by the state, after deleting
coordinates already pinned in the public terminal projection; in single-block
modes this is just the active block size.  The term $\Phi_{\rm rank}$ is the explicit form of the
``initial active-rank'' account used to birth the first active blocks and any
DRC/split child whose size is a fixed fraction of its parent active rank.
Let $I_{\rm old}$ be the raw cumulative
old-window incidence count already registered on the current branch, capped by
the high-revisit threshold through the coordinate counters $M(i)\le Q$.  After
the fixed-window parameters $q,h,Q,r_*$ are chosen, let $C_{\rm fr}(q)$,
$C_{\rm hash}$, and $C_{\rm old}(Q,r_*)$ denote the finite transcript constants
specified in Section~\ref{sec:constant-hierarchy}, and put
\[
  \omega_{\rm old}:=C_{\rm old}(Q,r_*)+B_{\rm rk}+B_s+C_{\rm hash}+10 .
\]
The old potential uses the weighted debit
$N_{\rm old}:=\omega_{\rm old}I_{\rm old}$, not the raw incidence count.
Choose constants
\[
  B_o>Q\,\omega_{\rm old}+C_{\rm aux}+10,
  \qquad
  B_f-B_o>C_{\rm fr}(q)+C_{\rm hash}+B_{\rm rk}+B_s+10.
\]
Put
\[
  \Phi_{\rm rank}=B_{\rm rk} m_{\rm act},
  \qquad
  \Phi_{\rm fresh}=B_f|[k]\setminus U|,
  \qquad
  \Phi_{\rm old}=B_o|U|-N_{\rm old},
  \qquad
  \Phi_{\rm scale}=B_s\sum_{\text{active }J}|J|.
\]
The potential is used as an account balance, not as an unlabelled sum of all
historical work.  When a new active scale is born, its $B_s|J|$ scale budget is
transferred from an existing recurrence account: the initial active-rank
balance, a parent active-rank drop, the parent active-scale balance, a fresh
drop, or registered old incidences.
Terminal shadows do not
fund nonterminal scale births, and a DRC certificate pays selector entropy only;
any scale deposit for the DRC-produced block is drawn from the parent active
rank/scale balance or from the fresh/old progress which created the block.  Fixed hash
returns have size $h<m_0$ in the hierarchy, so they are bounded-scale returns
unless their pending token is assigned to one of the existing hard-exit
accounts before a larger child is entered.  Thus a birth relabels already paid
recurrence budget as rank/scale budget and is not allowed to create additional
recurrence potential for free.  The branch quantity $\Delta\Phi$ below is the
charged expenditure from this account after such transfers, not the raw sum of
every temporary increase and decrease in the displayed formula.
When a DRC or split transition replaces an active rank $m$ by a child block of
size at most $\theta_{\rm rk}m$ with $\theta_{\rm rk}<1$, the transfer into the
child's scale balance, while keeping the child's inherited rank balance, is
taken from the rank drop
$B_{\rm rk}(m-\theta_{\rm rk}m)$ before any residual selector cost is charged.
Thus consuming old incidences decreases $\Phi_{\rm old}$ by the weighted debit,
large enough to reserve child rank/scale balance and pay the local old
transcript.
Promoting any
public fresh set $F\subseteq [k]\setminus U$, including residual coordinates
outside the current active block, to old support changes the first two terms by
\[
  -B_f|F|+B_o|F|=-(B_f-B_o)|F|,
\]
so promotion is a genuine potential drop.  The high-revisit rule gives
$I_{\rm old}\le Q|U|$ on nonterminal low-revisit branches, hence
$N_{\rm old}\le Q\omega_{\rm old}|U|<B_o|U|$ before a terminal or
bounded-support exit occurs.

\begin{lemma}[One-step account inequalities]
\label{lem:step-account-inequalities}
After the constants in Section~\ref{sec:constant-hierarchy} are fixed, the
following deterministic inequalities hold for every nonterminal step.  Let
$\varepsilon_{\rm hash}\in\{0,1\}$ indicate whether the step receives the
unique hash token assigned by Proposition~\ref{prop:hash-token-matching}.
\begin{enumerate}[label=\textup{(\roman*)}]
\item If a fresh step promotes a nonempty set $F\subseteq [k]\setminus U$ and
births an active child $J\subseteq F$, then
\[
  (B_f-B_o)|F|
  \ge
  (B_{\rm rk}+B_s)|J|+C_{\rm fr}(q)|F|
  +\varepsilon_{\rm hash}C_{\rm hash}+|F|.
\]
The same inequality without the child-balance term
$(B_{\rm rk}+B_s)|J|$ applies when no child is born.
\item If an old step registers $n\ge1$ old-coordinate incidences and births an
active child $J$ with $|J|\le n$, then
\[
  \omega_{\rm old}n
  \ge
  (B_{\rm rk}+B_s)|J|+C_{\rm old}(Q,r_*)n
  +\varepsilon_{\rm hash}C_{\rm hash}+n.
\]
\item If a parent active rank $m$ creates a public DRC/split child $J$ with
$|J|\le\theta_{\rm rk}m$, then
\[
  B_{\rm rk}(m-|J|)
  \ge
  (B_s+C_{\rm rk}+1)|J|+\varepsilon_{\rm hash}C_{\rm hash}.
\]
The child keeps the inherited $B_{\rm rk}|J|$ rank balance; only the displayed
rank drop is spent.
\item If a deletion-trapped scale step replaces an active block of size $m$ by
blocks of total size at most $4\sigma m$, then its scale drop is at least
\[
  B_s(1-4\sigma)m.
\]
For $m\ge m_0$ this drop dominates the fixed additive selector costs and any
single assigned $C_{\rm hash}$ token by the choice of $m_0$.
\end{enumerate}
\end{lemma}

\begin{proof}
For (i), use $|J|\le |F|$, $\varepsilon_{\rm hash}\le1$, and
\[
  B_f-B_o>C_{\rm fr}(q)+C_{\rm hash}+B_{\rm rk}+B_s+10 .
\]
For (ii), use $|J|\le n$, $\varepsilon_{\rm hash}\le1$, and
\[
  \omega_{\rm old}=C_{\rm old}(Q,r_*)+B_{\rm rk}+B_s+C_{\rm hash}+10 .
\]
For (iii), $|J|\le\theta_{\rm rk}m$ and the hierarchy gives
\[
  B_{\rm rk}(1-\theta_{\rm rk})m
  >(B_s+C_{\rm rk}+C_{\rm hash}+10)\theta_{\rm rk}m
  \ge (B_s+C_{\rm rk}+1)|J|+\varepsilon_{\rm hash}C_{\rm hash},
\]
using $|J|\ge1$ whenever a child is born.
The last item is exactly Lemma~\ref{lem:deletion-scale-descent} together with
the lower bound imposed on $m_0$.
\end{proof}

\begin{lemma}[Terminal and pruning charge case]
\label{lem:terminal-pruning-charge-case}
Let a step terminate by a projected value shadow, an amortized tuple prefix, a
low-symbol old value, or a codeword atom.  Let $T_v,R_v$ be the number of
coordinate pins and sample labels added to the terminal shadow at this step, and
let $\ell_v$ be any retained pruning-normalization loss charged at the same
step.  Then the step has no nonterminal child balance and no hash token to
transfer, and its selector contribution is charged as
\[
  c_v^{\rm sh}+\ell_v
  \le C_{\rm sh}(T_v+R_v)+\ell_v .
\]
The pruning term is included in $L_{\rm pr}$ in
Theorem~\ref{thm:shadow-potential}; the shadow term is included in
$T(\Theta)+R(\Theta)$.
\end{lemma}

\begin{proof}
Every listed terminal branch stops before a child state is entered.  The
terminal object retained for final counting is only the projected value shadow;
analysis-only data is marginalized.  The local selector entropy of a retained
coordinate value, codeword atom, or amortized prefix is exactly the
conditional-mass cost of the retained terminal event and is bounded by the
number of pinned coordinates and active sample labels up to the fixed terminal
constant $C_{\rm sh}$.  Any discarded complement is not inherited by a child and
is recorded as $\ell_v$ in the pruning ledger.
\end{proof}

\begin{lemma}[Fresh promotion charge case]
\label{lem:fresh-charge-case}
Suppose a nonterminal exit promotes a nonempty public set
$F\subseteq [k]\setminus U$ to the old support.  Let $J\subseteq F$ be the
continuing active child born from this promotion, with $J=\varnothing$ if there
is no child, and let $\varepsilon_{\rm hash}\in\{0,1\}$ denote whether this hard
exit receives the unique hash token assigned by
Proposition~\ref{prop:hash-token-matching}.  The available fresh drop is
\[
  A_{\rm fr}(F)=(B_f-B_o)|F|.
\]
It satisfies
\[
  A_{\rm fr}(F)
  \ge
  C_{\rm fr}(q)|F|+(B_{\rm rk}+B_s)|J|
  +\varepsilon_{\rm hash}C_{\rm hash}+|F|.
\]
Thus the same fresh exit pays its selector transcript, any rank/scale birth
deposit for $J$, and at most one assigned hash token.
\end{lemma}

\begin{proof}
This is Lemma~\ref{lem:step-account-inequalities}(i).  The change in
$\Phi_{\rm fresh}+\Phi_{\rm old}$ caused by promoting $F$ is
$-(B_f-B_o)|F|$, so $A_{\rm fr}(F)$ is precisely the spendable account drop
before old incidences are counted.
\end{proof}

\begin{lemma}[Old-incidence charge case]
\label{lem:old-charge-case}
Suppose a nonterminal old exit registers $n\ge1$ old-coordinate incidences and
births a continuing active child $J$ with $|J|\le n$, allowing $J=\varnothing$.
Let $\varepsilon_{\rm hash}\in\{0,1\}$ indicate whether this exit receives an
assigned hash token.  The available old debit is
\[
  A_{\rm old}(n)=\omega_{\rm old}n .
\]
It satisfies
\[
  A_{\rm old}(n)
  \ge
  C_{\rm old}(Q,r_*)n+(B_{\rm rk}+B_s)|J|
  +\varepsilon_{\rm hash}C_{\rm hash}+n .
\]
Consequently old-window role/signature/reservoir transcripts, old-born
rank/scale deposits, high-revisit selector costs, and one possible hash token
are all paid by the same registered-incidence debit.
\end{lemma}

\begin{proof}
This is Lemma~\ref{lem:step-account-inequalities}(ii).  The low-revisit branch
increments $N_{\rm old}$ by $\omega_{\rm old}n$, giving the displayed drop in
$\Phi_{\rm old}=B_o|U|-N_{\rm old}$ while $U$ is fixed.  When $U$ grows, the
increase $B_o|F|$ in $\Phi_{\rm old}$ has already been deducted in the fresh
case through the net drop $(B_f-B_o)|F|$.
\end{proof}

\begin{lemma}[Parent-rank DRC or split birth case]
\label{lem:rank-birth-charge-case}
Suppose a parent active rank $m$ creates a public DRC/split child
$J$ with $|J|\le\theta_{\rm rk}m$.  If the child is nonempty, the available
parent-rank drop
\[
  A_{\rm rk}(m,J)=B_{\rm rk}(m-|J|)
\]
satisfies
\[
  A_{\rm rk}(m,J)
  \ge
  (B_s+C_{\rm rk}+1)|J|+\varepsilon_{\rm hash}C_{\rm hash}.
\]
The child keeps its inherited rank balance $B_{\rm rk}|J|$; only the displayed
drop is spent to pay the DRC/split selector cost, the new scale balance, and at
most one assigned hash token.
\end{lemma}

\begin{proof}
This is Lemma~\ref{lem:step-account-inequalities}(iii).  The inequality
$|J|\le\theta_{\rm rk}m$ is part of the public DRC/split child rule, or follows
from the fallback $m_0>h/\theta_{\rm rk}$ for a fixed hash block treated as a
parent-rank child.
\end{proof}

\begin{lemma}[Deletion-trapped scale charge case]
\label{lem:scale-charge-case}
Suppose a deletion-trapped transition replaces an active block of size $m\ge
m_0$ by blocks of total size at most $4\sigma m$.  The available scale drop is
at least
\[
  A_{\rm sc}(m)=B_s(1-4\sigma)m .
\]
By the choice of $m_0$, this quantity pays the exact split/role selector cost
at that scale, any fixed additive scale selector cost, and at most one assigned
$C_{\rm hash}$ token.  The descendants inherit only the remaining deposited
scale/rank balances; no positive increase of $\Phi_{\rm scale}$ is counted as a
drop.
\end{lemma}

\begin{proof}
The drop formula is Lemma~\ref{lem:step-account-inequalities}(iv).  The
constant hierarchy chooses $m_0$ after $C_{\rm hash}$ and the fixed scale
selector constants, so the linear drop $B_s(1-4\sigma)m$ dominates every fixed
additive term for all $m\ge m_0$.
\end{proof}

\begin{lemma}[Hash-token charge case]
\label{lem:hash-token-charge-case}
Every completed nonterminal fixed-window hash return contributes at most one
token of value $C_{\rm hash}$.  This token is assigned injectively to the first
terminal, fresh, old, scale, or birth hard exit supplied by the returned
residual block.  At that exit the corresponding case lemma above is applied
with $\varepsilon_{\rm hash}=1$, unless the exit is terminal, in which case the
token is included in the terminal local charge.  No child state is entered with
an unpaid token.
\end{lemma}

\begin{proof}
This is Proposition~\ref{prop:hash-token-matching} together with
Lemma~\ref{lem:token-free-children}.  Bounded-support hash exits do not create a
pending token; they are classified immediately as fresh or old by
Proposition~\ref{prop:hash-ledger}.  Hence hash charges are never a reusable
branch budget and cannot pile up at one hard exit.
\end{proof}

\begin{lemma}[Reservoir and conditioned-copy charge case]
\label{lem:reservoir-charge-case}
Auxiliary conditioned-copy labels and reservoir labels do not create a new
potential source.  The Renyi copy tax of a fixed-window hash return is included
in $C_{\rm hash}$, and after that tax is paid
Lemma~\ref{lem:reservoir-marginalization} erases auxiliary labels without
changing the law of later selectors.  Terminal reservoir labels are bounded by
the $R(\Theta)$ term.
\end{lemma}

\begin{proof}
The copy tax is part of Theorem~\ref{thm:public-hash} and the definition of
$C_{\rm hash}$.  Proposition~\ref{prop:reservoir} bounds all retained
nonterminal sample labels by the fixed reservoir size; terminal labels are
counted in the projected shadow, and auxiliary labels are not retained after
marginalization.
\end{proof}

\begin{proposition}[Global charge summation]
\label{prop:global-charge}
For the finite truncation, suppose every reached nonterminal selector charge
$c_v$ is decomposed as
\[
  c_v=c_v^{\rm sh}+c_v^{\rm fr}+c_v^{\rm old}
      +c_v^{\rm sc}+c_v^{\rm hash},
\]
with the following ledger certificates on each retained branch:
\begin{enumerate}[label=\textup{(\roman*)}]
\item terminal-shadow charges satisfy
  \[
    \sum_v c_v^{\rm sh}\le C_{\rm sh}(T(\Theta)+R(\Theta));
  \]
\item fresh charges are attached to disjoint public sets newly promoted from
  $[k]\setminus U$ into $U$;
\item old charges are attached to old-coordinate incidences debited in
  $N_{\rm old}$, and no coordinate is used after its revisit counter reaches
  the threshold $Q$ unless the branch exits through the high-revisit rule;
\item scale charges at an active scale $m$ occur only on a child whose next
  scale is at most $\theta m$ for some fixed $\theta<1$, and every new active
  block $J$ receives a paid rank/scale certificate from the parent active-rank
  drop, the parent active-scale account, fresh progress, or old-incidence
  account which created it.  In the fresh/old-source cases this certificate
  includes $(B_{\rm rk}+B_s)|J|$.  Terminal shadows and DRC selector
  certificates are not recurrence-potential sources;
\item every nonterminal hash-window charge is either terminal, or is carried
as one pending token of value at most $C_{\rm hash}$ and assigned to the first
subsequent fresh promotion, old incidence, scale descent, or close-pair birth
produced by the hash return; bounded hash-support exits are classified
directly into fresh promotion or old-incidence registration by
Proposition~\ref{prop:hash-ledger}.  This assignment is made before any child
state or descendant hash window is entered from the returned residual block.
\end{enumerate}
Then, after increasing the constants in the potential hierarchy,
\[
  \sum_{v\ {\rm reached}} c_v
  \le
  C\,(T(\Theta)+R(\Theta)+\Delta\Phi)
\]
on every retained branch, and hence the same inequality holds in expectation.
\end{proposition}

\begin{proof}
The proof is the sum of the explicit case lemmas just proved.  Terminal and
pruning charges are Lemma~\ref{lem:terminal-pruning-charge-case}.  Fresh exits,
old-incidence exits, parent-rank DRC/split births, and deletion-trapped scale
steps are Lemmas~\ref{lem:fresh-charge-case}, \ref{lem:old-charge-case},
\ref{lem:rank-birth-charge-case}, and \ref{lem:scale-charge-case}.  Hash tokens
are assigned by Lemma~\ref{lem:hash-token-charge-case}, and reservoir/copy
labels by Lemma~\ref{lem:reservoir-charge-case}.  The remaining paragraphs only
record how the branchwise sums telescope.

The terminal-shadow part is (i).  The numerical inequalities used in the fresh,
old, parent-rank, and scale paragraphs are exactly the corresponding case
lemmas; the assertion that a hard exit pays at most one fixed hash token is
Lemma~\ref{lem:hash-token-charge-case}.
For the fresh part, if disjoint public sets
$F_a\subseteq [k]\setminus U$ are promoted to $U$, then
\[
  \Delta(\Phi_{\rm fresh}+\Phi_{\rm old})
  =
  \sum_a (B_f-B_o)|F_a|
\]
before old incidences are counted.  The definition of $C_{\rm fr}(q)$ includes
the bounded fresh-pattern and promotion transcript per promoted coordinate.
If a fresh promotion creates a continuing active child $J_a\subseteq F_a$,
the account first transfers $(B_{\rm rk}+B_s)|J_a|$ into that child's
rank and scale balances.  The
remaining spendable fresh drop is at least
\[
  (B_f-B_o)|F_a|-(B_{\rm rk}+B_s)|J_a|,
\]
which, by $B_f-B_o>C_{\rm fr}(q)+C_{\rm hash}+B_{\rm rk}+B_s+10$, pays the finite
fresh transcript and any single bounded hash charge assigned to that nonempty
fresh exit.  The transferred rank and scale balance remains in the child
potential and is not counted again as spent expenditure.

For old data, each charged role, pattern, bounded-support use, or old-window
auxiliary transcript contributes at most $C_{\rm old}(Q,r_*)$ per
old-coordinate incidence.  If an old exit registers $n$ raw incidences and
creates a continuing active child $J$, then $|J|\le n$.  The old account debits
$\omega_{\rm old}n$ from $N_{\rm old}$, transfers
$(B_{\rm rk}+B_s)|J|$ into the child's rank and
scale balance, and leaves at least
\[
  \omega_{\rm old}n-(B_{\rm rk}+B_s)|J|
  \ge C_{\rm old}(Q,r_*)n+C_{\rm hash}+10
\]
of spendable old expenditure whenever a hash token is assigned to that
nonempty exit.  This pays $\sum_v c_v^{\rm old}$ and the assigned token, while
the transferred rank and scale balance remains in the child potential.  On every
interval on which $U$ is fixed, the debit $\Delta N_{\rm old}$ is the raw drop
of $\Phi_{\rm old}$ before such transfers.  When $U$ increases, the increase
$B_o|F|$ in
$\Phi_{\rm old}$ is already included in the net fresh-promotion drop
$(B_f-B_o)|F|$ above.
The cap $N_{\rm old}\le Q\omega_{\rm old}|U|<B_o|U|$ keeps $\Phi_{\rm old}$
nonnegative on nonterminal low-revisit branches.

For scale data, a deletion-trapped transition from $m$ to at most
$\theta m$ has drop at least $B_s(1-\theta)m$ in the active scale potential.
Thus all charges $O(m)$ after a scale has been born are bounded by a geometric
series times the deposited active scale budget.  By (iv), births are not free:
before a new active block $J$ is entered, its inherited rank and scale balances
are assigned to the ledger which created it.  Proposition~\ref{prop:birth}
verifies the source in each case: the initial block is paid by the root
active-rank budget; DRC/split children of an existing active rank are paid by
the corresponding rank drop or by that parent's deposited scale balance; fresh
births by the fresh drop; old births by registered old incidences.  A terminal
shadow stops the branch and creates no nonterminal active child.  A DRC atom
and role refinement pay selector entropy, but not recurrence potential; the
scale deposit for a DRC-produced block is taken from the parent active-rank
drop, parent active-scale balance, or from the fresh/old progress which routed
to it.  Fixed hash returns
have size below $m_0$, and any later larger active child must receive its
deposit from the hard exit to which the hash token is assigned.  After
increasing the corresponding ledger constants, the birth budget is already
included in the preceding fresh, old, parent-rank, or parent-scale paragraph;
the bounded hash paragraph pays only the fixed hash token before it is assigned
to one of those hard-exit accounts.
Later scale transitions only spend that already deposited amount, so a positive
increase of $\Phi_{\rm scale}$ is never treated as a free negative drop.

Finally, (v) says a nonterminal hash window is not a separate reusable budget.
It is carried as a single pending charge of the returned block until the first
hard exit supplied by that block, or it is a bounded-support exit classified
as fresh or old by Proposition~\ref{prop:hash-ledger}.  When the first hard
exit occurs, the token value is included in that exit's local charge $c_v$ and
is removed from the returned block's state before any child state or descendant
hash window is entered.  Lemma~\ref{lem:token-free-children} gives the exact
state invariant: deletion-trapped scale children and close-pair birth children
start with no inherited hash debt.  This is a disjoint assignment of already
incurred entropy charges, not an operation which deletes entropy.  Since each
fixed hash window has cost at most $C_{\rm hash}$ and the assigned hard exits
are charged in the terminal-shadow, fresh, old, scale, or birth paragraphs,
adding $C_{\rm hash}$ to the corresponding ledger constant pays the hash part.
Summing the five paragraphs gives the branchwise inequality.
\end{proof}

\begin{proposition}[Terminal cylinder potential compatibility]
\label{prop:terminal-cylinder-potential}
Let a retained branch start from a parent public state $\mathsf S_{\rm parent}$
and terminate with projected shadow
\[
  \theta=((J_s,a_s))_{s\in S}.
\]
After all pending hash tokens, delayed old queues, and active-scale birth
deposits on that branch have been assigned, let $\Delta\Phi$ be the charged
potential expenditure of the branch.  For each $s\in S$, form the cylinder
child $\mathsf S_{\theta,s}$ whose underlying code is the full parent value
cylinder $\{x:x_{J_s}=a_s\}$ with $J_s$ deleted.  Its ledger/support data are
obtained from the terminal public state on the branch by deleting $J_s$ from the
current ambient coordinate set, active coordinates, old support, queues, and
incidence counters; non-value terminal atoms are not additional code
restrictions.  Then
\[
  \Delta\Phi+\sum_{s\in S}\Phi(\mathsf S_{\theta,s})
  \le
  R(\theta)\Phi(\mathsf S_{\rm parent}).
\]
\end{proposition}

\begin{proof}
The potential account used in Proposition~\ref{prop:global-charge} is a spent
ledger.  Hash tokens are assigned before a later hash window can open, delayed
old queues are either registered or charged to the hard exit which clears them,
and Proposition~\ref{prop:birth} assigns every scale birth to the ledger which
created it before the born block is entered.  After these assignments, the
branch has a spent amount $\Delta\Phi$ and a remaining terminal public ledger
$\Phi_{\rm leaf}$ with
\[
  \Delta\Phi+\Phi_{\rm leaf}\le \Phi(\mathsf S_{\rm parent}).
\]
The following table is a summary of the transition-by-transition check of this
account invariant.  The arithmetic entries are the inequalities of
Lemma~\ref{lem:step-account-inequalities}, with the hash-token multiplicity
controlled by Proposition~\ref{prop:hash-token-matching}.  ``Balance'' means
the remaining public potential available to cylinder children.
\begin{center}
\scriptsize
\begin{tabular}{p{0.22\linewidth}p{0.34\linewidth}p{0.32\linewidth}}
transition source & account action & compatibility effect\\
\hline
public deterministic choice
& no non-public selector and no potential transfer
& $\Delta\Phi$ and the balance are unchanged\\
terminal value/shadow pin
& charged to $T(\Theta)+R(\Theta)$, not to potential
& pinned coordinates are deleted in cylinder children\\
product-KL or collision exit
& terminal pins go to shadow; low-collision outputs enter a labelled hash token
& no active scale or old queue is born without a later assigned token\\
hash bounded support
& fresh part is promoted, or old incidences are registered
& the window closes and cannot be charged again\\
hash near-sunflower return
& $C_{\rm hash}$ is carried as one token and assigned to the first hard exit
& the token is paid before any child state or descendant hash window is entered\\
fresh promotion
& transfers any fresh-born $(B_{\rm rk}+B_s)|J|$ rank/scale balance, then spends the remaining
drop from $\Phi_{\rm fresh}+\Phi_{\rm old}$
& promoted coordinates are disjoint, and transferred rank/scale balance remains
in the child\\
old residual/window
& debits $N_{\rm old}$, transfers any old-born rank/scale balance, then spends the
remaining debit; or exits by high revisit
& delayed queues are registered or assigned before terminal counting\\
scale birth
& rank/scale balance is certified from parent rank drop, parent scale, fresh
drop, or old incidences
& terminal shadows and DRC selector certificates do not create recurrence
balance\\
scale descent
& spends the deposited scale balance geometrically
& the child has smaller active scale and no inherited hash debt\\
DRC/split block
& the public DRC atom and role refinement pay selector entropy
& any rank/scale balance comes from parent rank/scale or fresh/old progress\\
reservoir/copy/pruning
& copy tax or pruning mass is paid locally
& auxiliary labels and pruned branches are not inherited by cylinder children\\
bounded leaf
& absorbed by $D_0$ in the finite base case
& no recursive cylinder child with unpaid balance is created
\end{tabular}
\end{center}
Formally, the account invariant is imposed at token-free checkpoints.  A
nonterminal hash return may create one bounded pending token, but this is a
temporary one-step debt attached to the returned residual block, not a new
recurrence balance.  By Lemma~\ref{lem:token-free-children}, the token must be
assigned to the exit which leaves that residual block before any child state or
descendant hash window is entered.  Let $S_t$ be the accumulated spent potential
after the first $t$ checkpoint-to-checkpoint transitions of the branch, and let
$B_t$ be the public recurrence balance still carried by the checkpoint state,
including inherited rank/scale balances but excluding selector entropy already
charged to $T(\Theta)+R(\Theta)$ or to the local proof-tree ledger.  Initially
\[
  S_0=0,\qquad B_0=\Phi(\mathsf S_{\rm parent}).
\]
For each nonterminal checkpoint transition, Lemmas~\ref{lem:step-account-inequalities}
and \ref{lem:token-free-children}
give
\[
  S_{t+1}-S_t+B_{t+1}\le B_t
\]
with any unique possible hash token from
Proposition~\ref{prop:hash-token-matching} included in the hard-exit charge
which restores the next checkpoint.  During the intermediate residual state
immediately after a hash return, the ordinary checkpoint inequality is not
asserted; equivalently one may view the branch as carrying a bounded overdraft
of size at most $C_{\rm hash}$.  That overdraft is cleared at the next checkpoint
and cannot be inherited by a child.  Public deterministic choices have equality;
terminal value pins and projected shadows add no recurrence child; copying and
pruning charges are local transcript or discarded-mass charges and do not
increase $B_{t+1}$; bounded leaves stop the recurrence and are absorbed by
$D_0$.  Induction on token-free checkpoints yields the account invariant
\[
  \Delta\Phi+\Phi_{\rm leaf}\le \Phi(\mathsf S_{\rm parent}).
\]
The child code itself is the full value cylinder, so the non-value terminal
transcript does not affect the cylinder-size estimate in
Theorem~\ref{thm:shadow-potential}.  For the inherited ledger, deleting pinned
coordinates from the terminal public state cannot increase the rank, fresh,
old-incidence, or scale components of $\Phi$.  The rank and scale terms decrease
because the child ambient and active blocks are smaller.  A
deleted coordinate outside $U$ is removed from the current ambient
$[k]\setminus U$, so the fresh term decreases by $B_f$.  The only nontrivial
case is the old term.  If a deleted old coordinate $i$ has raw revisit counter
$m_i\le Q$, deleting it changes
\[
  \Phi_{\rm old}=B_o|U|-\omega_{\rm old}I_{\rm old}
\]
by
\[
  -B_o+\omega_{\rm old}m_i\le -B_o+Q\omega_{\rm old}<0.
\]
Deleting queued incidences is no worse: queued incidences which have not yet
been debited were either assigned to the hard exit before terminal counting, or
else are absent from both the terminal child and $I_{\rm old}$.  Thus each
inherited cylinder child satisfies
\[
  \Phi(\mathsf S_{\theta,s})\le \Phi_{\rm leaf}.
\]
Therefore
\[
  \Delta\Phi+\sum_{s\in S}\Phi(\mathsf S_{\theta,s})
  \le
  \Delta\Phi+R(\theta)\Phi_{\rm leaf}
  \le
  R(\theta)\Phi(\mathsf S_{\rm parent}),
\]
because $\Delta\Phi+\Phi_{\rm leaf}\le\Phi(\mathsf S_{\rm parent})$ by the
account invariant, $R(\theta)\ge1$, and
$\Phi_{\rm leaf}\le\Phi(\mathsf S_{\rm parent})$.
\end{proof}

\begin{lemma}[Selector entropy primitives]
\label{lem:selector-primitives}
All selector estimates used in Proposition~\ref{prop:local-selector-certificates}
reduce to the following elementary bounds.
\begin{enumerate}[label=\textup{(\roman*)}]
\item Let $p$ be a finite law and $0<\theta\le1$.  If
$H_\theta=\{a:p(a)\ge\theta\}$, then a selector which chooses an atom of
$H_\theta$ or the complement has entropy at most
\[
  \log(1+\lfloor\theta^{-1}\rfloor).
\]
On an atom branch the value has conditional mass at least $\theta$ and hence
is a fixed-cost projected tuple pin.
\item Let $0<\rho,\tau\le1$ and put
\[
  H=\{a:p(a)\ge\rho\tau\},\qquad L=H^c .
\]
The atom branches in $H$ are fixed-cost terminal pins.  If $p(L)<\rho$, then
discarding $L$ has pruning loss $-\log(1-p(L))\le2\rho$ for
$\rho\le1/2$.  If $p(L)\ge\rho$, then
\[
  \Coll(p(\cdot\mid L))
  =
  \sum_{a\in L}\left(\frac{p(a)}{p(L)}\right)^2
  \le \frac{\max_{a\in L}p(a)}{p(L)}
  \le \tau .
\]
\item If a stopped prefix selector outputs a word $u$ with length $\ell(u)$
and every retained output satisfies
\[
  -\log P[U=u]\le \lambda\ell(u)+b,
\]
then
\[
  \bfH(U)\le \lambda\,\EE\ell(U)+b.
\]
For tuple-coordinate prefixes, deterministic repeats have zero conditional
entropy and inconsistent repeats have zero mass; after discarding those
repeats the number of entropy-bearing values is at most the projected shadow
coordinate count $T(\Theta)$.
\item A retained public set event of conditional mass at least $\rho$ adds at
most $\log(1/\rho)$ to a labelled hash transcript.  If appending that set
event would exceed the remaining $B_{\rm win}$ budget, the set event is not
appended and the branch terminalizes on the already paid value prefix.
\item At fixed density parameters, a public DRC atom on a block of size $s$
has cost $O(s)$, and a three-role refinement on that block has entropy at most
$s\log 3$.
\item For one labelled hash-window key $K$, if $\bfH(K)\le B_{\rm win}$ and
$G$ is the good-key event with conditional probability at least $1/2$, then
\[
  \bfH_3(K\mid G)\le \bfH(K\mid G)\le 2B_{\rm win}+O(1).
\]
Thus the three-copy Renyi tax is $2\bfH_3(K\mid G)+O(1)=O(B_{\rm win})$, and
the exact split exposure on a fixed $h$-block costs at most
\[
  \log\sum_{r\le6\beta h}\binom hr\le h\log 2.
\]
\item With $Q,r_*$, the old scale, and the sample-label signature fixed, the
continuing old-window selectors have a finite per-incidence alphabet once
high-symbol value vectors are not exposed.  Hence their total entropy is at
most $C_{\rm old}(Q,r_*)n$ on a branch registering $n$ old incidences.
\end{enumerate}
\end{lemma}

\begin{proof}
For (i), $|H_\theta|\le\theta^{-1}$.  Item (ii) is the usual heavy-or-light
decomposition; the pruning bound is $-\log(1-x)\le2x$ for $0\le x\le1/2$.
Item (iii) follows by averaging the displayed pointwise prefix bound.  The
projection statement is Lemma~\ref{lem:tuple-prefix-shadow}.  Item (iv) is the
definition of the stopped per-labelled-window transcript.  Item (v) is
Lemma~\ref{lem:drc-row-block} and the bound $|\{XY|Z,XZ|Y,YZ|X\}^s|=3^s$.
For (vi), Renyi entropy decreases with order and conditioning on
$G$ of mass at least $1/2$ gives
\[
  \bfH(K\mid G)\le \bfH(K)/\Prb(G)+\log2\le2B_{\rm win}+O(1);
\]
Lemma~\ref{lem:random-copy} gives the factor $2\bfH_3$ for three copies.
Item (vii) is the finite-alphabet accounting in Theorem~\ref{hyp:old-window}:
the continuing data consist only of sample-label signature, equality pattern,
role, scale, bounded reservoir labels, and queue-flush markers, all charged per
registered incidence before the revisit cap.
\end{proof}

\begin{proposition}[Local selector entropy certificates]
\label{prop:local-selector-certificates}
Run the finite truncation of the normal-form transition system.  At a reached
nonterminal node $v$, let $E_v$ be the event that the branch reaches $v$ and
let $\Gamma_v$ be the first non-public child selector after all deterministic
public tie-breaks at $v$ have been applied.  Let $\ell_v^{\rm pr}$ be the
pruning-normalization loss at $v$, equal to $-\log q_v$ on a retained pruning
step of retained conditional mass $q_v$ and equal to $0$ otherwise.  There are nonnegative local
charges
\[
  c_v=c_v^{\rm sh}+c_v^{\rm fr}+c_v^{\rm old}
      +c_v^{\rm sc}+c_v^{\rm hash}
\]
such that
\[
  \bfH(\Gamma_v\mid E_v)+\EE[\ell_v^{\rm pr}\mid E_v]\le \EE[c_v\mid E_v]
\]
for every reached node, and on every retained branch these charges satisfy the
five ledger certificates of Proposition~\ref{prop:global-charge}.
\end{proposition}

\begin{proof}
The elementary entropy estimates used in the rows are collected in
Lemma~\ref{lem:selector-primitives}.  The following table gives the selector
and the local entropy bound.  All
coordinate identities, block orders, and tie-breaks not shown in the second
column are $\mathcal F_v$-measurable and therefore have zero conditional
entropy.
\begin{center}
\scriptsize
\begin{tabular}{p{0.18\linewidth}p{0.25\linewidth}p{0.23\linewidth}p{0.16\linewidth}}
transition & $\Gamma_v$ and pruning loss & entropy/loss bound & ledger\\[1mm]
\hline
public choice &
deterministic block, scale, or tie-break &
0 &
none\\[1mm]
terminal atom/value pin &
sample labels and pinned coordinate values &
$C(T_v+R_v)$ &
shadow\\[1mm]
low common-symbol entropy &
$A_J=X_J=Y_J$ on the terminal block &
$\bfH(A_J\mid E_v)\le h|J|$ &
shadow\\[1mm]
collision exit &
heavy atoms, light complement, or one hash coordinate &
$O(1)$ plus possible pruning/hash transcript &
sh./pr./hash\\[1mm]
product-KL exit &
interleaved prefix, tail/prune bit, or value atom from $(A_j,B_j)$ &
$\lambda|\text{prefix}|+O(1)$, or terminal value cost &
shadow/hash\\[1mm]
hash window &
public sample label, transcript, split exposure, and copy key &
$B_{\rm win}+2\bfH_3(K)+O(h)$ &
hash, then assigned\\[1mm]
fresh residual &
fresh equality pattern/support promotion &
$C_{\rm fr}(q)|F|$ &
fresh potential\\[1mm]
old residual/window &
sample-label signature, five-state pattern, role vector, incidence data; no
high-symbol values &
$\le \Delta N_{\rm old}$ after the fixed old debit is applied &
old potential\\[1mm]
low-symbol old terminal &
old common-value vector on $O'$ &
$\bfH(A_{O'}\mid E_v)\le h|O'|$ &
shadow\\[1mm]
scale split/descent &
DRC atom, exact split set, and role refinement on an active block &
$O(|J|)+C_\sigma |J|+|J|\log 3$ &
rank/scale birth/drop\\[1mm]
reservoir/copy labels &
bounded active sample labels; auxiliary copies not retained in terminal data &
$O(R_v)$ plus paid copy tax &
reservoir/hash
\end{tabular}
\end{center}

We verify the rows in the order in which they can occur.  Public rows have
zero charge by the definition of the public sigma-algebra.  Terminal atoms,
terminal tuple pins, immediate value pins, and low common-symbol exits are
terminal projected shadows.  The cross/high-collision exits of
Lemma~\ref{lem:collision-exit-pin} are heavy/light routings: their heavy
branches are terminal projected shadows, their small complements are charged
pruning branches, and their large complements are actual low-collision hash
coordinates.  The terminal value entropy is exactly the entropy of the pinned
coordinates and the sample labels retained in the shadow.  A one-coordinate
heavy pin has fixed selector entropy by Lemma~\ref{lem:selector-primitives}(i),
while a terminalized prefix is paid by the amortized prefix estimate
Lemma~\ref{lem:selector-primitives}(iii).  Hence terminal value data is
charged to $c_v^{\rm sh}$ and contributes to $T(\Theta)+R(\Theta)$.

The product-KL row is Theorem~\ref{thm:product-kl-branch}.  The prefix
value is either amortized, with cost at most $\lambda|\text{prefix}|$, or the
branch is routed through the first-crossing heavy/light scan using the previous
amortized prefix.  If no positive-KL or first-crossing stop occurs, the full
tuple prefix is amortized and terminal.  Positive-tail and first-crossing
	light-complement events have fixed cost when retained in the hash transcript;
	if they would overflow the window they are not appended.  Any
	low-collision output is the actual conditional law of one tuple coordinate
	under the paid transcript.  Thus the row is charged to terminal shadow,
	pruning, or to the hash ledger, never to an unpaid pair-alphabet selector.
	At a retained pruning step the child selector itself is deterministic after
	the public set is known, but the normalization loss
	$\ell_v^{\rm pr}=-\log q_v$ is nonzero; the estimates
	in Lemma~\ref{lem:selector-primitives}(ii)--(iv), equivalently the
	corresponding estimates in Lemmas~\ref{lem:positive-tail} and
	\ref{lem:heavy-or-light-hash}, put that loss into the same local charge.

The hash row is Theorem~\ref{thm:public-hash} together with
Propositions~\ref{prop:hash-ledger} and \ref{prop:hash-accumulation}.  The
sample label is public, and the
selected coordinate block is a deterministic function of the paid window
transcript.  The local charge consists of the window transcript entropy, the
order-three Renyi copying tax for that label's conditional law, and the exact
split exposure, bounded by Lemma~\ref{lem:selector-primitives}(vi).  If the
block is returned to residual mode, this charge is
carried as one pending token and assigned to the first subsequent terminal,
fresh, old, scale, or birth exit before another hash window can be opened from
that returned block or from a child born inside it; a bounded leaf or terminal
shadow pays it directly.

Fresh residual charges are supplied by
Lemma~\ref{lem:residual-support-promotion}.  The promoted fresh coordinates
form a public set $F$ disjoint from all earlier fresh promotions, and the
finite pattern/role transcript costs at most $C_{\rm fr}(q)|F|$.  This is
$c_v^{\rm fr}$.

Old residual and old-window charges are supplied by
Lemma~\ref{lem:old-residual-routing}, Proposition~\ref{prop:old-return-public},
Theorem~\ref{thm:old-symbol-budget}, Theorem~\ref{hyp:old-window}, and
Proposition~\ref{prop:delayed-old-queue}.  Low-symbol old values terminalize
and are charged to the shadow row.  High-symbol old values are not selectors
of the continuing old-window tree.  The remaining sample-label signature,
equality-pattern, role, scale, reservoir-label, and incidence data are bounded
by $C_{\rm old}(Q,r_*)$ per registered old incidence.  The weighted debit
$\omega_{\rm old}$ also reserves any old-born rank and scale balances and any
single hash token assigned to the nonempty old exit.  The delayed queue prevents
the same bounded short window for a fixed signature from being repeatedly
recharged without increasing the raw counter and hence $N_{\rm old}$; if the
signature changes, the bounded queue is charged to that hard exit before it is
cleared.
When a high-revisit branch is routed by Lemma~\ref{lem:high-revisit} to a
one-coordinate hash output rather than to terminal value pins or a charged
pruned complement, the selector has fixed cost $O(1)$ by
Lemma~\ref{lem:selector-primitives}(ii).  This cost is charged to the same
old-incidence/high-revisit account which detects the cap; if the coordinate
later participates in a hash return, the window transcript and copy cost are
separately paid by the single hash token assigned to the next hard exit.  Thus
queue increments before a completed window are not hidden inside the final
token, and the complement of a heavy old value is never silently discarded.

Scale rows come from the DRC and split-exposure steps in residual repair and
from deletion-trapped descent.  Lemma~\ref{lem:drc-row-block} gives a DRC atom
of cost $O(|J|)$ because the density parameters are fixed.  The exact sparse
split exposure costs at most $C_\sigma |J|$, and role refinement costs
$|J|\log 3$.  Lemma~\ref{lem:deletion-scale-descent} gives the geometric
scale drop, while Proposition~\ref{prop:birth} supplies the birth certificate
for every new active block.  When the block is born by replacing an active
rank-$m$ state by a fixed-fraction DRC or split child, the certificate is drawn
from the active-rank drop before the child rank potential is inherited.

Finally, reservoir and conditioned-copy labels are controlled by
Proposition~\ref{prop:reservoir} and Lemmas~\ref{lem:fixed-copy} and
\ref{lem:random-copy}.  Auxiliary copy labels are paid by their copy tax and
are not retained in the terminal transcript.  The public fresh-reset pruning
loss of Lemma~\ref{lem:fresh-reset-pruning} is included in the auxiliary
constant $C_{\rm aux}$ and hence in the hard-exit charge which clears that
window.  This proves the displayed
conditional entropy inequalities and also verifies the five branchwise ledger
certificates required by Proposition~\ref{prop:global-charge}.
\end{proof}

\begin{theorem}[Delayed-window shadow-potential assembly]
\label{thm:assembly}
Run the finite truncation of the transition system from
Definition~\ref{def:transition-system} and Proposition~\ref{prop:finite-truncation},
enforce Definition~\ref{hyp:public-drc}, use Theorem~\ref{hyp:old-window},
use Definition~\ref{hyp:reservoir}, use the product-KL exit of
Theorem~\ref{thm:product-kl-branch}, and choose the constants as in
Section~\ref{sec:constant-hierarchy}.  Assign every new active scale the birth
certificate specified in Definition~\ref{hyp:birth}; for the finite-state
system this assignment is verified in Proposition~\ref{prop:birth}.  The
branch invariants are those of Proposition~\ref{prop:branch-invariants}.
Then
the complete proof tree has pruning loss
\[
  L_{\rm pr}
  =
  \sum_{\textup{pruning steps on the branch}} -\log(\textup{retained mass}),
\]
with pruning steps interpreted as in
Proposition~\ref{prop:pruning-normalization}, and satisfies
\[
  \bfH(\Theta)+\EE L_{\rm pr}
  \le
  C\,\EE[T(\Theta)+R(\Theta)]
  + C\,\EE[\Delta\Phi].
\]
Consequently Theorem~\ref{thm:shadow-potential} and
Theorem~\ref{thm:multi-shadow} give a constant-base recurrence.
\end{theorem}

\begin{proof}
Let $\Gamma$ be the full terminal leaf transcript of the finite truncation:
it records every non-public child selector, every exact-exposure atom, every
role-refinement atom, and the terminal projected shadow before unpaid
analysis-only data are projected away.  Proposition~\ref{prop:tree-entropy-ledger}
gives the retained transcript entropy, and
Proposition~\ref{prop:pruning-normalization} adds exactly the accumulated
normalization loss of charged pruning steps.  Thus
\[
  \bfH(\Theta)+\EE L_{\rm pr}
  \le
  \sum_v \Prb(E_v)\,
  \bigl(\bfH(\Gamma_v\mid E_v)+\EE[\ell_v^{\rm pr}\mid E_v]\bigr).
\]
By Proposition~\ref{prop:local-selector-certificates}, the local charges
satisfy
\[
  \bfH(\Gamma_v\mid E_v)+\EE[\ell_v^{\rm pr}\mid E_v]\le \EE[c_v\mid E_v],
\]
and their branchwise decomposition satisfies the five hypotheses of
Proposition~\ref{prop:global-charge}.  Therefore
Proposition~\ref{prop:global-charge}, followed by expectation, gives
\[
  \bfH(\Theta)
  \le
  C\,\EE[T(\Theta)+R(\Theta)]
  + C\,\EE[\Delta\Phi]
\]
for a constant $C$ depending only on the fixed thresholds.  This estimate is
obtained before any use of shadow descent.

It remains to check that the potential in Theorem~\ref{thm:shadow-potential}
is the recurrence potential of the cylinder children, not the potential of an
intermediate transition child.  This is exactly
Proposition~\ref{prop:terminal-cylinder-potential}: the cylinder children are
formed from the terminal public state by deleting the pinned shadow
coordinates, and their inherited potentials satisfy the
recurrence-compatibility hypothesis in Theorem~\ref{thm:shadow-potential}.

Now Theorem~\ref{thm:shadow-potential} converts the entropy estimate to a
pointwise shadow atom and combines it with the stateful version of
Theorem~\ref{thm:multi-shadow}.  The $T(\Theta),R(\Theta)$ terms are absorbed
after $D$ is chosen in Section~\ref{sec:constant-hierarchy}.
\end{proof}

\section{Active Scale Birth}

It remains to prevent free increases of $\Phi_{\rm scale}$.

\begin{definition}[Rank/scale birth certificate]
\label{hyp:birth}
Fix the rank and scale coefficients $B_{\rm rk},B_s$.  Every new active block
$J$ is assigned a public certificate for the rank and scale balances it carries.
If $J$ is a fixed-fraction child of a parent active rank $m$, the certificate
keeps the inherited $B_{\rm rk}|J|$ rank balance in the child and pays the
selector and $B_s|J|$ scale balance from the parent rank drop or parent
active-scale account.  If $J$ is born from a fresh or old hard exit, the
certificate transfers
\[
  (B_{\rm rk}+B_s)|J|
\]
from the corresponding fresh/rank or old-incidence account into the child.
The only allowed public sources are: initial active-rank, parent active-rank
drop, parent active-scale, fresh/rank progress, or old-incidence.
Terminal shadows are terminal and do not create
nonterminal active children.  Public DRC certificates pay selector entropy but
are not independent sources of recurrence potential.  Fixed hash-window
returns are below the bounded-leaf scale $m_0$ unless their pending token is
assigned to one of the listed hard-exit accounts before any larger active child
is entered; by Lemma~\ref{lem:token-free-children} the child itself is
token-free.
\end{definition}

\begin{proposition}[Scale birth audit]
\label{prop:birth}
Under the public DRC convention, immediate terminalization, public hash
returns, and fixed-scale old-window reduction, Definition~\ref{hyp:birth}
holds for all active blocks produced by the finite-state system.
\end{proposition}

\begin{proof}
Initial blocks are paid by the active-rank budget of the root state.  A
terminal shadow stops the branch, so it never pays for a continuing active
scale.  Residual DRC blocks and close-pair activation blocks produced from
Lemma~\ref{lem:drc-row-block} have public selector cost $O(|J|)$ because all
density parameters are fixed constants, but that selector cost is not a scale
deposit.  If such a block is a fixed-fraction child of a parent active rank
$m$, the constant hierarchy makes the rank drop
$B_{\rm rk}(m-|J|)$ large enough to reserve $B_s|J|$ and pay the selector cost
while the $B_{\rm rk}|J|$ rank balance remains inherited in the child.  If the
block is instead produced by a fresh or old escape, the full
$(B_{\rm rk}+B_s)|J|$ rank/scale balance is drawn from the corresponding fresh
drop or old-incidence account.  Near-sunflower activation split exposures are paid by
the same active-rank certificate.
Fresh residual blocks are paid by fresh/rank progress; in the bounded
$|F|\le q$ case, Lemma~\ref{lem:residual-support-promotion} first promotes the
nonempty split part $F_{\rm sp}$ before the fresh close-pair child is entered.
Old residual blocks are paid by weighted old incidence, including the old-supported
residual witness routing of Lemma~\ref{lem:old-residual-routing}.  Pooled close-pair blocks are paid by the
incidence window which produced them.  High-revisit bounded-support blocks are
paid by the repeated old incidences that triggered the high-revisit exit.
One-coordinate KL increments are not active scale chains.  A fixed hash return
has size $h<m_0$ and is therefore bounded at the recurrence interface unless a
later hard exit creates a larger active child; that later child is funded by the
hard-exit account to which the hash token has been assigned, and enters
token-free by Lemma~\ref{lem:token-free-children}.  If a fixed hash
block is ever treated as a child of a parent active rank $m\ge m_0$, the choice
$m_0>h/\theta_{\rm rk}$ gives $h\le\theta_{\rm rk}m$, so the parent-rank
inequality of Lemma~\ref{lem:step-account-inequalities} applies as a separate
fallback.  This exhausts the block sources in the finite-state normal form
without using terminal shadows or DRC selector certificates as independent
recurrence-potential sources.
\end{proof}

\section{Constant Hierarchy}
\label{sec:constant-hierarchy}

The constants are chosen in the following order.
\begin{enumerate}[label=\textup{(\arabic*)}]
\item Choose a cross-collision threshold $0<\gamma_0<1$, then choose the
activation collision threshold $\alpha$ and the residual density/hash survival
threshold $0<\beta<1/48$ small enough for the $\sigma$ below to satisfy
$4\sigma<1$.  Choose a product-transfer tolerance
$\epsilon>0$ with $1/2-3\sqrt{\epsilon/2}>0$.  Put
\[
  \sigma
  :=
  \max\left\{6\alpha,\ \frac{6\alpha}{(1-\gamma_0)^2},\ 6\beta\right\}
\]
for Lemma~\ref{lem:deletion-scale-descent}.  Fix the scale-budget coefficient
$B_s$ in
Definition~\ref{hyp:birth} large enough to pay the exact split exposure cost
$\log\sum_{r\le\sigma m}\binom mr$ and the associated role-refinement cost on
each active block of size $m$.  Then choose
$\tau\le\min(\alpha,\gamma_0^2,\beta)$ for the activation, collision-exit, and
hash survival scans.  Define
\[
  \theta_{\rm rk}:=\max\{\sigma,\ (1-\tau)\alpha/4,\ \beta/4,\ 4\sigma\}<1.
\]
Positive-KL value pins have the fixed mass threshold
$\rho\tau$, where $\rho=1/2-3\sqrt{\epsilon/2}$.  There is no
constant-hierarchy condition
involving the binary entropy of subsets of $L$: residual repair first promotes
all possible fresh witnesses globally, and the remaining old-supported
witnesses are paid by the old-incidence ledger rather than by exposing
$W_L$.
\item Choose the KL prefix threshold, an absolute fixed hash-window length
$h$, and a cumulative hash transcript budget $B_{\rm win}$ so that
Theorem~\ref{thm:positive-kl-public} feeds only bounded-prefix coordinates
into the fixed hash window while the cumulative budget remains available;
amortized prefixes exceeding the remaining transcript budget terminalize as
projected shadows.  Definition~\ref{def:hash-survival} terminalizes any later
atom of mass $\ge\beta$.
\item Choose the residual support threshold $q$, the old-window queue
threshold, the old revisit threshold $Q$, and reservoir size $r_*$.  These
determine the finite constants
\[
  C_{\rm fr}(q),\qquad C_{\rm hash},\qquad C_{\rm old}(Q,r_*),
\]
as follows.  $C_{\rm fr}(q)$ is chosen as a per-promoted-coordinate constant
large enough to pay the maximum fresh selector/transcript entropy over the
finite list of residual transitions exposing at most $q$ fresh coordinates.  In
particular, when $|F|\le q$ but only a nonempty split part $F_{\rm sp}$ is
promoted, the whole $q\log 5$ five-state pattern exposure and the bounded
close-pair routing produced from it are bounded by
$C_{\rm fr}(q)|F_{\rm sp}|$.  It excludes rank/scale deposits and hash tokens,
which are paid separately by the displayed inequalities for $B_f-B_o$.  The
hash constant is taken explicitly as
\[
  C_{\rm hash}
  =
  C_0\bigl(B_{\rm win}+h+1\bigr),
\]
where $C_0$ is the absolute constant covering the proof-tree transcript
entropy $B_{\rm win}$, the Renyi copy tax
$2\bfH_3(T_G\mid P_0)\le 4B_{\rm win}+O(1)$, the exact split exposure
$\log\sum_{r\le 6\beta h}\binom hr$ (or the cruder $h\log2$), and the bounded
survival-scan bookkeeping in Theorem~\ref{thm:public-hash} and
Proposition~\ref{prop:hash-ledger}.  $C_{\rm old}(Q,r_*)$ is the maximum
old-window auxiliary transcript cost per old incidence before the revisit cap
$Q$, including the sample-label signature, five-state pattern, three-role
refinement, bounded reservoir labels, and the bounded queue-flush charge when
the active signature changes.  It excludes old-born rank/scale deposits and a
possible assigned hash token, which are paid by the later choice of
$\omega_{\rm old}$.  All three are finite because $q,h,Q,r_*$ and the window
length are fixed constants.
\item Choose the bounded-leaf scale
$m_0>\max(q,h/\theta_{\rm rk},4/\beta,8/((1-\tau)\alpha))$ large enough that, for
every scale $m\ge m_0$, the geometric scale drop
$B_s(1-4\sigma)m$ also absorbs any fixed additive token assigned to that scale
transition, including $C_{\rm hash}$ and the finite one-step scale selector
constants.  Then choose a crude constant $D_0$ large enough to dominate all
states of active rank and block size below $m_0$.
\item Choose the public fresh-reset gate parameter $\eta$ larger than the
finite auxiliary transcript constant $C_{\rm aux}$ in
Lemma~\ref{lem:fresh-reset-pruning}.
\item Let $C_{\rm rk}$ be the finite per-coordinate selector constant for
root activation, DRC row-block selection, sparse split exposure, and role
refinement at the fixed thresholds above.  Choose the active-rank coefficient
$B_{\rm rk}$ so that
\[
  B_{\rm rk}(1-\theta_{\rm rk})
  >(B_s+C_{\rm rk}+C_{\rm hash}+10)\theta_{\rm rk}.
\]
Then choose the old-incidence capacity coefficient
  \[
    \omega_{\rm old}:=C_{\rm old}(Q,r_*)+B_{\rm rk}+B_s+C_{\rm hash}+10,
    \qquad
    B_o>Q\,\omega_{\rm old}+C_{\rm aux}+10,
  \]
  and then choose the fresh
coefficient $B_f$ so that
\[
  B_f-B_o > C_{\rm fr}(q)+C_{\rm hash}+B_{\rm rk}+B_s+10 .
\]
With this choice, every public residual support promotion in
Lemma~\ref{lem:residual-support-promotion} has positive net potential drop and
bounded hash-support promotion has a fresh certificate after reserving the
rank and scale balances of the active child born from that promotion.
\item Let $C_*$ be the maximum of the finitely many row constants in the local
certificate table of Proposition~\ref{prop:local-selector-certificates}, after the choices above have
fixed $q,h,Q,r_*,B_s,B_{\rm rk},B_o,B_f$ and the window thresholds.  Finally choose
$A\ge C_*$ and
\[
  D\ge \max\{D_0,\exp(2C_*)\}
\]
large enough to absorb the terminal shadow, reservoir, birth-certificate, and
potential-drop constants in Theorems~\ref{thm:multi-shadow} and
\ref{thm:shadow-potential}.
\end{enumerate}

\begin{proposition}[Feasibility of the constant hierarchy]
\label{prop:constant-feasibility}
The constants in Section~\ref{sec:constant-hierarchy} can be chosen so that
all displayed inequalities and all fixed-window requirements used earlier in
the proof hold simultaneously.
\end{proposition}

\begin{proof}
The displayed requirements reduce to the following finite list.

First choose any fixed $0<\gamma_0<1$, then choose $\alpha>0$ so small that
$24\alpha/(1-\gamma_0)^2<1/2$, and choose $\epsilon>0$ with
\[
  \rho:=1/2-3\sqrt{\epsilon/2}>0 .
\]
Next choose $0<\beta<1/48$ small enough that $24\beta<1/2$, and put
\[
  \sigma:=\max\left\{6\alpha,\frac{6\alpha}{(1-\gamma_0)^2},6\beta\right\}.
\]
Then $4\sigma<1$.  Since $\sigma<1/4$,
\[
  C_\sigma=\sup_{m\ge1}m^{-1}
  \log\sum_{r\le\sigma m}\binom mr
\]
is finite.  The scale-budget coefficient $B_s$ is chosen to exceed the fixed
multiple of $C_\sigma+\log 3$ used for split-set exposure and role
refinement.

Choose
\[
  0<\tau\le\min(\alpha,\gamma_0^2,\beta).
\]
Then
\[
  \theta_{\rm rk}:=\max\{\sigma,\ (1-\tau)\alpha/4,\ \beta/4,\ 4\sigma\}<1 .
\]
This single inequality is the value-pin condition used by activation,
collision exits, and hash survival; the positive-KL value-pin threshold is
the smaller fixed constant $\rho\tau$.  It also removes the former binary
entropy constraint on subsets of $L$, because residual repair no longer
exposes a random old-supported set $W_L\subseteq L$ as an unpaid sparse
selector.

Choose the KL prefix threshold and the cumulative transcript budget
$B_{\rm win}$ for one fixed labelled hash window, with $B_{\rm win}$ larger than the
fixed tail/pruning increments $\log(1/\rho)+O(1)$ and the corresponding
positive-tail and first-crossing light-complement pruning charges.  This is a
per-window, per-sample-label budget.  Multiple sample labels may open distinct
windows on one branch; their transcripts are not concatenated into one global
$B_{\rm win}$ account, and each completed or bounded window is paid separately
by the hash-token matching ledger.  The construction of
Theorem~\ref{thm:positive-kl-public} sends a coordinate into its labelled hash mode only
when its prefix and tail/pruning transcript fits in the remaining budget of
the current window.  If a set-valued tail/pruning increment would overflow the
window, it is not appended; the branch terminalizes on the already paid value
prefix, and any larger prefix is terminalized as a projected shadow.
If the budget closes before $2h$ distinct live coordinates are accumulated,
the branch either has already terminalized on the projected prefix or is a
bounded-support hash exit.  Hence this step imposes no upper bound on
$B_{\rm win}$ from later constants.

After $q,h,Q,r_*$ and the old-window queue threshold are fixed, the constants
$C_{\rm fr}(q)$, $C_{\rm hash}$, and $C_{\rm old}(Q,r_*)$ defined above are
finite.  Since $4\sigma<1$, choose
$m_0>\max(q,h/\theta_{\rm rk},4/\beta,8/((1-\tau)\alpha))$ so large that
$B_s(1-4\sigma)m_0$ dominates $C_{\rm hash}$ and the finite additive constants
of a single scale transition, and then choose $D_0$ for the finitely many
bounded active ranks and block sizes below $m_0$.  The fresh-reset parameter is
chosen so that
\[
  \eta>C_{\rm aux}.
\]
Let $C_{\rm rk}$ be the resulting finite per-coordinate constant for
active-rank DRC/split selectors.  Choose
\[
  B_{\rm rk}(1-\theta_{\rm rk})
  >(B_s+C_{\rm rk}+C_{\rm hash}+10)\theta_{\rm rk}.
\]
Then choose
\[
  \omega_{\rm old}:=C_{\rm old}(Q,r_*)+B_{\rm rk}+B_s+C_{\rm hash}+10,
  \qquad
  B_o>Q\,\omega_{\rm old}+C_{\rm aux}+10
\]
and
\[
  B_f>B_o+C_{\rm fr}(q)+C_{\rm hash}+B_{\rm rk}+B_s+10 .
\]
These two inequalities are exactly the potential requirements used in
Proposition~\ref{prop:global-charge}: $B_o>Q\omega_{\rm old}$ keeps
$\Phi_{\rm old}=B_o|U|-N_{\rm old}$ nonnegative before high revisit because
$N_{\rm old}=\omega_{\rm old}I_{\rm old}$ and $I_{\rm old}\le Q|U|$.  The
weighted debit $\omega_{\rm old}$ pays old-incidence transcript, reserves any
old-born rank and scale balances, and pays any single hash token assigned to that
nonempty old exit, while
$B_f-B_o>C_{\rm fr}(q)+C_{\rm hash}+B_{\rm rk}+B_s+10$ pays every fresh
promotion together with any bounded hash charge after reserving the fresh-born
rank and scale balances assigned to it.

At this stage every local conditional-entropy estimate in
Proposition~\ref{prop:local-selector-certificates}
has a fixed finite coefficient.  More explicitly, after the threshold choices
above one may take
\[
\begin{aligned}
  C_*
  = C_{\rm abs}\max\{&
  C_{\rm sh},\ C_{\rm aux},\ C_\sigma+\log 3,\ B_s,\ 
  C_{\rm rk},\ C_{\rm fr}(q),\ C_{\rm hash},\ C_{\rm old}(Q,r_*),\\
  &B_{\rm rk},\ \omega_{\rm old},\ B_o,\ B_f,\ r_*,\ 1\},
\end{aligned}
\]
where $C_{\rm abs}$ is an absolute numerical constant and $C_{\rm sh}$ denotes
the terminal projected-shadow constant from
Theorem~\ref{thm:multi-shadow}.  This displayed maximum is not intended to
optimize constants; it records the dependency graph.  In particular,
\begin{align*}
  C_{\rm hash}&=C_0(B_{\rm win}+h+1),\\
  B_{\rm rk}&=O((B_s+C_{\rm rk}+C_{\rm hash}+1)/(1-\theta_{\rm rk})),\\
  B_o&=O(Q(C_{\rm old}+B_{\rm rk}+B_s+C_{\rm hash}+1)+C_{\rm aux}+1),\\
  B_f&=B_o+O(C_{\rm fr}+C_{\rm hash}+B_{\rm rk}+B_s+1).
\end{align*}
Thus $C_*$ may be very large, but it depends only on the fixed hierarchy and
not on $k$, the alphabet sizes, or $|\C|$.  It also does not depend on $A$ or
$D$.  Choose
\[
  A\ge C_*,
  \qquad
  D\ge \max\{D_0,\exp(2C_*)\}.
\]
Then the base cases below $m_0$ are dominated by $D^j$, and the shadow
potential theorem may be used with $\lambda=c=C_*$, since
\[
  \log D\ge 2C_*=\lambda+c .
\]
Every right-hand side above is already fixed at the moment it is used, and no
later displayed requirement asks for an upper bound on $B_o,B_f,A,D$ or a
lower bound on $\alpha,\epsilon,\tau$.
\end{proof}

\section{Interface Verification}

The public-normal-form reduction above uses explicit interfaces for root
activation, public DRC, residual support promotion, old-supported residual
routing, hash windows, projected shadows, active reservoirs, global charge
summation, and scale birth.  Theorem~\ref{thm:product-kl-branch}
routes the product-KL excess branch through a genuine sample coordinate, so
the verification never treats a low-collision law on the pair alphabet as a code
coordinate.

Corollary~\ref{cor:bounded-alphabet} is used only as a bounded-alphabet and
variable-profile check; it is not the alphabet-free mechanism.  Its bound
contains $Q$ because slice rank pays for exposed alphabet values.  In the
normal-form proof, high-alphabet old values are not exposed as selectors:
low-symbol value vectors terminalize locally, high-symbol old data is routed
through incidence and scale ledgers, and product-KL creates only a
one-coordinate pin or hash coordinate under the actual tuple law.  Thus the
alphabet size can enter only through terminal projected shadows or bounded
local transcripts already counted in $T(\Theta),R(\Theta)$, not as a factor in
the active-rank recurrence.

The fixed-window mechanism is therefore not a random restriction over the
alphabet.  A coordinate contributes to a hash key only if its whole value-prefix
and tail/pruning transcript fits inside the fixed budget $B_{\rm win}$.  If the
KL search reaches a long prefix, the framework stops on that value prefix and pays
it as a projected shadow; it is not copied through the Renyi lemma.  Thus the
copy key has entropy at most $B_{\rm win}$ by construction, while the cost of a
long alphabet-dependent value description is routed to $T(\Theta)$ instead of
to the active-rank recurrence.  This is the formal point at which the proposed
framework departs from ALWZ-style random restriction accounting.

The old-supported sparse residual selector is no longer exposed as an unpaid
atom of $W_L\subseteq L$.  Lemma~\ref{lem:residual-support-promotion} first
promotes all possible fresh witnesses globally, and
Lemma~\ref{lem:old-residual-routing} handles the remaining old-supported
positive support by exposing only old-coordinate pattern data through the
old-incidence ledger.

The global projected-shadow mass interface is supplied by
Proposition~\ref{prop:tree-entropy-ledger},
Proposition~\ref{prop:local-selector-certificates}, and
Theorem~\ref{thm:shadow-potential}.  The probability of reaching a local
terminal choice is not discarded: it is part of the terminal leaf transcript
before projection, and after merging leaves with the same projected shadow the
mass of that shadow can only increase.

The product-KL exit creates only a one-coordinate terminal pin or hash
coordinate.  Later uses of the close-pair theorem enter either terminal
shadow mode or the fixed-window hash ledger, and do not require a density
increment from product KL.

The hash ledger is label-specific: a low-collision output from a tuple node is
entered into the window of the actual public sample label which produced it.
The fixed-window theorem then works with the conditional law of that one
sample label and realizes the required independent copies by the Renyi copy
lemma.  Thus no step converts a low-collision statement for a mixed tuple
alphabet, or for the reference law $Q_E$, into an unproved statement about the
ambient code law.

\section{Conclusion}

Combining Theorem~\ref{thm:assembly}, Proposition~\ref{prop:birth}, and the
multi-shadow recurrence gives the conditional constant-base bookkeeping theorem
of Theorem~\ref{thm:main-reduction}.  The local exits are organized so that the
only permitted outputs are public blocks, terminal projected shadows, potential
drops, or pruning charges; private terminal transcripts are not allowed to enter
the final counting.  If the interface verifications recorded above are accepted
without additional hypotheses, the front reduction in
Section~\ref{sec:front-reduction} transfers the resulting encoded code bound
back to ordinary rank-$k$ set families.  Proposition~\ref{prop:constant-feasibility}
records the corresponding non-circular constant hierarchy.

\begin{thebibliography}{99}

\bibitem{AHS72}
H.~L. Abbott, D.~Hanson, and N.~Sauer,
\emph{Intersection theorems for systems of sets},
J. Combin. Theory Ser. A \textbf{12} (1972), 381--389.

\bibitem{AlonSpencer16}
N.~Alon and J.~H. Spencer,
\emph{The Probabilistic Method},
4th ed., Wiley, Hoboken, 2016.

\bibitem{ALWZ21}
R.~Alweiss, S.~Lovett, K.~Wu, and J.~Zhang,
\emph{Improved bounds for the sunflower lemma},
Ann. of Math. (2) \textbf{194} (2021), no.~3, 795--815.

\bibitem{ASU13}
N.~Alon, A.~Shpilka, and C.~Umans,
\emph{On sunflowers and matrix multiplication},
Comput. Complexity \textbf{22} (2013), no.~2, 219--243.

\bibitem{BCW21}
T.~Bell, S.~Chueluecha, and L.~Warnke,
\emph{Note on sunflowers},
Discrete Math. \textbf{344} (2021), no.~7, 112367.

\bibitem{CoverThomas06}
T.~M. Cover and J.~A. Thomas,
\emph{Elements of Information Theory},
2nd ed., Wiley-Interscience, Hoboken, 2006.

\bibitem{CLP17}
E.~Croot, V.~F. Lev, and P.~P. Pach,
\emph{Progression-free sets in $\mathbb Z_4^n$ are exponentially small},
Ann. of Math. (2) \textbf{185} (2017), no.~1, 331--337.

\bibitem{CsiszarKorner11}
I.~Csiszar and J.~Korner,
\emph{Information Theory: Coding Theorems for Discrete Memoryless Systems},
2nd ed., Cambridge University Press, Cambridge, 2011.

\bibitem{EG17}
J.~S. Ellenberg and D.~Gijswijt,
\emph{On large subsets of $\mathbb F_q^n$ with no three-term arithmetic
progression},
Ann. of Math. (2) \textbf{185} (2017), no.~1, 339--343.

\bibitem{ER60}
P.~Erdos and R.~Rado,
\emph{Intersection theorems for systems of sets},
J. London Math. Soc. \textbf{35} (1960), 85--90.

\bibitem{ES78}
P.~Erdos and E.~Szemeredi,
\emph{Combinatorial properties of systems of sets},
J. Combin. Theory Ser. A \textbf{24} (1978), no.~3, 308--313.

\bibitem{FKNP21}
K.~Frankston, J.~Kahn, B.~Narayanan, and J.~Park,
\emph{Thresholds versus fractional expectation-thresholds},
Ann. of Math. (2) \textbf{194} (2021), no.~2, 475--495.

\bibitem{FoxSudakov11}
J.~Fox and B.~Sudakov,
\emph{Dependent random choice},
Random Structures Algorithms \textbf{38} (2011), no.~1--2, 68--99.

\bibitem{Kostochka99}
A.~V. Kostochka,
\emph{Extremal problems on $\Delta$-systems},
in \emph{Numbers, Information and Complexity} (Bielefeld, 1998),
Kluwer Academic Publishers, Boston, 2000, 143--150.

\bibitem{KullbackLeibler51}
S.~Kullback and R.~A. Leibler,
\emph{On information and sufficiency},
Ann. Math. Statist. \textbf{22} (1951), no.~1, 79--86.

\bibitem{Mantel07}
W.~Mantel,
\emph{Problem 28},
Wiskundige Opgaven \textbf{10} (1907), 60--61.

\bibitem{NS17}
E.~Naslund and W.~Sawin,
\emph{Upper bounds for sunflower-free sets},
Forum Math. Sigma \textbf{5} (2017), e15, 10 pp.

\bibitem{Pinsker64}
M.~S. Pinsker,
\emph{Information and Information Stability of Random Variables and Processes},
Holden-Day, San Francisco, 1964.

\bibitem{Rao20}
A.~Rao,
\emph{Coding for sunflowers},
Discrete Anal. (2020), Paper No.~2, 8 pp.

\bibitem{Renyi61}
A.~Renyi,
\emph{On measures of entropy and information},
in \emph{Proceedings of the Fourth Berkeley Symposium on Mathematical
Statistics and Probability, Vol. I}, University of California Press, Berkeley,
1961, 547--561.

\bibitem{Shannon48}
C.~E. Shannon,
\emph{A mathematical theory of communication},
Bell System Tech. J. \textbf{27} (1948), 379--423, 623--656.

\bibitem{Tao16}
T.~Tao,
\emph{A symmetric formulation of the Croot--Lev--Pach--Ellenberg--Gijswijt
capset bound},
blog post, 18 May 2016,
\href{https://terrytao.wordpress.com/2016/05/18/a-symmetric-formulation-of-the-croot-lev-pach-ellenberg-gijswijt-capset-bound/}{terrytao.wordpress.com, 18 May 2016}.

\bibitem{Tao20}
T.~Tao,
\emph{The sunflower lemma via Shannon entropy},
blog post, 20 July 2020,
\href{https://terrytao.wordpress.com/2020/07/20/the-sunflower-lemma-via-shannon-entropy/}{terrytao.wordpress.com, 20 July 2020}.

\end{thebibliography}

\end{document}
